% Generated mechanically from frozen input p07/ja/curves.tex.
% Do not edit this generated file; edit the bound locale source instead.

% OK, start here.
%

\setcounter{chapter}{52}
\chapter{代数曲線}



\phantomsection
\label{curves-section-phantom}


\section{序論}
\label{curves-section-introduction}

\noindent
本章では代数曲線の理論の一部を展開する．
複素数体上の代数曲線を扱う参考文献として，
書籍 \cite{ACGH} がある．

\medskip\noindent
既知の結果．一般的な代数幾何に加えて，
代数曲線に関するいくつかの具体的な結果をすでに証明している．
以下にそれらを列挙する．
\begin{enumerate}
\item $1$ 次元 Noether スキームのアフィン開集合およびその上の
豊富な可逆加群について，Varieties，節
\ref{varieties-section-dimension-one} で論じた．
\item 曲線はアフィンまたは射影的のいずれかであることを，
Varieties，節 \ref{varieties-section-curves} で見た．
\item 固有曲線上の局所自由加群の次数について，Varieties，節
\ref{varieties-section-divisors-curves} で論じた．
\item 代数閉体上の非特異射影曲線の Picard スキームについて，
Picard Schemes of Curves，節 \ref{pic-section-introduction} で論じた．
\end{enumerate}





\section{曲線と関数体}
\label{curves-section-curves-function-fields}

\noindent
本節では，Varieties，節
\ref{varieties-section-varieties-rational-maps} の結果を
曲線の場合について詳述する．

\begin{lemma}
\label{curves-lemma-extend-over-dvr}
$k$ を体とする．$X$ を曲線，$Y$ を固有多様体とする．
$U \subset X$ を空でない開集合とし，$f : U \to Y$ を射とする．
$x \in X$ が閉点で，$\mathcal{O}_{X, x}$ が離散付値環であるならば，
開集合 $U \subset U' \subset X$ が存在して $x$ を含み，さらに
多様体の射 $f' : U' \to Y$ が存在して $f$ を延長する．
\end{lemma}

\begin{proof}
これは Morphisms，補題 \ref{morphisms-lemma-extend-across}
の特別な場合である．
\end{proof}

\begin{lemma}
\label{curves-lemma-extend-over-normal-curve}
$k$ を体とする．$X$ を正規曲線，$Y$ を固有多様体とする．
$X$ から $Y$ への有理写像の集合は，射 $X \to Y$ の集合と同じである．
\end{lemma}

\begin{proof}
$X$ から $Y$ への有理写像は射 $X \to Y$ へ延長できる．
これは補題 \ref{curves-lemma-extend-over-dvr} による．実際，すべての局所環は
離散付値環である（例えば Varieties，補題
\ref{varieties-lemma-regular-point-on-curve} による）．
逆に，二つの射 $f,g: X \to Y$ が有理写像として同値ならば，
Morphisms，補題 \ref{morphisms-lemma-equality-of-morphisms} により
$f = g$ である．
\end{proof}

\begin{lemma}
\label{curves-lemma-flat}
$k$ を体とする．$f : X \to Y$ を $k$ 上の曲線の間の
非定数射とする．$Y$ が正規ならば，$f$ は平坦である．
\end{lemma}

\begin{proof}
$x \in X$ を取り，その像を $y \in Y$ とする．このとき
$\mathcal{O}_{Y, y}$ は体または離散付値環である
（Varieties，補題 \ref{varieties-lemma-regular-point-on-curve}）．
$f$ は非定数なので支配的である（実際，$X$ の一般点を $Y$ の
一般点へ写さなければならない）．したがって
$\mathcal{O}_{Y, y} \to \mathcal{O}_{X, x}$ は単射である
（Morphisms，補題 \ref{morphisms-lemma-dominant-between-integral}）．
ゆえに $\mathcal{O}_{X, x}$ は $\mathcal{O}_{Y, y}$-加群として捩れをもたず，
したがって $\mathcal{O}_{X, x}$ は $\mathcal{O}_{Y, y}$-加群として平坦である．
これは More on Algebra，補題
\ref{more-algebra-lemma-valuation-ring-torsion-free-flat} による．
\end{proof}

\begin{lemma}
\label{curves-lemma-finite}
$k$ を体とする．$f : X \to Y$ を $k$ 上のスキームの射とする．
次を仮定する．
\begin{enumerate}
\item $Y$ は $k$ 上分離的である；
\item $X$ は次元 $\leq 1$ で $k$ 上固有である；
\item $f(Z)$ は少なくとも二点をもつ．ここで $Z \subset X$ は
次元 $1$ の任意の既約成分である．
\end{enumerate}
このとき $f$ は有限である．
\end{lemma}

\begin{proof}
$f$ は Morphisms，補題
\ref{morphisms-lemma-image-proper-scheme-closed} により固有である．
したがって $f(X)$ は閉であり，閉点の像も閉点である．
$y \in Y$ を $X$ の閉点の像とする．条件 (3) により，
$f^{-1}(\{y\})$ は $X$ の閉部分集合であり，$1$ 次元既約成分の
一般点を一つも含まない．したがって $f^{-1}(\{y\})$ は有限である．
ゆえに $f$ は $y$ のある開近傍上で有限である．これは
More on Morphisms，補題
\ref{more-morphisms-lemma-proper-finite-fibre-finite-in-neighbourhood} による
（$Y$ が Noether スキームならば，より簡単な Cohomology of Schemes，
補題 \ref{coherent-lemma-proper-finite-fibre-finite-in-neighbourhood}
を用いてもよい）．上で見たように，このような点 $y$ は十分多く存在するので，
証明は完了する．
\end{proof}

\begin{lemma}
\label{curves-lemma-extend-to-completion}
$k$ を体とする．$X \to Y$ を多様体の射とし，
$Y$ は固有，$X$ は曲線であるとする．
このとき分解 $X \to \overline{X} \to Y$ が存在し，
$X \to \overline{X}$ は開埋め込み，$\overline{X}$ は射影曲線である．
\end{lemma}

\begin{proof}
これは補題 \ref{curves-lemma-extend-over-dvr} と Varieties，補題
\ref{varieties-lemma-reduced-dim-1-projective-completion} から明らかである．
\end{proof}

\noindent
以下が本節の主定理である．多様体の射 $f : X \to Y$ について，
像 $f(X)$ が一点 $y$ のみからなり，その点が $Y$ に属するとき，
この射は {\it 定数}であるという．このとき $y$ は $Y$ の閉点である
（$X$ の閉点の像が $Y$ の閉点となるためである）．

\begin{theorem}
\label{curves-theorem-curves-rational-maps}
$k$ を体とする．次の圏は標準的に同値である．
\begin{enumerate}
\item 有限生成体拡大 $K/k$ で超越次数が $1$ のものの圏；
\item 曲線と支配的有理写像の圏；
\item 正規射影曲線と非定数射の圏；
\item 非特異射影曲線と非定数射の圏；
\item 正則射影曲線と非定数射の圏；
\item 正規固有曲線と非定数射の圏．
\end{enumerate}
\end{theorem}

\begin{proof}
圏 (1) と (2) の間の同値は，Varieties，定理
\ref{varieties-theorem-varieties-rational-maps} の同値の制限である．
実際，多様体が曲線であることと，その関数体の超越次数が $1$ であることは
同値である．例えば Varieties，補題
\ref{varieties-lemma-dimension-locally-algebraic} を参照されたい．

\medskip\noindent
圏 (3)，(4)，(5)，(6) は同じである．まず，「正則」と「非特異」は
同義語である．Properties，定義
\ref{properties-definition-regular} を参照されたい．
Noether $1$ 次元スキームについては，正規であることと正則であることは
同値である（Properties，補題 \ref{properties-lemma-regular-normal} および
\ref{properties-lemma-normal-dimension-1-regular}）．曲線の場合については
Varieties，補題 \ref{varieties-lemma-regular-point-on-curve} を参照されたい．
したがって (3) と (5) は同じである．最後に Varieties，補題
\ref{varieties-lemma-dim-1-proper-projective} により，(6) と (3) は同じである．

\medskip\noindent
非特異射影曲線の非定数射 $f : X \to Y$ が与えられると，$f$ は
一般点 $\eta$（$X$ の一般点）を一般点 $\xi$（$Y$ の一般点）へ写す．したがって射
$k(Y) = \mathcal{O}_{Y, \xi} \to \mathcal{O}_{X, \eta} = k(X)$
を圏 (1) において得る．二つの射 $f,g: X \to Y$ が同じ射
$k(Y) \to k(X)$ を与えるならば，(1) と (2) の同値により，
$f$ と $g$ は有理写像として同値である．ゆえに補題
\ref{curves-lemma-extend-over-normal-curve} により $f=g$ である．
逆に，圏 (1) の写像 $k(Y) \to k(X)$ が与えられたとする．
すると射 $U \to Y$ を，ある空でない開集合 $U \subset X$ に対して得る．
補題 \ref{curves-lemma-extend-over-dvr} により，これは $X$ 全体へ延長され，
圏 (5) の射が得られる．したがって完全忠実関手 (5)$\to$(1) が存在する．

\medskip\noindent
証明を終えるには，(1) の任意の $K/k$ が正規射影曲線の関数体であることを
示せばよい．$K = k(X)$ となる曲線 $X$ が存在することはすでに分かっている．
$X$ をその正規化（これは $X$ と双有理な多様体である）で置き換えることにより，
$X$ は正規であると仮定してよい（Varieties，補題
\ref{varieties-lemma-normalization-locally-algebraic}）．次に Varieties，補題
\ref{varieties-lemma-reduced-dim-1-projective-completion} のように，
$X \to \overline{X}$ であって
$\overline{X} \setminus X = \{x_1, \ldots, x_n\}$ を満たすものを選ぶ．
$X$ は正規であり，各局所環 $\mathcal{O}_{\overline{X}, x_i}$ も正規なので，
$\overline{X}$ は求める正規射影曲線である．
（注意：Varieties，補題
\ref{varieties-lemma-dim-1-projective-completion} を用いてまずコンパクト化し，
次に Varieties，補題
\ref{varieties-lemma-normalization-locally-algebraic} を用いて正規化してもよい．
この順序なら，やや扱いにくい Morphisms，補題
\ref{morphisms-lemma-relative-normalization-normal-codim-1} を用いずに済む．）
\end{proof}

\begin{definition}
\label{curves-definition-normal-projective-model}
$k$ を体とし，$X$ を曲線とする．
{\it $X$ の非特異射影モデル}とは，組 $(Y, \varphi)$ であって，
$Y$ が非特異射影曲線であり，$\varphi : k(X) \to k(Y)$ が
関数体の同型であるものをいう．
\end{definition}

\noindent
定理 \ref{curves-theorem-curves-rational-maps} により，非特異射影モデルは
一意的な同型を除いて定まる．したがってしばしば「その非特異射影モデル」
という．通常，記法から $\varphi$ を省略する．注意すべきことに，
$Y$ は $k$ 上滑らかとは限らないが，補題
\ref{curves-lemma-nonsingular-model-smooth} により，これは正標数でのみ起こり得る．

\begin{lemma}
\label{curves-lemma-nonsingular-model-smooth}
$k$ を体とする．$X$ を曲線とし，$Y$ を $X$ の非特異射影モデルとする．
$k$ が完全体ならば，$Y$ は滑らかな射影曲線である．
\end{lemma}

\begin{proof}
例えば Varieties，補題
\ref{varieties-lemma-regular-point-on-curve} を参照されたい．
\end{proof}

\begin{lemma}
\label{curves-lemma-smooth-models}
$k$ を体とする．$X$ を $k$ 上幾何学的既約な曲線とする．
体拡大 $K/k$ に対し，$Y_K$ を $(X_K)_{red}$ の非特異射影モデルと書く．
\begin{enumerate}
\item $X$ が固有ならば，$Y_K$ は $X_K$ の正規化である．
\item 有限純非分離拡大 $K/k$ が存在し，$Y_K$ は滑らかである．
\item $Y_K$ が滑らかであるとき\footnote{幾何学的に被約であるだけでもよい．}，
$H^0(Y_K, \mathcal{O}_{Y_K}) = K$ が成り立つ．
\item 可換図式
$$
\xymatrix{
\Omega & K' \ar[l] \\
K \ar[u] & k \ar[l] \ar[u]
}
$$
をなす体が与えられ，$Y_K$ と $Y_{K'}$ が滑らかならば，
$Y_\Omega = (Y_K)_\Omega = (Y_{K'})_\Omega$.
\end{enumerate}
\end{lemma}

\begin{proof}
$X'$ を $X$ の非特異射影モデルとする．このとき $X'$ と $X$ は
互いに同型な空でない開部分スキームをもつ．特に，$X'$ は幾何学的既約である．
これは $X$ が幾何学的既約だからである（細部は省略する）．したがって $X$ は射影的であると
仮定してよい．

\medskip\noindent
$X$ が固有であると仮定する．すると $X_K$ は固有であり，したがって正規化
$(X_K)^\nu$ は，固有スキーム上有限なスキームとして固有である
（Varieties，補題 \ref{varieties-lemma-normalization-locally-algebraic}，
および Morphisms，補題 \ref{morphisms-lemma-finite-proper} と
\ref{morphisms-lemma-composition-proper}）．一方，$X_K$ は既約である．
これは $X$ が幾何学的既約だからである．したがって $X_K^\nu$ は固有，正規，既約であり，
$(X_K)_{red}$ と双有理である．固有曲線は射影的なので，これで (1) が従う
（Varieties，補題 \ref{varieties-lemma-dim-1-proper-projective}）．

\medskip\noindent
(2) の証明．$X$ は固有で (1) が成り立つので，Varieties，補題
\ref{varieties-lemma-finite-extension-geometrically-normal} を適用し，
有限純非分離拡大 $K/k$ であって $Y_K$ が幾何学的に正規となるものを得る．
曲線について正規性と正則性は同値なので，$Y_K$ は幾何学的に正則である
（Properties，補題 \ref{properties-lemma-normal-dimension-1-regular}）．
このとき $Y$ は Varieties，補題
\ref{varieties-lemma-geometrically-regular-smooth} により滑らかな多様体である．

\medskip\noindent
$Y_K$ が幾何学的に被約ならば，$Y_K$ は幾何学的に整であり
（Varieties，補題 \ref{varieties-lemma-geometrically-integral}），
Varieties，補題 \ref{varieties-lemma-regular-functions-proper-variety} により
$H^0(Y_K, \mathcal{O}_{Y_K}) = K$ である．滑らかな多様体は幾何学的に被約
（実際には幾何学的に正則）なので，これで (3) が従う．Varieties，補題
\ref{varieties-lemma-geometrically-regular-smooth} を参照されたい．

\medskip\noindent
$Y_K$ が滑らかならば，任意の拡大 $\Omega/K$ に対して基底変換
$(Y_K)_\Omega$ は $\Omega$ 上滑らかである
（Morphisms，補題 \ref{morphisms-lemma-base-change-smooth}）．
したがって $Y_\Omega = (Y_K)_\Omega$ は明らかである．これで (4) が従う．
\end{proof}






\section{線形系列}
\label{curves-section-linear-series}

\noindent
古典的な扱いから少し離れ（注意 \ref{curves-remark-classical-linear-series} を参照），
線形系列を次のように定義する．

\begin{definition}
\label{curves-definition-linear-series}
$k$ を体とする．$X$ を次元 $\leq 1$ の $k$ 上の固有スキームとする．
$d \geq 0$，$r \geq 0$ とする．
{\it 次数 $d$，次元 $r$ の線形系列}とは組 $(\mathcal{L}, V)$ であって，
$\mathcal{L}$ が可逆 $\mathcal{O}_X$-加群で，次数が $d$ であり
（Varieties，定義 \ref{varieties-definition-degree-invertible-sheaf}）であり，
$V \subset H^0(X, \mathcal{L})$ が $k$-部分ベクトル空間で，次元が $r + 1$ である
ものをいう．これを略して，$(\mathcal{L}, V)$ を {\it $\mathfrak g^r_d$}
（$X$ 上）と書く．
\end{definition}

\noindent
主として $X$ が非特異固有曲線である場合にこの定義を用いる．
実際，上の定義は古典的な $\mathfrak g^r_d$ の定義を一般化する一つの方法に
すぎない．例えば $X$ が固有曲線ならば，$\mathcal{L}$ として捩れなし連接
$\mathcal{O}_X$-加群で階数 $1$ のものを許すことで線形系列を一般化できる．
非特異曲線上では捩れなし連接加群はすべて局所自由なので，これは非特異
固有曲線に対する本書の概念と一致する．

\medskip\noindent
次の補題は，固有非特異曲線に対する線形系列の幾何学的意味を説明する．

\begin{lemma}
\label{curves-lemma-linear-series}
$k$ を体とする．$X$ を $k$ 上の非特異固有曲線とする．
$(\mathcal{L}, V)$ を $\mathfrak g^r_d$ とし，$X$ 上で考える．このとき射
$$
\varphi : X \longrightarrow \mathbf{P}^r_k = \text{Proj}(k[T_0, \ldots, T_r])
$$
である $k$ 上の多様体の射と，写像
$\alpha : \varphi^*\mathcal{O}_{\mathbf{P}^r_k}(1) \to \mathcal{L}$
が存在し，$\varphi^*T_0, \ldots, \varphi^*T_r$ は $V$ の基底へ送られる．
この写像は $\alpha$ によるものである．
\end{lemma}

\begin{proof}
$s_0, \ldots, s_r \in V$ を $k$-基底とする．$X$ は非特異なので，像
$\mathcal{L}' \subset \mathcal{L}$ は，写像
$s_0, \ldots, s_r : \mathcal{O}_X^{\oplus r + 1} \to \mathcal{L}$
の像として，可逆 $\mathcal{O}_X$-加群である．
これは例えば Divisors，補題
\ref{divisors-lemma-torsion-free-over-regular-dim-1} による．次に
Constructions，補題 \ref{constructions-lemma-projective-space} を用いて射
$$
\varphi = \varphi_{(\mathcal{L}', (s_0, \ldots, s_r))} :
X \longrightarrow \mathbf{P}^r_k
$$
を得る．これは補題の主張にある射である．
\end{proof}

\begin{lemma}
\label{curves-lemma-linear-series-trivial-existence}
$k$ を体とする．$X$ を次元 $\leq 1$ の $k$ 上の固有スキームとする．
$X$ が $\mathfrak g^r_d$ をもつならば，$X$ は $\mathfrak g^s_d$ をもち，
これは任意の $0 \leq s \leq r$ に対して成り立つ．
\end{lemma}

\begin{proof}
これは，ベクトル空間 $V$ が次元 $r + 1$ で $k$ 上定義されていれば，
次元 $s + 1$ の線形部分空間を任意の $0 \leq s \leq r$ に対して
もつためである．
\end{proof}

\begin{lemma}
\label{curves-lemma-g1d}
$k$ を体とする．$X$ を $k$ 上の非特異固有曲線とする．
$(\mathcal{L}, V)$ を $\mathfrak g^1_d$ とし，$X$ 上で考える．このとき補題
\ref{curves-lemma-linear-series} の射 $\varphi : X \to \mathbf{P}^1_k$ は，
次のいずれかを満たす．
\begin{enumerate}
\item 非定数で，次数が $\leq d$ である；
\item $\mathbf{P}^1_k$ の閉点を経由して分解し，この場合
$H^0(X, \mathcal{O}_X) \not = k$ である．
\end{enumerate}
\end{lemma}

\begin{proof}
補題 \ref{curves-lemma-linear-series} により，
$\mathcal{L}' = \varphi^*\mathcal{O}_{\mathbf{P}^1_k}(1)$
から $\mathcal{L}' \to \mathcal{L}$ への非零写像が存在する．
したがって Varieties，補題
\ref{varieties-lemma-check-invertible-sheaf-trivial} により
$0 \leq \deg(\mathcal{L}') \leq d$ である．
$\deg(\mathcal{L}') = 0$ ならば，同じ補題により
$\mathcal{L}' \cong \mathcal{O}_X$ であり，線形独立な二つの切断が
存在するので場合 (2) となる．$\deg(\mathcal{L}') > 0$ ならば，
$\varphi$ は非定数である（定数射による可逆加群の引戻しは自明だからである）．
したがって
$$
\deg(\mathcal{L}') =
\deg(X/\mathbf{P}^1_k) \deg(\mathcal{O}_{\mathbf{P}^1_k}(1))
$$
が Varieties，補題 \ref{varieties-lemma-degree-pullback-map-proper-curves}
により成り立つ．$\mathcal{O}_{\mathbf{P}^1_k}(1)$ の次数は $1$ なので，
これで証明が完了する．
\end{proof}

\begin{lemma}
\label{curves-lemma-grd-inequalities}
$k$ を体とする．$X$ を $k$ 上の固有曲線で
$H^0(X, \mathcal{O}_X) = k$ を満たすものとする．$X$ が
$\mathfrak g^r_d$ をもつならば $r \leq d$ である．等号が成り立つならば
$H^1(X, \mathcal{O}_X) = 0$，すなわち $X$ の種数
（定義 \ref{curves-definition-genus}）は $0$ である．
\end{lemma}

\begin{proof}
$(\mathcal{L}, V)$ を $\mathfrak g^r_d$ とする．これは $r$ を増加させる
だけなので，$V = H^0(X, \mathcal{L})$ と仮定してよい．非零元
$s \in V$ を選ぶ．すると $s$ の零点スキームは有効 Cartier 因子
$D \subset X$ であり，$\mathcal{L} = \mathcal{O}_X(D)$ が成り立ち，
短完全列
$$
0 \to \mathcal{O}_X \to \mathcal{L} \to \mathcal{L}|_D \to 0
$$
を得る．Divisors，補題 \ref{divisors-lemma-characterize-OD} および
注意 \ref{divisors-remark-ses-regular-section} を参照されたい．
Varieties，補題 \ref{varieties-lemma-degree-effective-Cartier-divisor} により
$\deg(D) = \deg(\mathcal{L}) = d$ である．$D$ は Artin スキームなので
$\mathcal{L}|_D \cong \mathcal{O}_D$ である\footnote{ここでは，
$D \to \Spec(k)$ が有限なので，これは Divisors，補題
\ref{divisors-lemma-finite-trivialize-invertible-upstairs} から従う．}．
したがって
$$
\dim_k H^0(D, \mathcal{L}|_D) = \dim_k H^0(D, \mathcal{O}_D) = \deg(D) = d
$$
である．一方，仮定により $\dim_k H^0(X, \mathcal{O}_X) = 1$ かつ
$\dim H^0(X, \mathcal{L}) = r + 1$ である．したがって
$r + 1 \leq 1 + d$，すなわち補題のとおり $r \leq d$ である．

\medskip\noindent
等号が成り立つと仮定する．このとき
$H^0(X, \mathcal{L}) \to H^0(X, \mathcal{L}|_D)$ は全射である．
$H^1(X, \mathcal{L})$ が零だと分かれば，コホモロジー長完全列により
$H^1(X, \mathcal{O}_X)$ が零と分かり，証明は完了する．そこで
$\mathcal{L}$ を，消滅をもつ十分大きな冪で置き換える．上で見たように
$\mathcal{L}|_D$ は自明な可逆加群なので，切断 $t$ を $\mathcal{L}$ から取り，
$D$ への制限が $\mathcal{L}|_D$ を生成するものを取れる．乗法写像
$$
\mu :
H^0(X, \mathcal{L}) \otimes_k H^0(X, \mathcal{L})
\longrightarrow
H^0(X, \mathcal{L}^{\otimes 2})
$$
と短完全列
$$
0 \to \mathcal{L} \xrightarrow{s}
\mathcal{L}^{\otimes 2} \to \mathcal{L}^{\otimes 2}|_D \to 0
$$
を考える．$H^0(\mathcal{L}) \to H^0(\mathcal{L}|_D)$ は全射であり，
$t$ は $\mathcal{L}|_D$ の自明化へ写るので，
$\mu(H^0(X, \mathcal{L}) \otimes t)$ は
$H^0(X, \mathcal{L}^{\otimes 2})$ の部分空間であって，
$\mathcal{L}^{\otimes 2}|_D$ の大域切断へ全射となる．したがって
$$
\dim H^0(X, \mathcal{L}^{\otimes 2}) = r + 1 + d = 2r + 1 =
\deg(\mathcal{L}^{\otimes 2}) + 1
$$
である．ゆえに $\mathcal{L}^{\otimes 2}$ は $\mathcal{L}$ と同じ性質，
すなわち大域切断空間の次元が次数に一を加えたものに等しいという性質をもつ．
$\mathcal{L}$ は豊富であるから（Varieties，補題
\ref{varieties-lemma-ample-curve}），ある $n_0$ が存在して，
$\mathcal{L}^{\otimes n}$ の $H^1$ は任意の $n \geq n_0$ に対して消滅する
（Cohomology of Schemes，補題
\ref{coherent-lemma-coherent-proper-ample}）．したがって，上の議論を
$\mathcal{L}^{\otimes n}$ に適用し，$n = 2^m$ と取る．ここで $m$ は十分大きい．
適用すれば補題が従う．
\end{proof}

\begin{remark}[古典的定義]
\label{curves-remark-classical-linear-series}
$X$ を代数閉体 $k$ 上の滑らかな射影曲線とする．二つの有効 Cartier 因子
$D, D' \subset X$ が {\it 線形同値}であるとは，
$\mathcal{O}_X(D) \cong \mathcal{O}_X(D')$ が $\mathcal{O}_X$-加群として
成り立つことをいう．$\Pic(X) = \text{Cl}(X)$ なので
（Divisors，補題 \ref{divisors-lemma-local-rings-UFD-c1-bijective}），
$D$ と $D'$ が線形同値であることと，$D$ と $D'$ に付随する Weil 因子が
$\text{Cl}(X)$ の同じ元を定めることは同値である．有効 Cartier 因子
$D \subset X$ で次数が $d$ のものが与えられたとき，
{\it 完全線形系}または {\it 完全線形系列} $|D|$（$D$ に付随するもの）とは，
有効 Cartier 因子 $E \subset X$ で $D$ と線形同値なものの集合である．
別の言い方をすれば，$|D|$ は射
$$
\gamma_d :
\underline{\Hilbfunctor}^d_{X/k}
\longrightarrow
\underline{\Picardfunctor}^d_{X/k}
$$
（Picard Schemes of Curves，補題 \ref{pic-lemma-picard-pieces}）の，
$\mathcal{O}_X(D)$ に対応する閉点上のファイバーの閉点集合である．
これにより $|D|$ は自然なスキーム構造をもち，実際
$|D| \cong \mathbf{P}^m_k$ であって
$m + 1 = h^0(\mathcal{O}_X(D))$ となる．より標準的には
$$
|D| = \mathbf{P}(H^0(X, \mathcal{O}_X(D))^\vee)
$$
である．ここで $(-)^\vee$ は $k$-線形双対を表し，$\mathbf{P}$ は
Constructions，例 \ref{constructions-example-projective-space} のものである．
この言葉では，$X$ 上の {\it 線形系}または {\it 線形系列}とは，
線形方程式で切り出せる閉部分多様体 $L \subset |D|$ をいう．
$L$ の次元が $r$ ならば $L = \mathbf{P}(V^\vee)$ であり，ここで
$V \subset H^0(X, \mathcal{O}_X(D))$ は次元 $r + 1$ の線形部分空間である．
したがって古典的線形系列 $L \subset |D|$ は，上で定義した線形系列
$(\mathcal{O}_X(D), V)$ に対応する．
\end{remark}







\section{双対性}
\label{curves-section-duality}

\noindent
本節では，体上の固有な $1$ 次元概形に対する双対化複体および双対性についての，
きわめて一般的な理論から得られる帰結を詳しく述べる．体上の高次元固有概形に
対する類似の議論に関心がある場合は，Duality for Schemes，節
\ref{duality-section-duality-proper-over-field} を参照されたい．

\begin{lemma}
\label{curves-lemma-duality-dim-1}
$X$ を次元 $\leq 1$ の，体 $k$ 上の固有概形とする．次の性質をもつ
双対化複体 $\omega_X^\bullet$ が存在する．
\begin{enumerate}
\item $H^i(\omega_X^\bullet)$ が非零となり得るのは $i = -1, 0$ の場合だけである．
\item $\omega_X = H^{-1}(\omega_X^\bullet)$ は，支えが次元 $1$ の
既約成分からなる連接 Cohen--Macaulay 加群である．
\item 閉点 $x \in X$ に対し，加群 $H^0(\omega_{X, x}^\bullet)$ が非零で
あるための必要十分条件は，次のいずれかが成り立つことである：
\begin{enumerate}
\item $\dim(\mathcal{O}_{X, x}) = 0$，または
\item $\dim(\mathcal{O}_{X, x}) = 1$ であり，かつ
$\mathcal{O}_{X, x}$ が Cohen--Macaulay でない．
\end{enumerate}
\item $K \in D_\QCoh(\mathcal{O}_X)$ に対し，シフトおよび
識別三角形と両立する関手的同型\footnote{この性質は $\omega_X^\bullet$ を
$D_\QCoh(\mathcal{O}_X)$ において，Yoneda の補題により一意的同型を除いて
特徴づける．$\omega_X^\bullet$ は $D^b_{\textit{Coh}}(\mathcal{O}_X)$ に
属するので，実際には $K \in D^b_{\textit{Coh}}(\mathcal{O}_X)$ の場合を
考えれば十分である．}
$$
\Ext^i_X(K, \omega_X^\bullet) = \Hom_k(H^{-i}(X, K), k)
$$
が存在する．
\item 関手的同型
$\Hom(\mathcal{F}, \omega_X) = \Hom_k(H^1(X, \mathcal{F}), k)$ が存在する．
ここで $\mathcal{F}$ は $X$ 上の準連接加群である．
\item $X \to \Spec(k)$ が相対次元 $1$ の滑らかな射ならば，
$\omega_X \cong \Omega_{X/k}$ である．
\end{enumerate}
\end{lemma}

\begin{proof}
構造射を $f : X \to \Spec(k)$ と書く．まず，Duality for Schemes，備考
\ref{duality-remark-relative-dualizing-complex} で述べられている相対双対化複体
$$
\omega_X^\bullet = \omega_{X/k}^\bullet = a(\mathcal{O}_{\Spec(k)})
$$
から始める．$a$ は
$f_* : D_\QCoh(\mathcal{O}_X) \to D(\mathcal{O}_{\Spec(k)})$ の右随伴なので，
性質 (4) は構成から成り立つ．$f$ は固有であるから，定義により
$f^!(\mathcal{O}_{\Spec(k)}) = a(\mathcal{O}_{\Spec(k)})$ である．
Duality for Schemes，節 \ref{duality-section-upper-shriek} を参照されたい．
したがって $\omega_X^\bullet$ と $\omega_X$ は
Duality for Schemes，例 \ref{duality-example-proper-over-local} および
Duality for Schemes，例
\ref{duality-example-equidimensional-over-field} に現れるものに一致する．
(1) と (2) は Duality for Schemes，補題
\ref{duality-lemma-vanishing-good-dualizing} から従う．
閉点 $x \in X$ に対し，$\omega_{X, x}^\bullet$ は
$\mathcal{O}_{X, x}$ 上の正規化された双対化複体である．
Duality for Schemes，補題
\ref{duality-lemma-good-dualizing-normalized} を参照されたい．
したがって (3) は Dualizing Complexes，補題
\ref{dualizing-lemma-apply-CM} から従う．
(5) は，連接な $\mathcal{F}$ については Duality for Schemes，補題
\ref{duality-lemma-dualizing-module-proper-over-A} から従い，一般の場合は
$K = \mathcal{F}[0]$ および $i = -1$ として (4) を展開すれば従う．
(6) は Duality for Schemes，補題
\ref{duality-lemma-smooth-proper} から従う．
\end{proof}

\begin{lemma}
\label{curves-lemma-duality-dim-1-CM}
$X$ を体 $k$ 上の固有概形で，Cohen--Macaulay かつ次元 $1$ の等次元な
ものとする．補題 \ref{curves-lemma-duality-dim-1} の加群 $\omega_X$ は
次の性質をもつ．
\begin{enumerate}
\item $\omega_X$ は $X$ 上の双対化加群である
（Duality for Schemes，節 \ref{duality-section-dualizing-module}）．
\item $\omega_X$ は支えが $X$ である連接 Cohen--Macaulay 加群である．
\item 関手的同型
$\Ext^i_X(K, \omega_X[1]) = \Hom_k(H^{-i}(X, K), k)$ が存在し，
$K \in D_\QCoh(X)$ に対してシフトと両立する．
\item 関手的同型
$\Ext^{1 + i}(\mathcal{F}, \omega_X) = \Hom_k(H^{-i}(X, \mathcal{F}), k)$
が存在する．ここで $\mathcal{F}$ は $X$ 上の準連接加群である．
\end{enumerate}
\end{lemma}

\begin{proof}
補題 \ref{curves-lemma-duality-dim-1} の証明から，$\omega_X$ は
Duality for Schemes，例 \ref{duality-example-proper-over-local} に
現れるものなので，双対化加群であることを思い出そう．他の主張は
補題 \ref{curves-lemma-duality-dim-1} と，$\omega_X^\bullet = \omega_X[1]$ となる
ことから従う．後者は $X$ が Cohen--Macaulay であるためである
（Duality for Schemes，補題 \ref{duality-lemma-dualizing-module-CM-scheme}）．
\end{proof}

\begin{remark}
\label{curves-remark-rework-duality-locally-free}
$X$ を次元 $\leq 1$ の，体 $k$ 上の固有概形とし，
$\omega_X^\bullet$ と $\omega_X$ を補題 \ref{curves-lemma-duality-dim-1} の
ものとする．$\mathcal{E}$ を有限局所自由 $\mathcal{O}_X$-加群，
$\mathcal{E}^\vee$ をその双対とすると，標準同型
$$
\Hom_k(H^{-i}(X, \mathcal{E}), k) =
H^i(X, \mathcal{E}^\vee \otimes_{\mathcal{O}_X}^\mathbf{L} \omega_X^\bullet)
$$
がある．これは補題と Cohomology，補題
\ref{cohomology-lemma-dual-perfect-complex} から従う．
$X$ が Cohen--Macaulay かつ次元 $1$ の等次元ならば，
補題 \ref{curves-lemma-duality-dim-1-CM} により標準同型
$$
\Hom_k(H^{-i}(X, \mathcal{E}), k) =
H^{1 + i}(X, \mathcal{E}^\vee \otimes_{\mathcal{O}_X} \omega_X)
$$
がある．特に $\mathcal{L}$ が可逆 $\mathcal{O}_X$-加群ならば，
$$
\dim_k H^0(X, \mathcal{L}) =
\dim_k H^1(X, \mathcal{L}^{\otimes -1} \otimes_{\mathcal{O}_X} \omega_X)
$$
および
$$
\dim_k H^1(X, \mathcal{L}) =
\dim_k H^0(X, \mathcal{L}^{\otimes -1} \otimes_{\mathcal{O}_X} \omega_X)
$$
を得る．
\end{remark}

\noindent
ここで双対化複体の整合性を確認しておく．

\begin{lemma}
\label{curves-lemma-sanity-check-duality}
$X$ を次元 $\leq 1$ の，体 $k$ 上の固有概形とし，
$\omega_X^\bullet$ と $\omega_X$ を補題 \ref{curves-lemma-duality-dim-1} の
ものとする．
\begin{enumerate}
\item $X \to \Spec(k)$ が $X \to \Spec(k') \to \Spec(k)$ と分解し，
ここで $k'$ が体であるならば，
$\omega_X^\bullet$ と $\omega_X$ は，$k$ を $k'$ に置き換えた
性質 (4)，(5)，(6) を満たす．
\item $K/k$ が体拡大ならば，$\omega_X^\bullet$ と $\omega_X$ の
基底変換 $X_K$ への引戻しは，射 $X_K \to \Spec(K)$ に対する
補題 \ref{curves-lemma-duality-dim-1} のものとなる．
\end{enumerate}
\end{lemma}

\begin{proof}
構造射を $f : X \to \Spec(k)$ と書く．主張 (1) の意味は，
$\omega_X^\bullet$ と $\omega_X$ が射
$f' : X \to \Spec(k')$ に対する補題 \ref{curves-lemma-duality-dim-1} の
ものだということである．補題 \ref{curves-lemma-duality-dim-1} の証明では
$\omega_X^\bullet = a(\mathcal{O}_{\Spec(k)})$ とおいた．ここで $a$ は
$f$ に対する Duality for Schemes，補題
\ref{duality-lemma-twisted-inverse-image} の右随伴である．
したがって
$a(\mathcal{O}_{\Spec(k)}) \cong a'(\mathcal{O}_{\Spec(k)})$
を示せばよい．ここで $a'$ は $f'$ に対する
Duality for Schemes，補題
\ref{duality-lemma-twisted-inverse-image} の右随伴である．
$k' \subset H^0(X, \mathcal{O}_X)$ なので，$k'/k$ は有限拡大である
（Cohomology of Schemes，補題
\ref{coherent-lemma-proper-over-affine-cohomology-finite}）．
随伴の一意性により $a = a' \circ b$ である．ここで $b$ は
$g : \Spec(k') \to \Spec(k)$ に対する Duality for Schemes，補題
\ref{duality-lemma-twisted-inverse-image} の右随伴である．
言い換えれば $f^! = (f')^! \circ g^!$ である．したがって
$\Hom_k(k', k) \cong k'$ が $k'$-加群として成り立つことを示せば十分である．
Duality for Schemes，例
\ref{duality-example-affine-twisted-inverse-image} を参照されたい．
これは両辺が同じ次元の $k'$-ベクトル空間であり，その次元が $1$ なので成り立つ．

\medskip\noindent
(2) の証明．Duality for Schemes，補題
\ref{duality-lemma-more-base-change} による $a$ の基底変換があるので成り立つ．
Duality for Schemes，備考
\ref{duality-remark-relative-dualizing-complex} の議論も参照されたい．
\end{proof}

\begin{lemma}
\label{curves-lemma-closed-immersion-dim-1-CM}
$X$ を次元 $\leq 1$ の，体 $k$ 上の固有概形とする．
$i : Y \to X$ を閉埋め込みとする．
$\omega_X^\bullet$, $\omega_X$, $\omega_Y^\bullet$, $\omega_Y$ を
補題 \ref{curves-lemma-duality-dim-1} のものとする．このとき
\begin{enumerate}
\item $\omega_Y^\bullet = R\SheafHom(\mathcal{O}_Y, \omega_X^\bullet)$ である．
\item $\omega_Y = \SheafHom(\mathcal{O}_Y, \omega_X)$ かつ
$i_*\omega_Y = \SheafHom_{\mathcal{O}_X}(i_*\mathcal{O}_Y, \omega_X)$ である．
\end{enumerate}
\end{lemma}

\begin{proof}
$g : Y \to \Spec(k)$ と $f : X \to \Spec(k)$ をそれぞれ構造射と書く．
このとき $g = f \circ i$ である．$a, b, c$ をそれぞれ
Duality for Schemes，補題
\ref{duality-lemma-twisted-inverse-image} による $f, g, i$ の右随伴と書く．
右随伴の一意性により $b = c \circ a$ であり，ここでは
$Rg_* = Rf_* \circ Ri_*$ を用いた．
補題 \ref{curves-lemma-duality-dim-1} の証明では
$\omega_X^\bullet = a(\mathcal{O}_{\Spec(k)})$ および
$\omega_Y^\bullet = b(\mathcal{O}_{\Spec(k)})$ とおいた．
したがって $\omega_Y^\bullet = c(\omega_X^\bullet)$ であり，
Duality for Schemes，補題
\ref{duality-lemma-twisted-inverse-image-closed} により (1) が従う．
$\omega_X = H^{-1}(\omega_X^\bullet)$ および
$\omega_Y = H^{-1}(\omega_Y^\bullet)$ なので，
$\omega_Y = \SheafHom(\mathcal{O}_Y, \omega_X)$ を得る．
さらに Duality for Schemes，補題
\ref{duality-lemma-sheaf-with-exact-support-ext} により
$i_*\omega_Y = \SheafHom_{\mathcal{O}_X}(i_*\mathcal{O}_Y, \omega_X)$
が従う．
\end{proof}

\begin{lemma}
\label{curves-lemma-closed-subscheme-reduced-gorenstein}
$X$ を体 $k$ 上の固有概形で，Gorenstein，簡約，かつ次元 $1$ の
等次元なものとする．$i : Y \to X$ を，次元 $1$ の等次元な
簡約閉部分概形とする．$j : Z \to X$ を $X \setminus Y$ の
概形論的閉包とする．このとき
\begin{enumerate}
\item $Y$ と $Z$ は Cohen--Macaulay である．
\item $\mathcal{I} \subset \mathcal{O}_X$，
$\mathcal{J} \subset \mathcal{O}_X$ をそれぞれ $Y$，$Z$ の
$X$ におけるイデアル層とすると，
$$
\mathcal{I} = i_*\mathcal{I}'
\quad\text{かつ}\quad
\mathcal{J} = j_*\mathcal{J}'
$$
である．ここで $\mathcal{I}' \subset \mathcal{O}_Z$，
$\mathcal{J}' \subset \mathcal{O}_Y$ は，$Y \cap Z$ の，それぞれ
$Z$，$Y$ におけるイデアル層である．
\item $\omega_Y = \mathcal{J}'(i^*\omega_X)$ かつ
$i_*(\omega_Y) = \mathcal{J}\omega_X$ である．
\item $\omega_Z = \mathcal{I}'(i^*\omega_X)$ かつ
$i_*(\omega_Z) = \mathcal{I}\omega_X$ である．
\item 次の短完全列がある：
\begin{align*}
0 \to \omega_X \to i_*i^*\omega_X \oplus j_*j^*\omega_X \to
\mathcal{O}_{Y \cap Z} \to 0 \\
0 \to i_*\omega_Y \to \omega_X \to j_*j^*\omega_X \to 0 \\
0 \to j_*\omega_Z \to \omega_X \to i_*i^*\omega_X \to 0 \\
0 \to i_*\omega_Y \oplus j_*\omega_Z \to \omega_X \to
\mathcal{O}_{Y \cap Z} \to 0 \\
0 \to \omega_Y \to i^*\omega_X \to \mathcal{O}_{Y \cap Z} \to 0 \\
0 \to \omega_Z \to j^*\omega_X \to \mathcal{O}_{Y \cap Z} \to 0
\end{align*}
\end{enumerate}
ここで $\omega_X$, $\omega_Y$, $\omega_Z$ は補題
\ref{curves-lemma-duality-dim-1} のものである．
\end{lemma}

\begin{proof}
簡約な $1$ 次元 Noether 概形は Cohen--Macaulay なので (1) が成り立つ．
$X$ は簡約だから，概形論的に $X = Y \cup Z$ である．
Morphisms，補題
\ref{morphisms-lemma-scheme-theoretic-union} の記法と主張により，短完全列
$$
0 \to \mathcal{O}_X \to
\mathcal{O}_Y \oplus \mathcal{O}_Z \to \mathcal{O}_{Y \cap Z} \to 0
$$
を得る．
$\mathcal{J} = \Ker(\mathcal{O}_X \to \mathcal{O}_Z)$，
$\mathcal{J}' = \Ker(\mathcal{O}_Y \to \mathcal{O}_{Y \cap Z})$，
$\mathcal{I} = \Ker(\mathcal{O}_X \to \mathcal{O}_Y)$，および
$\mathcal{I}' = \Ker(\mathcal{O}_Z \to \mathcal{O}_{Y \cap Z})$
なので，図式追跡から (2) が従う．
$\mathcal{I} + \mathcal{J}$ は $Y \cap Z$ のイデアル層であり，
$\mathcal{I} \cap \mathcal{J} = 0$ であることに注意しよう．
したがって次の完全列がある：
\begin{align*}
0 \to \mathcal{O}_X \to \mathcal{O}_Y \oplus \mathcal{O}_Z \to
\mathcal{O}_{Y \cap Z} \to 0 \\
0 \to \mathcal{J} \to \mathcal{O}_X \to \mathcal{O}_Z \to 0 \\
0 \to \mathcal{I} \to \mathcal{O}_X \to \mathcal{O}_Y \to 0 \\
0 \to \mathcal{J} \oplus \mathcal{I} \to \mathcal{O}_X \to
\mathcal{O}_{Y \cap Z} \to 0 \\
0 \to \mathcal{J}' \to \mathcal{O}_Y \to \mathcal{O}_{Y \cap Z} \to 0 \\
0 \to \mathcal{I}' \to \mathcal{O}_Z \to \mathcal{O}_{Y \cap Z} \to 0
\end{align*}
$X$ は Gorenstein なので，$\omega_X$ は可逆 $\mathcal{O}_X$-加群である
（Duality for Schemes，補題 \ref{duality-lemma-gorenstein}）．
$Y \cap Z$ は次元 $0$ なので
$\omega_X|_{Y \cap Z} \cong \mathcal{O}_{Y \cap Z}$ である．
したがって (3) と (4) を示せば，上の短完全列を可逆加群
$\omega_X$ とテンソルすることにより，補題の短完全列を得る．
対称性により (3) を示せば十分であり，さらに (2) により
$i_*(\omega_Y) = \mathcal{J}\omega_X$ を示せば十分である．

\medskip\noindent
補題 \ref{curves-lemma-closed-immersion-dim-1-CM} により
$i_*\omega_Y = \SheafHom_{\mathcal{O}_X}(i_*\mathcal{O}_Y, \omega_X)$
である．再び $\omega_X$ が可逆であることを用いると，最終的には
$\SheafHom_{\mathcal{O}_X}(\mathcal{O}_X/\mathcal{I}, \mathcal{O}_X)$
が $\mathcal{J}$ へ，$1$ での評価により同型に写ることを示せば十分である．
言い換えれば，$\mathcal{J}$ が $\mathcal{I}$ の零化イデアルであることを
示せばよい．これは上の議論から従う．
\end{proof}







\section{Riemann--Roch}
\label{curves-section-Riemann-Roch}

\noindent
$k$ を体とする．$X$ を次元 $\leq 1$ の $k$ 上の固有概形とする．
Varieties，節 \ref{varieties-section-divisors-curves} では，定数階数の
局所自由 $\mathcal{O}_X$-加群 $\mathcal{E}$ の次数を公式
\begin{equation}
\label{curves-equation-degree}
\deg(\mathcal{E}) =
\chi(X, \mathcal{E}) - \text{rank}(\mathcal{E})\chi(X, \mathcal{O}_X)
\end{equation}
で定義した．Varieties，定義
\ref{varieties-definition-degree-invertible-sheaf} を参照されたい．
Chow Homology の章では，$\mathcal{E}$ の第一 Chern 類をサイクル上の
作用として定義し（Chow Homology，節
\ref{chow-section-intersecting-chern-classes}），
\begin{equation}
\label{curves-equation-degree-c1}
\deg(\mathcal{E}) = \deg(c_1(\mathcal{E}) \cap [X]_1)
\end{equation}
を証明した．Chow Homology，補題
\ref{chow-lemma-degree-vector-bundle} を参照されたい．
(\ref{curves-equation-degree}) と (\ref{curves-equation-degree-c1}) を組み合わせると，
Riemann--Roch 公式の最初の形
\begin{equation}
\label{curves-equation-rr}
\chi(X, \mathcal{E}) =
\deg(c_1(\mathcal{E}) \cap [X]_1) +
\text{rank}(\mathcal{E})\chi(X, \mathcal{O}_X)
\end{equation}
を得る．$\mathcal{L}$ が可逆 $\mathcal{O}_X$-加群ならば，
Varieties，定義 \ref{varieties-definition-intersection-number} で定義された
数値的交点数 $(\mathcal{L} \cdot X)$ も考えられる．しかし
Varieties，補題
\ref{varieties-lemma-intersection-numbers-and-degrees-on-curves} により
\begin{equation}
\label{curves-equation-numerical-degree}
(\mathcal{L} \cdot X) = \deg(\mathcal{L})
\end{equation}
なので，これは新しいものを与えない．$\mathcal{L}$ が豊富ならば，
この整数は正であり，$X$ の $\mathcal{L}$ に関する次数と呼ばれる：
\begin{equation}
\label{curves-equation-degree-X}
\deg_\mathcal{L}(X) = (\mathcal{L} \cdot X) = \deg(\mathcal{L})
\end{equation}
Varieties，定義 \ref{varieties-definition-degree} を参照されたい．

\medskip\noindent
真の Riemann--Roch 定理を得るには，$\chi(X, \mathcal{O}_X)$ を
$X$ 上の標準的な零サイクルの次数として書きたい．完全に一般的な形については
\cite{F} を参照されたい．ここでは双対性を用いて，$X$ が Gorenstein である
場合の公式を得る．ただし，これはある意味では便法である
（例えば，この方法は高次元では使えない）．

\medskip\noindent
まず補題 \ref{curves-lemma-duality-dim-1} と \ref{curves-lemma-duality-dim-1-CM} を用いて，
$\mathcal{O}_X$ のオイラー標数と，双対化複体または双対化加群の
オイラー標数との関係を得る．

\begin{lemma}
\label{curves-lemma-euler}
$X$ を次元 $\leq 1$ の，体 $k$ 上の固有概形とする．
補題 \ref{curves-lemma-duality-dim-1} の $\omega_X^\bullet$ と $\omega_X$ に対し，
$$
\chi(X, \mathcal{O}_X) = \chi(X, \omega_X^\bullet)
$$
が成り立つ．$X$ が Cohen--Macaulay かつ次元 $1$ の等次元ならば，
$$
\chi(X, \mathcal{O}_X) = - \chi(X, \omega_X)
$$
が成り立つ．
\end{lemma}

\begin{proof}
第一の公式の右辺を次のように定義する：
$$
\chi(X, \omega_X^\bullet) =
\sum\nolimits_{i \in \mathbf{Z}} (-1)^i\dim_k H^i(X, \omega_X^\bullet)
$$
これは $\omega_X^\bullet$ が $D^b_{\textit{Coh}}(\mathcal{O}_X)$ に属するため
良定義であるが，さらに
$$
H^i(X, \omega_X^\bullet) =
\Ext^i(\mathcal{O}_X, \omega_X^\bullet) =
H^{-i}(X, \mathcal{O}_X)
$$
が常に有限次元で，$i = 0, -1$ の場合にしか非零にならないことからも
分かる．これはもちろん第一の公式も証明する．第二の公式は第一の公式と，
Cohen--Macaulay の場合に $\omega_X^\bullet = \omega_X[1]$ となることから
従う．補題 \ref{curves-lemma-duality-dim-1-CM} を参照されたい．
\end{proof}

\noindent
補題 \ref{curves-lemma-euler} を用いて，$\chi(X, \mathcal{O}_X)$ の求める公式を，
$\omega_X$ が可逆，すなわち $X$ が Gorenstein である場合に得る．
その主張は，$-1/2$ と $\omega_X$ の第一 Chern 類との積を，
サイクル $[X]_1$（$X$ に付随するもの）と交わらせたものが，$X$ 上の半整数係数をもつ
自然な零サイクルであり，その次数が $\chi(X, \mathcal{O}_X)$ になるという
ものである．Riemann--Roch の主張に現れる分数は避けられない．

\begin{lemma}[Riemann--Roch]
\label{curves-lemma-rr}
$X$ を体 $k$ 上の固有概形で，Gorenstein かつ次元 $1$ の等次元なものとする．
$\omega_X$ を補題 \ref{curves-lemma-duality-dim-1} のものとする．このとき
\begin{enumerate}
\item $\omega_X$ は可逆 $\mathcal{O}_X$-加群である．
\item $\deg(\omega_X) = -2\chi(X, \mathcal{O}_X)$ である．
\item 定数階数の局所自由 $\mathcal{O}_X$-加群 $\mathcal{E}$ に対し，
$$
\chi(X, \mathcal{E}) = \deg(\mathcal{E}) -
\textstyle{\frac{1}{2}} \text{rank}(\mathcal{E}) \deg(\omega_X)
$$
であり，
$\dim_k(H^i(X, \mathcal{E})) =
\dim_k(H^{1 - i}(X, \mathcal{E}^\vee \otimes_{\mathcal{O}_X} \omega_X))$
がすべての $i \in \mathbf{Z}$ に対して成り立つ．
\end{enumerate}
\end{lemma}

\noindent
非特異（正規）曲線は Gorenstein である．Duality for Schemes，補題
\ref{duality-lemma-regular-gorenstein} を参照されたい．

\begin{proof}
Gorenstein 概形は Cohen--Macaulay であることを思い出そう
（Duality for Schemes，補題 \ref{duality-lemma-gorenstein-CM}）．
したがって $\omega_X$ は $X$ 上の双対化加群である．
補題 \ref{curves-lemma-duality-dim-1-CM} を参照されたい．
Gorenstein 性の定義から，双対化層が可逆であることはほぼ直ちに従う．
Duality for Schemes，節 \ref{duality-section-gorenstein} を参照されたい．
(\ref{curves-equation-rr}) を $\omega_X$ に適用すると，
$$
\chi(X, \omega_X) = 
\deg(c_1(\omega_X) \cap [X]_1) + \chi(X, \mathcal{O}_X)
$$
を得る．補題 \ref{curves-lemma-euler} と組み合わせると，
$$
2\chi(X, \mathcal{O}_X) = - \deg(c_1(\omega_X) \cap [X]_1) = - \deg(\omega_X)
$$
を得る．第二の等号は (\ref{curves-equation-degree-c1}) による．
これを $\mathcal{E}$ に対する (\ref{curves-equation-rr}) に戻せば，
補題の表示公式を得る．次元についての対称性は $X$ に対する双対性から従う．
備考 \ref{curves-remark-rework-duality-locally-free} を参照されたい．
\end{proof}





\section{いくつかの消滅結果}
\label{curves-section-vanishing}

\noindent
本節では，いくつかの非常に弱い消滅結果を述べる．もう少し多く，かつ
より興味深い結果については節 \ref{curves-section-more-vanishing} を参照されたい．

\begin{lemma}
\label{curves-lemma-automatic}
$k$ を体とする．$X$ を $k$ 上の次元 $1$ の固有概形で，
$H^0(X, \mathcal{O}_X) = k$ を満たすものとする．このとき $X$ は連結で，
Cohen--Macaulay かつ次元 $1$ の等次元である．
\end{lemma}

\begin{proof}
$\Gamma(X, \mathcal{O}_X) = k$ は非自明な冪等元をもたないので，
$X$ は連結である．これだけで $X$ が次元 $1$ の等次元であることも分かる
（次元 $0$ の既約成分は連結成分になってしまう）．
$\mathcal{I} \subset \mathcal{O}_X$ を，閉点に支えをもつ最大の
連接部分加群とする．この $\mathcal{I}$ は存在し
（Divisors，補題 \ref{divisors-lemma-remove-embedded-points}），
大域生成される
（Varieties，補題 \ref{varieties-lemma-chi-tensor-finite}）．
$1 \in \Gamma(X, \mathcal{O}_X)$ は $\mathcal{I}$ の切断ではないので，
$\mathcal{I} = 0$ を得る．したがって $X$ は埋め込まれた点をもたない
（Divisors，補題 \ref{divisors-lemma-remove-embedded-points}）．
ゆえに Divisors，補題 \ref{divisors-lemma-S1-no-embedded} により
$X$ は $(S_1)$ を満たす．従って $X$ は Cohen--Macaulay である．
\end{proof}

\noindent
本節では以下の状況で議論する．

\begin{situation}
\label{curves-situation-Cohen-Macaulay-curve}
ここで $k$ は体，$X$ は $k$ 上の次元 $1$ の固有概形であり，
$H^0(X, \mathcal{O}_X) = k$ を満たす．
\end{situation}

\noindent
補題 \ref{curves-lemma-automatic} により，概形 $X$ は Cohen--Macaulay かつ
次元 $1$ の等次元である．補題 \ref{curves-lemma-duality-dim-1} および
\ref{curves-lemma-duality-dim-1-CM} で論じた双対化加群 $\omega_X$ は
非消滅な $H^1$ をもつ．実際，
$\dim_k H^1(X, \omega_X) = \dim_k H^0(X, \mathcal{O}_X) = 1$ である．
$\omega_X$ より少しでも「正」であるものは $H^1$ が消滅することが分かる．

\begin{lemma}
\label{curves-lemma-vanishing}
状況 \ref{curves-situation-Cohen-Macaulay-curve} にあるとする．完全列
$$
\omega_X \to \mathcal{F} \to \mathcal{Q} \to 0
$$
が連接 $\mathcal{O}_X$-加群の列として与えられ，$H^1(X, \mathcal{Q}) = 0$
（例えば $\dim(\text{Supp}(\mathcal{Q})) = 0$ ならばそうである）
とする．このとき $H^1(X, \mathcal{F}) = 0$ または
$\mathcal{F} = \omega_X \oplus \mathcal{Q}$ である．
\end{lemma}

\begin{proof}
（括弧内の主張は Cohomology of Schemes，補題
\ref{coherent-lemma-coherent-support-dimension-0} から従う．）
$H^0(X, \mathcal{O}_X) = k$ は $H^1(X, \omega_X)$ と双対であるから
（節 \ref{curves-section-Riemann-Roch} を参照），
$\dim H^1(X, \omega_X) = 1$ である．層 $\omega_X$ は，関手
$\mathcal{F} \mapsto \Hom_k(H^1(X, \mathcal{F}), k)$ を連接
$\mathcal{O}_X$-加群の圏上で表現する
（Duality for Schemes，補題
\ref{duality-lemma-dualizing-module-proper-over-A}）．
補題の主張にある完全列を考え，$H^1(X, \mathcal{F}) \not = 0$ と仮定する．
$H^1(X, \mathcal{Q}) = 0$ なので，
$H^1(X, \omega_X) \to H^1(X, \mathcal{F})$ は同型である．
上で述べた $\omega_X$ の普遍性により，写像 $\mathcal{F} \to \omega_X$ が
存在し，$H^1$ 上でこの同型の逆を誘導する．合成
$\omega_X \to \mathcal{F} \to \omega_X$ は（普遍性により）恒等写像である．
これで補題が証明された．
\end{proof}

\begin{lemma}
\label{curves-lemma-vanishing-twist}
状況 \ref{curves-situation-Cohen-Macaulay-curve} にあるとする．
$\mathcal{L}$ を，大域生成され $\mathcal{O}_X$ と同型でない
可逆 $\mathcal{O}_X$-加群とする．このとき
$H^1(X, \omega_X \otimes \mathcal{L}) = 0$ である．
\end{lemma}

\begin{proof}
節 \ref{curves-section-Riemann-Roch} で論じた双対性により，
$H^0(X, \mathcal{L}^{\otimes - 1}) = 0$ を示せばよい．
そうでなければ，大域切断 $t$ を $\mathcal{L}^{\otimes - 1}$ から，
大域切断 $s$ を $\mathcal{L}$ から，$st \not = 0$ となるように選べる．
しかし $st$ は $1$ の定数倍である．これは
$H^0(X, \mathcal{O}_X) = k$ という仮定による．従って
$\mathcal{L} \cong \mathcal{O}_X$ となり，矛盾する．
\end{proof}

\begin{lemma}
\label{curves-lemma-globally-generated}
状況 \ref{curves-situation-Cohen-Macaulay-curve} にあるとする．完全列
$$
\omega_X \to \mathcal{F} \to \mathcal{Q} \to 0
$$
が連接 $\mathcal{O}_X$-加群の列として与えられ，
$\dim(\text{Supp}(\mathcal{Q})) = 0$ かつ
$\dim_k H^0(X, \mathcal{Q}) \geq 2$ とする．さらに，
非零部分加群 $\mathcal{Q}' \subset \mathcal{F}$ で，
$\mathcal{Q}' \to \mathcal{Q}$ が単射となるものは存在しないとする．
このとき，$\mathcal{F}$ の大域切断で生成される部分加群は
$\mathcal{Q}$ へ全射する．
\end{lemma}

\begin{proof}
$\mathcal{F}' \subset \mathcal{F}$ を，大域切断と
$\omega_X \to \mathcal{F}$ の像で生成される部分加群とする．
$\dim_k H^0(X, \mathcal{Q}) \geq 2$ および
$\dim_k H^1(X, \omega_X) = \dim_k H^0(X, \mathcal{O}_X) = 1$ なので，
$\mathcal{F}' \to \mathcal{Q}$ は零でなく，
$\omega_X \to \mathcal{F}'$ は同型でない．従って
$H^1(X, \mathcal{F}') = 0$ である．これは補題
\ref{curves-lemma-vanishing} と $\mathcal{F}$ に関する仮定による．短完全列
$$
0 \to \mathcal{F}' \to \mathcal{F} \to
\mathcal{Q}/\Im(\mathcal{F}' \to \mathcal{Q}) \to 0
$$
を考える．右辺の商が非零ならば，
$H^0(X, \mathcal{F})$ が $H^0(X, \mathcal{F}')$ より大きくなるため矛盾する．
\end{proof}

\noindent
大域生成に関する例を一つ挙げる．

\begin{lemma}
\label{curves-lemma-globally-generated-curve}
状況 \ref{curves-situation-Cohen-Macaulay-curve} にあり，$X$ が整であると仮定する．
$0 \to \omega_X \to \mathcal{F} \to \mathcal{Q} \to 0$ を連接
$\mathcal{O}_X$-加群の短完全列で，$\mathcal{F}$ は捩れなし，
$\dim(\text{Supp}(\mathcal{Q})) = 0$，かつ
$\dim_k H^0(X, \mathcal{Q}) \geq 2$ を満たすものとする．このとき
$\mathcal{F}$ は大域生成される．
\end{lemma}

\begin{proof}
大域切断で生成される部分加群 $\mathcal{F}'$ を考える．
補題 \ref{curves-lemma-globally-generated} により $\mathcal{F}' \to \mathcal{Q}$ は
全射であり，特に $\mathcal{F}' \not = 0$ である．$X$ は曲線なので，
$\mathcal{F}' \subset \mathcal{F}$ は階数 $1$ の層の包含である．従って
$\mathcal{Q}' = \mathcal{F}/\mathcal{F}'$ は有限個の点に支えをもつ．
矛盾を得るため，$\mathcal{Q}'$ が非零であると仮定する．すると
$H^1(X, \mathcal{F}') \not = 0$ である．普遍性により非零写像
$\mathcal{F}' \to \omega_X$ を得る
（Duality for Schemes，補題
\ref{duality-lemma-dualizing-module-proper-over-A}）．合成
$\mathcal{F}' \to \omega_X \to \mathcal{F}$ の像は大域切断で生成されるので，
$\mathcal{F}'$ に含まれる．従って非零自己写像
$\mathcal{F}' \to \mathcal{F}'$ を得る．
$\mathcal{F}'$ は固有曲線上の階数 $1$ の捩れなし層なので，
これは自己同型でなければならない（詳細は省略する）．しかし，これは
$\mathcal{F}'$ が $\omega_X \subset \mathcal{F}$ に含まれることを意味し，
$\mathcal{F}' \to \mathcal{Q}$ の全射性に矛盾する．
\end{proof}

\begin{lemma}
\label{curves-lemma-tensor-omega-with-globally-generated-invertible}
状況 \ref{curves-situation-Cohen-Macaulay-curve} にあるとする．
$\mathcal{L}$ を非常に豊富な可逆 $\mathcal{O}_X$-加群で，
$\deg(\mathcal{L}) \geq 2$ を満たすものとする．このとき
$\omega_X \otimes_{\mathcal{O}_X} \mathcal{L}$ は大域生成される．
\end{lemma}

\begin{proof}
$k$ が代数閉体であると仮定する．$x \in X$ を閉点とする．
$C_i \subset X$ を既約成分とし，各 $i$ に対して
$x_i \in C_i$ をその一般点とする．Varieties，補題
\ref{varieties-lemma-very-ample-vanish-at-point} により，切断
$s \in H^0(X, \mathcal{L})$ を選べる．この $s$ は $x$ では消えるが，
$x_i$ ではすべての $i$ に対して消えない．対応する加群写像
$s : \mathcal{O}_X \to \mathcal{L}$ は単射であり，余核
$\mathcal{Q}$ は有限個の点に支えをもち，$H^0(X, \mathcal{Q}) \geq 2$
を満たす．対応する完全列
$$
0 \to \omega_X \to \omega_X \otimes \mathcal{L} \to
\omega_X \otimes \mathcal{Q} \to 0
$$
を考える．補題 \ref{curves-lemma-globally-generated} により，大域切断で生成される
加群は $\omega_X \otimes \mathcal{Q}$ へ全射する．$x$ は任意だったので，
これで補題が証明された．一部の詳細は省略する．

\medskip\noindent
$k$ が代数閉体でない場合を，代数閉体の場合に帰着する．
以下の証明は読み飛ばしてもよい．代数閉包 $\overline{k}$ を $k$ に対して選び，
基底変換 $X_{\overline{k}}$ を考える．
$X_{\overline{k}} \to \Spec(\overline{k})$ が状況
\ref{curves-situation-Cohen-Macaulay-curve} の例であることを確認しよう．
平坦基底変換（Cohomology of Schemes，補題
\ref{coherent-lemma-flat-base-change-cohomology}）により
$H^0(X_{\overline{k}}, \mathcal{O}) = \overline{k}$ である．
概形 $X_{\overline{k}}$ は $\overline{k}$ 上固有であり
（Morphisms，補題 \ref{morphisms-lemma-base-change-proper}），
次元 $1$ の等次元である
（Morphisms，補題
\ref{morphisms-lemma-dimension-fibre-after-base-change}）．
補題 \ref{curves-lemma-sanity-check-duality} により，$\omega_X$ の
$X_{\overline{k}}$ への引戻しは $X_{\overline{k}}$ の双対化加群である．
$\mathcal{L}$ の $X_{\overline{k}}$ への引戻しは非常に豊富である
（Morphisms，補題 \ref{morphisms-lemma-very-ample-base-change}）．
$\mathcal{L}$ の $X_{\overline{k}}$ への引戻しの次数は，
$\mathcal{L}$ の $X$ 上での次数に等しい
（Varieties，補題 \ref{varieties-lemma-degree-base-change}）．
最後に，$\omega_X \otimes \mathcal{L}$ が大域生成されることと，
その基底変換が大域生成されることは同値である
（Varieties，補題 \ref{varieties-lemma-globally-generated-base-change}）．
以上により，結果は基礎体が代数閉体である場合の結果から従う．
\end{proof}




\section{非常に豊富な可逆層}
\label{curves-section-very-ample}

\noindent
代数閉体上有限型な概形について，可逆加群 $\mathcal{L}$ が非常に豊富である
ためによく用いられる判定法は，概形 $X$ の点と接ベクトルを
$\mathcal{L}$ の切断が分離するというものである
（Varieties，節 \ref{varieties-section-separating-points-tangent-vectors}）．
曲線に対する別の判定法を述べる．Varieties，小節
\ref{varieties-subsection-regularity} と比較されたい．

\begin{lemma}
\label{curves-lemma-criterion-very-ample}
$k$ を体とする．$X$ を $k$ 上の次元 $1$ の固有概形で，
$H^0(X, \mathcal{O}_X) = k$ を満たすものとする．
$\mathcal{L}$ を可逆 $\mathcal{O}_X$-加群とし，次を仮定する：
\begin{enumerate}
\item $\mathcal{L}$ は正則な大域切断をもつ．
\item $H^1(X, \mathcal{L}) = 0$ である．
\item $\mathcal{L}$ は豊富である．
\end{enumerate}
このとき $\mathcal{L}^{\otimes 6}$ は $X$ 上 $k$ に関して非常に豊富である．
\end{lemma}

\begin{proof}
$s$ を $\mathcal{L}$ の正則な大域切断とする．
$i : Z = Z(s) \to X$ を $s$ の零点概形とする．Divisors，節
\ref{divisors-section-effective-Cartier-invertible} を参照されたい．
条件 (3) により $Z \not = \emptyset$ である（細部を一つ省略する）．
短完全列
$$
0 \to \mathcal{O}_X \xrightarrow{s} \mathcal{L} \to
i_*(\mathcal{L}|_Z) \to 0
$$
を考える．$\mathcal{L}$ とテンソルすると，
$$
0 \to \mathcal{L} \to \mathcal{L}^{\otimes 2} \to
i_*(\mathcal{L}^{\otimes 2}|_Z) \to 0
$$
を得る．$Z$ は次元 $0$ であることに注意しよう
（Divisors，補題 \ref{divisors-lemma-effective-Cartier-makes-dimension-drop}）．
従って Artin 環のスペクトルであり
（Varieties，補題 \ref{varieties-lemma-algebraic-scheme-dim-0}），
ゆえに $\mathcal{L}|_Z \cong \mathcal{O}_Z$ である
（Algebra，補題 \ref{algebra-lemma-locally-free-semi-local-free}）．
この短完全列から
$H^1(X, \mathcal{L}^{\otimes 2}) = 0$ も従う
（例えば Varieties，補題 \ref{varieties-lemma-chi-tensor-finite} を用いて，
右端の該当箇所が消滅することを見ればよい）．
$n \geq 1$ に関する帰納法と完全列
$$
0 \to \mathcal{L}^{\otimes n} \xrightarrow{s}
\mathcal{L}^{\otimes n + 1} \to
i_*(\mathcal{L}^{\otimes n + 1}|_Z) \to 0
$$
を用いると，$H^1(X, \mathcal{L}^{\otimes n}) = 0$ が $n > 0$ に対して
成り立つこと，および大域切断 $t_{n + 1}$ が
$\mathcal{L}^{\otimes n + 1}$ に存在して，
$\mathcal{L}^{\otimes n + 1}|_Z \cong \mathcal{O}_Z$ の自明化を与えることが分かる．

\medskip\noindent
乗法写像
$$
\mu_n :
H^0(X, \mathcal{L}) \otimes_k H^0(X, \mathcal{L}^{\otimes n})
\oplus
H^0(X, \mathcal{L}^{\otimes 2}) \otimes_k H^0(X, \mathcal{L}^{\otimes n - 1})
\longrightarrow
H^0(X, \mathcal{L}^{\otimes n + 1})
$$
を考える．これは $n \geq 3$ に対して全射であると主張する．
これを見るため，短完全列
$$
0 \to \mathcal{L}^{\otimes n} \xrightarrow{s}
\mathcal{L}^{\otimes n + 1} \to i_*(\mathcal{L}^{\otimes n + 1}|_Z) \to 0
$$
を考える．この列の左側から来る $\mathcal{L}^{\otimes n + 1}$ の切断は
$\mu_n$ の像に含まれる．一方，
$H^0(\mathcal{L}^{\otimes 2}) \to H^0(\mathcal{L}^{\otimes 2}|_Z)$ は
全射であり（上を参照），$t_{n - 1}$ は
$\mathcal{L}^{\otimes n - 1}|_Z$ の自明化に写るので，
$\mu_n(H^0(X, \mathcal{L}^{\otimes 2}) \otimes t_{n - 1})$ は，
$H^0(X, \mathcal{L}^{\otimes n + 1})$ の部分空間であって，
$\mathcal{L}^{\otimes n + 1}|_Z$ の大域切断へ全射するものを与える．
これで主張が証明された．

\medskip\noindent
前段落の主張から，次数付き $k$-代数
$$
S = \bigoplus\nolimits_{n \geq 0} H^0(X, \mathcal{L}^{\otimes n})
$$
は次数 $0, 1, 2, 3$ の部分によって $k$ 上生成されることが分かる．
$X = \text{Proj}(S)$ であることを思い出そう
（Morphisms，補題 \ref{morphisms-lemma-proper-ample-is-proj}）．
従って $S^{(6)} = \bigoplus_{n} S_{6n}$ は次数 $1$ で生成される．
これは，望んだ通り $\mathcal{L}^{\otimes 6}$ が非常に豊富であることを意味する．
\end{proof}

\begin{lemma}
\label{curves-lemma-criterion-very-ample-bis}
$k$ を体とする．$X$ を $k$ 上の次元 $1$ の固有概形で，
$H^0(X, \mathcal{O}_X) = k$ を満たすものとする．
$\mathcal{L}$ を可逆 $\mathcal{O}_X$-加群とし，次を仮定する：
\begin{enumerate}
\item $\mathcal{L}$ は大域生成される．
\item $H^1(X, \mathcal{L}) = 0$ である．
\item $\mathcal{L}$ は豊富である．
\end{enumerate}
このとき $\mathcal{L}^{\otimes 2}$ は $X$ 上 $k$ に関して非常に豊富である．
\end{lemma}

\begin{proof}
基底 $s_0, \ldots, s_n$ を，$H^0(X, \mathcal{L}^{\otimes 2})$ の
$k$ 上の基底として選ぶ．性質 (1) により
$\mathcal{L}^{\otimes 2}$ は大域生成され，射
$$
\varphi_{\mathcal{L}^{\otimes 2}, (s_0, \ldots, s_n)} :
X \longrightarrow \mathbf{P}^n_k
$$
を得る．Constructions，節
\ref{constructions-section-projective-space} を参照されたい．
補題の主張は，この射が閉埋め込みであるということである．
これを確かめるため，$k$ をその代数閉包で置き換えてよい
（Descent，補題 \ref{descent-lemma-descending-property-closed-immersion}）．
従って $k$ は代数閉体であると仮定してよい．

\medskip\noindent
$k$ が代数閉体であると仮定する．各一般点 $\eta_i \in X$ に対し，
$V_i \subset H^0(X, \mathcal{L})$ を，$k$-ベクトル部分空間であって
$\eta_i$ で消える切断からなるものとする．$\mathcal{L}$ は大域生成されるので，
$V_i \not = H^0(X, \mathcal{L})$ である．$X$ は有限個の既約成分しか
もたず，$k$ は無限体なので，切断 $s \in H^0(X, \mathcal{L})$ であって，
$\eta_i$ でどの $i$ に対しても消えないものを見つけられる．このとき $s$ は
$\mathcal{L}$ の正則切断である（補題 \ref{curves-lemma-automatic} により
$X$ は Cohen--Macaulay であり，従って $\mathcal{L}$ は埋め込まれた
随伴点をもたないからである）．

\medskip\noindent
特に，補題 \ref{curves-lemma-criterion-very-ample} の証明で与えた主張はすべて，
この $s$ に対して成り立つ．さらに $\mathcal{L}$ は大域生成されるので，
大域切断 $t \in H^0(X, \mathcal{L})$ であって，$t|_Z$ が消えないものを
見つけられる（$Z$ の点が有限個であることを用い，上と同様に論じる）．
すると補題 \ref{curves-lemma-criterion-very-ample} の証明では，$t$ を用いてさらに
乗法写像
$$
\mu_n :
H^0(X, \mathcal{L}) \otimes_k H^0(X, \mathcal{L}^{\otimes 2})
\longrightarrow
H^0(X, \mathcal{L}^{\otimes 3})
$$
が全射であることを示せる．従って
$$
S = \bigoplus\nolimits_{n \geq 0} H^0(X, \mathcal{L}^{\otimes n})
$$
は次数 $0, 1, 2$ の部分によって $k$ 上生成される．
補題 \ref{curves-lemma-criterion-very-ample} の証明と同様に論じると，
$S^{(2)} = \bigoplus_{n} S_{2n}$ は次数 $1$ で生成される．
これは望んだ通り，$\mathcal{L}^{\otimes 2}$ が非常に豊富であることを意味する．
一部の詳細は省略する．
\end{proof}









\section{曲線の種数}
\label{curves-section-genus}

\noindent
$X$ が体 $k$ 上の滑らかな射影的幾何学的既約曲線ならば，以前に $X$ の
種数を $H^1(X, \mathcal{O}_X)$ の次元として定義した
（Picard Schemes of Curves，定義 \ref{pic-definition-genus}）．
この場合には $H^0(X, \mathcal{O}_X) = k$ であることに注意しよう
（Varieties，補題 \ref{varieties-lemma-regular-functions-proper-variety}）．
これを次のように一般化する．

\begin{definition}
\label{curves-definition-genus}
$k$ を体とする．$X$ を $k$ 上の次元 $1$ の固有概形で，
$H^0(X, \mathcal{O}_X) = k$ を満たすものとする．このとき
$X$ の {\it 種数}とは $g = \dim_k H^1(X, \mathcal{O}_X)$ のことである．
\end{definition}

\noindent
これは $X$ の {\it 算術種数}と呼ばれることもある．文献では，固有曲線
$X$ の，$k$ 上の算術種数を
$$
p_a(X) = 1 - \chi(X, \mathcal{O}_X) =
1 - \dim_k H^0(X, \mathcal{O}_X) + \dim_k H^1(X, \mathcal{O}_X)
$$
と定義することもある．これは適用できる場合には本書の定義と一致する．
実際，$H^0(X, \mathcal{O}_X) = k$ を仮定しているからである．ただし，
\begin{enumerate}
\item $p_a(X)$ は負になり得る．
\item $p_a(X)$ は基礎体 $k$ に依存するので，$p_a(X/k)$ と書くべきである．
\end{enumerate}
例えば $k = \mathbf{Q}$ かつ $X = \mathbf{P}^1_{\mathbf{Q}(i)}$ ならば，
$p_a(X/\mathbf{Q}) = -1$ かつ $p_a(X/\mathbf{Q}(i)) = 0$ である．

\medskip\noindent
本書の定義における仮定 $H^0(X, \mathcal{O}_X) = k$ には二つの帰結がある．
一方では，基礎体について曖昧さがないことを意味する．他方では，概形 $X$ が
Cohen--Macaulay かつ次元 $1$ の等次元であることを含意する
（補題 \ref{curves-lemma-automatic}）．$\omega_X$ を補題
\ref{curves-lemma-duality-dim-1} および \ref{curves-lemma-duality-dim-1-CM} の
双対化加群とすると，双対性により
\begin{equation}
\label{curves-equation-genus}
g = \dim_k H^1(X, \mathcal{O}_X) = \dim_k H^0(X, \omega_X)
\end{equation}
を得る．備考 \ref{curves-remark-rework-duality-locally-free} を参照されたい．

\medskip\noindent
$X$ が $k$ 上固有で次元 $\leq 1$ であり，
$H^0(X, \mathcal{O}_X)$ が基礎体 $k$ と等しくない場合には，上の表示公式で
与えた算術種数 $p_a(X)$ の代わりに不変量 $\chi(X, \mathcal{O}_X)$ を用いる．
実際，\cite[276ページ]{FAC} および \cite[序論]{Hirzebruch} では
$\chi(X, \mathcal{O}_X)$ を算術種数と呼ぶべきだと提唱されている．

\begin{lemma}
\label{curves-lemma-genus-base-change}
$k'/k$ を体拡大とする．$X$ を $k$ 上の次元 $1$ の固有概形で，
$H^0(X, \mathcal{O}_X) = k$ を満たすものとする．このとき $X_{k'}$ は
$k'$ 上の次元 $1$ の固有概形であり，
$H^0(X_{k'}, \mathcal{O}_{X_{k'}}) = k'$ を満たす．さらに，
$X_{k'}$ の種数は $X$ の種数に等しい．
\end{lemma}

\begin{proof}
$X_{k'}$ の次元は，例えば Morphisms，補題
\ref{morphisms-lemma-dimension-fibre-after-base-change} により $1$ である．
射 $X_{k'} \to \Spec(k')$ は Morphisms，補題
\ref{morphisms-lemma-base-change-proper} により固有である．
等式 $H^0(X_{k'}, \mathcal{O}_{X_{k'}}) = k'$ は Cohomology of Schemes，補題
\ref{coherent-lemma-flat-base-change-cohomology} から従う．
種数の等式も同じ補題から従う．
\end{proof}

\begin{lemma}
\label{curves-lemma-genus-gorenstein}
$k$ を体とする．$X$ を $k$ 上の次元 $1$ の固有概形で，
$H^0(X, \mathcal{O}_X) = k$ を満たすものとする．$X$ が Gorenstein ならば，
$$
\deg(\omega_X) = 2g - 2
$$
である．ここで $g$ は $X$ の種数で，$\omega_X$ は補題
\ref{curves-lemma-duality-dim-1} のものである．
\end{lemma}

\begin{proof}
補題 \ref{curves-lemma-rr} から直ちに従う．
\end{proof}

\begin{lemma}
\label{curves-lemma-genus-smooth}
$X$ を体 $k$ 上の滑らかな固有曲線で，
$H^0(X, \mathcal{O}_X) = k$ を満たすものとする．このとき
$$
\dim_k H^0(X, \Omega_{X/k}) = g
\quad\text{かつ}\quad
\deg(\Omega_{X/k}) = 2g - 2
$$
である．ここで $g$ は $X$ の種数である．
\end{lemma}

\begin{proof}
補題 \ref{curves-lemma-duality-dim-1} により $\Omega_{X/k} = \omega_X$ である．
従って公式は (\ref{curves-equation-genus}) と補題
\ref{curves-lemma-genus-gorenstein} から従う．
\end{proof}






\section{平面曲線}
\label{curves-section-plane-curves}

\noindent
$k$ を体とする．{\it 平面曲線}とは，曲線 $X$ であって
$\mathbf{P}^2_k$ の閉部分概形と同型なものをいう．埋め込み
$X \to \mathbf{P}^2_k$ は与えられているものとみなすことが多い．
Divisors，例 \ref{divisors-example-closed-subscheme-of-proj} により，
曲線は対応する斉次イデアル
$$
I(X) =
\Ker\left(
k[T_0, T_2, T_2] \longrightarrow \bigoplus \Gamma(X, \mathcal{O}_X(n))
\right)
$$
によって定まる．この状況では，
$$
X = \text{Proj}(k[T_0, T_2, T_2]/I)
$$
が $\mathbf{P}^2_k$ の閉部分概形として成り立つことを思い出そう．
これらの構成に関するより一般的な情報については，
Divisors，例 \ref{divisors-example-closed-subscheme-of-proj} と
そこに挙げられた参考文献を参照されたい．ある斉次多項式に対して
$I(X) = (F)$ となることが分かる．
$F \in k[T_0, T_1, T_2]$ である（補題 \ref{curves-lemma-equation-plane-curve}）．
$X$ は既約なので $F$ も既約である
（補題 \ref{curves-lemma-plane-curve}）．さらに，短完全列
$$
0 \to \mathcal{O}_{\mathbf{P}^2_k}(-d) \xrightarrow{F}
\mathcal{O}_{\mathbf{P}^2_k} \to \mathcal{O}_X \to 0
$$
を調べると，$d = \deg(F)$ のとき $H^0(X, \mathcal{O}_X) = k$ であり，
$X$ の種数が $(d - 1)(d - 2)/2$ であることが分かる．
補題 \ref{curves-lemma-genus-plane-curve} の証明を参照されたい．

\medskip\noindent
滑らかな平面曲線を見つける最も簡単な方法は，方程式を具体的に書くことである．
$p$ を $k$ の標数とする．$p$ が $d$ を割らないならば，
$$
F = T_0^d + T_1^d + T_2^d
$$
と取れる．対応する曲線 $X = V_+(F)$ を次数 $d$ の
{\it Fermat 曲線}という．これは滑らかである．実際，各標準アフィン開
$D_+(T_i)$ 上で，アフィン曲線
$$
\Spec(k[x, y]/(x^d + y^d + 1))
$$
と同型な曲線を得る．環写像
$k \to k[x, y]/(x^d + y^d + 1)$ は Algebra，補題
\ref{algebra-lemma-relative-global-complete-intersection-smooth} により
滑らかである．というのも，$d x^{d - 1}$ と $d y^{d - 1}$ が
$k[x, y]/(x^d + y^d + 1)$ の単位イデアルを生成するからである．
$p | d$ だが $p \not = 3$ ならば，方程式
$$
F = T_0^{d - 1}T_1 + T_1^{d - 1}T_2 + T_2^{d - 1}T_0
$$
を用いればよい．実際，アフィン開上では
$x + x^{d - 1}y + y^{d - 1}$ が得られ，その偏導関数は
$1 - x^{d - 2}y$ と $x^{d - 1} - y^{d - 2}$ である．
三者すべての共通零点集合は空である\footnote{実際，
$x^{d - 1} = y^{d - 2}$ ならば
$0 = x + x^{d - 1}y + y^{d - 1} = x + 2 x^{d - 1} y$ である．
$x \not = 0$ である．実際 $1 = x^{d - 2}y$ である．従って
$0 = 1 + 2x^{d - 2}y = 3$ を得る．これは $3 = 0$ でなければ矛盾である．}．
標数 $3$ の例を作ることは読者に委ねる．

\medskip\noindent
より一般に，任意の体 $k$ と任意の $n$ および $d$ に対して，
次数 $d$ の滑らかな超曲面が $\mathbf{P}^n_k$ 内に存在する．
例えば \cite{Poonen} を参照されたい．

\medskip\noindent
もちろん，この方法で得られる滑らかな曲線の種数は三角数に限られる．
任意の種数の滑らかな曲線を得るには，
$\mathbf{P}^1 \times \mathbf{P}^1$ 上にある滑らかな曲線を探せばよい
（将来の参照をここに挿入する）．

\begin{lemma}
\label{curves-lemma-equation-plane-curve}
$Z \subset \mathbf{P}^2_k$ を，次元 $1$ の等次元で埋め込まれた点をもたない
閉部分概形とする（同値な言い方では，$Z$ は Cohen--Macaulay である）．
このとき，イデアル $I(Z) \subset k[T_0, T_1, T_2]$（$Z$ に対応するもの）は
単項である．
\end{lemma}

\begin{proof}
これは Divisors，補題
\ref{divisors-lemma-equation-codim-1-in-projective-space} の特別な場合である
（Varieties，補題
\ref{varieties-lemma-equation-codim-1-in-projective-space} も参照）．
括弧内の主張は，$1$ 次元 Noether 概形が Cohen--Macaulay であることと
埋め込まれた点をもたないことが同値であることから従う．Divisors，補題
\ref{divisors-lemma-noetherian-dim-1-CM-no-embedded-points} を参照されたい．
\end{proof}

\begin{lemma}
\label{curves-lemma-plane-curve}
$Z \subset \mathbf{P}^2_k$ を補題 \ref{curves-lemma-equation-plane-curve} のものとし，
$I(Z) = (F)$ とする．ここで $F \in k[T_0, T_1, T_2]$ である．
このとき，$Z$ が曲線であるための必要十分条件は $F$ が既約なことである．
\end{lemma}

\begin{proof}
$F$ が可約で，$F = F' F''$ と書けるとする．
$Z'$ を $\mathbf{P}^2_k$ の閉部分概形で，$F'$ により定義されるものとする．
明らかに $Z' \subset Z$ かつ $Z' \not = Z$ である．
$Z'$ も次元 $1$ なので，$Z$ が簡約でないか，または
$Z$ が既約でないかのいずれかである．逆に，
$Z = \sum a_i D_i$ と書く．ここで $D_i$ は $Z$ の既約成分である
（Divisors，補題 \ref{divisors-lemma-codim-1-part} および
\ref{divisors-lemma-codimension-1-is-effective-Cartier}）．
$F_i \in k[T_0, T_1, T_2]$ を $D_i$ のイデアルを生成する斉次多項式とする．
このとき，$F$ と $\prod F_i^{a_i}$ が
$\mathbf{P}^2_k$ の同じ閉部分概形を切り出すことは明らかである．
従って $F = \lambda \prod F_i^{a_i}$ である．ここで
$\lambda \in k^*$ である．これは両者が $Z$ のイデアルを生成するためである．
ゆえに，$F$ が既約ならば $Z$ は素因子，すなわち曲線である．
\end{proof}

\begin{lemma}
\label{curves-lemma-genus-plane-curve}
$Z \subset \mathbf{P}^2_k$ を補題 \ref{curves-lemma-equation-plane-curve} のものとし，
$I(Z) = (F)$ とする．ここで $F \in k[T_0, T_1, T_2]$ である．
このとき $H^0(Z, \mathcal{O}_Z) = k$ であり，$Z$ の種数は
$(d - 1)(d - 2)/2$ である．ここで $d = \deg(F)$ である．
\end{lemma}

\begin{proof}
$S = k[T_0, T_1, T_2]$ とおく．次数付き加群の完全列
$$
0 \to S(-d) \xrightarrow{F} S \to S/(F) \to 0
$$
がある．与えられた閉埋め込みを $i : Z \to \mathbf{P}^2_k$ と書く．
完全関手 $\widetilde{\ }$ を適用すると
（Constructions，補題 \ref{constructions-lemma-proj-sheaves}），
$$
0 \to \mathcal{O}_{\mathbf{P}^2_k}(-d) \to
\mathcal{O}_{\mathbf{P}^2_k} \to i_*\mathcal{O}_Z \to 0
$$
を得る．これは $F$ が $Z$ のイデアルを生成するためである．
$\mathcal{O}_{\mathbf{P}^2_k}(-d)$ と
$\mathcal{O}_{\mathbf{P}^2_k}$ のコホモロジー群は Cohomology of Schemes，補題
\ref{coherent-lemma-cohomology-projective-space-over-ring} に与えられている．
一方，Cohomology of Schemes，補題
\ref{coherent-lemma-relative-affine-cohomology} により
$H^q(Z, \mathcal{O}_Z) = H^q(\mathbf{P}^2_k, i_*\mathcal{O}_Z)$ である．
長完全コホモロジー列を適用すると，まず
$$
k = H^0(\mathbf{P}^2_k, \mathcal{O}_{\mathbf{P}^2_k}) \longrightarrow
H^0(Z, \mathcal{O}_Z)
$$
が同型であり，次に境界写像
$$
H^1(Z, \mathcal{O}_Z) \longrightarrow
H^2(\mathbf{P}^2_k, \mathcal{O}_{\mathbf{P}^2_k}(-d)) \cong
k[T_0, T_1, T_2]_{d - 3}
$$
が同型であることを得る．後者の次元が
$(d - 1)(d - 2)/2$ であることは容易に分かるので，証明は完了する．
\end{proof}

\begin{lemma}
\label{curves-lemma-smooth-plane-curve-point-over-separable}
$Z \subset \mathbf{P}^2_k$ を補題 \ref{curves-lemma-equation-plane-curve} のものとし，
$I(Z) = (F)$ とする．ここで $F \in k[T_0, T_1, T_2]$ である．
$Z \to \Spec(k)$ が少なくとも一点で滑らかで，$k$ が無限体ならば，
滑らかな軌跡に含まれる閉点 $z \in Z$ であって，$\kappa(z)/k$ が
次数高々 $d$ の有限分離拡大となるものが存在する．
\end{lemma}

\begin{proof}
$z' \in Z$ が $Z \to \Spec(k)$ の滑らかな点であるとする．
必要なら座標の番号を付け替えて，$z'$ が $D_+(T_0)$ に含まれると仮定できる．
$f = F(1, x, y) \in k[x, y]$ とおく．すると
$Z \cap D_+(X_0)$ は $k[x, y]/(f)$ のスペクトルと同型である．
$f_x, f_y$ を $f$ の $x, y$ に関する偏導関数とする．
$z'$ は $Z/k$ の滑らかな点なので，$f_x$ または $f_y$ のいずれかは
$z'$ で非零である（Algebra，節 \ref{algebra-section-smooth} の議論を参照）．
座標の番号を付け替えて，$f_y$ は $z'$ で非零であると仮定してよい．
このとき，非空な開部分概形 $V \subset Z \cap D_{+}(X_0)$ であって，射影
$$
p : V \longrightarrow \Spec(k[x])
$$
がエタールとなるものがある．$f$ の $y$ に関する多項式としての次数は
高々 $d$ なので，射影 $p$ のファイバーの次数は高々 $d$ である
（Morphisms，節 \ref{morphisms-section-universally-bounded} の議論を参照）．
さらに $p$ はエタールなので，$p$ の像は開部分
$U \subset \Spec(k[x])$ である．最後に，$k$ は無限体なので，
$k$-有理点の集合 $U(k)$（$U$ の有理点）は無限，従って特に非空である．
任意の $t \in U(k)$ を選び，点 $z \in V$ を $t$ に写るように取る．
この $z$ が求める点である．
\end{proof}





\section{種数 0 の曲線}
\label{curves-section-genus-zero}

\noindent
後で，固有な種数 0 の曲線がどのような形をしているかを知る必要が生じる．
Gorenstein な固有種数 0 曲線は次数 $2$ の平面曲線，すなわち二次曲線である
ことが分かる（補題 \ref{curves-lemma-genus-zero}）．一般の固有種数 0 曲線は，
（より大きな体上の）非特異曲線から押し出しの手続きによって得られる
（補題 \ref{curves-lemma-genus-zero-singular}）．非特異曲線は Gorenstein
なので，これら二つの結果ですべての場合が尽くされる．

\begin{lemma}
\label{curves-lemma-genus-zero-pic}
$X$ を体 $k$ 上の固有曲線で，$H^0(X, \mathcal{O}_X) = k$ を
満たすものとする．$X$ の種数が $0$ ならば，任意の可逆
$\mathcal{O}_X$-加群 $\mathcal{L}$ で次数 $0$ のものは自明である．
\end{lemma}

\begin{proof}
実際，Riemann--Roch（補題 \ref{curves-lemma-rr}）により
$\dim_k H^0(X, \mathcal{L}) \geq 0 + 1 - 0 = 1$ である．従って
$\mathcal{L}$ は非零切断をもち，Varieties，補題
\ref{varieties-lemma-check-invertible-sheaf-trivial} により
$\mathcal{L} \cong \mathcal{O}_X$ である．
\end{proof}

\begin{lemma}
\label{curves-lemma-genus-zero-positive-degree}
$X$ を体 $k$ 上の固有曲線で，$H^0(X, \mathcal{O}_X) = k$ を
満たすものとする．$X$ の種数は $0$ であると仮定する．
$\mathcal{L}$ を可逆 $\mathcal{O}_X$-加群で次数 $d > 0$ のものとする．
このとき，次が成り立つ：
\begin{enumerate}
\item $\dim_k H^0(X, \mathcal{L}) = d + 1$ かつ
$\dim_k H^1(X, \mathcal{L}) = 0$ である；
\item $\mathcal{L}$ は非常に豊富であり，$\mathbf{P}^d_k$ への
閉埋め込みを定める．
\end{enumerate}
\end{lemma}

\begin{proof}
次数と種数の定義により，
$$
\dim_k H^0(X, \mathcal{L}) - \dim_k H^1(X, \mathcal{L}) = d + 1
$$
である．$s$ を $\mathcal{L}$ の非零切断とする．このとき $s$ の零点概形は
有効 Cartier 因子 $D \subset X$ であり，
$\mathcal{L} = \mathcal{O}_X(D)$ で，短完全列
$$
0 \to \mathcal{O}_X \to \mathcal{L} \to \mathcal{L}|_D \to 0
$$
がある．Divisors，補題 \ref{divisors-lemma-characterize-OD} および
備考 \ref{divisors-remark-ses-regular-section} を参照されたい．仮定により
$H^1(X, \mathcal{O}_X) = 0$ なので，写像
$H^0(X, \mathcal{L}) \to H^0(X, \mathcal{L}|_D)$ は全射である．
$\mathcal{L}|_D$ は大域切断で生成される
（$\dim(D) = 0$ であるからで，Varieties，補題
\ref{varieties-lemma-chi-tensor-finite} を参照）ので，可逆加群
$\mathcal{L}$ は大域切断で生成されると結論できる．実際，$D$ は
Artin 概形なので，
$\mathcal{L}|_D \cong \mathcal{O}_D$\footnote{この場合には
$D \to \Spec(k)$ が有限であるため，これは Divisors，補題
\ref{divisors-lemma-finite-trivialize-invertible-upstairs}
から従う．} である．従って，切断 $t$ を $\mathcal{L}$ から取り，
その $D$ への制限が $\mathcal{L}|_D$ を生成するようにできる．
この短完全列から $H^1(X, \mathcal{L}) = 0$ も分かる．

\medskip\noindent
$n \geq 1$ に対して乗法写像
$$
\mu_n :
H^0(X, \mathcal{L}) \otimes_k H^0(X, \mathcal{L}^{\otimes n})
\longrightarrow
H^0(X, \mathcal{L}^{\otimes n + 1})
$$
を考える．これは全射であると主張する．これを見るために短完全列
$$
0 \to \mathcal{L}^{\otimes n} \xrightarrow{s}
\mathcal{L}^{\otimes n + 1} \to \mathcal{L}^{\otimes n + 1}|_D \to 0
$$
を考える．この列の左から来る $\mathcal{L}^{\otimes n + 1}$ の切断は
$\mu_n$ の像に入る．他方，
$H^0(\mathcal{L}) \to H^0(\mathcal{L}|_D)$ は全射であり，
$t^n$ は $\mathcal{L}^{\otimes n}|_D$ の自明化へ写る．従って
$\mu_n(H^0(X, \mathcal{L}) \otimes t^n)$ は
$H^0(X, \mathcal{L}^{\otimes n + 1})$ の部分空間であって，
$\mathcal{L}^{\otimes n + 1}|_D$ の大域切断へ全射するものを与える．
これで主張が証明された．

\medskip\noindent
Varieties，補題 \ref{varieties-lemma-ample-curve} により
$\mathcal{L}$ は豊富であることに注意する．従って Morphisms，補題
\ref{morphisms-lemma-proper-ample-is-proj} は同型
$$
X \longrightarrow
\text{Proj}\left(
\bigoplus\nolimits_{n \geq 0} H^0(X, \mathcal{L}^{\otimes n})\right)
$$
を与える．写像 $\mu_n$ はすべての $n \geq 1$ に対して全射なので，
右辺の次数付き代数は $H^0(X, \mathcal{L})$ 上の対称代数の商である．
$k$-基底 $s_0, \ldots, s_d$ を $H^0(X, \mathcal{L})$ に選ぶと，
これは $d + 1$ 変数多項式代数の商であることが分かる．次数付き環の商は
$\text{Proj}$ の閉埋め込みに対応する
（Constructions，補題
\ref{constructions-lemma-surjective-graded-rings-generated-degree-1-map-proj}）
ので，閉埋め込み $X \to \mathbf{P}^d_k$ を得る．この射が
切断 $s_0, \ldots, s_d$（$\mathcal{L}$ のもの）に付随する Constructions，補題
\ref{constructions-lemma-projective-space} の射であることの確認は省略する．
\end{proof}

\begin{lemma}
\label{curves-lemma-genus-zero}
$X$ を体 $k$ 上の固有曲線で，$H^0(X, \mathcal{O}_X) = k$ を
満たすものとする．$X$ が Gorenstein で種数 $0$ ならば，
$X$ は次数 $2$ の平面曲線と同型である．
\end{lemma}

\begin{proof}
$\mathcal{L} = \omega_X^{\otimes -1}$ という可逆層を考える．ここで
$\omega_X$ は補題 \ref{curves-lemma-duality-dim-1} のものである．補題
\ref{curves-lemma-genus-gorenstein} により $\deg(\omega_X) = -2$ なので，
$\deg(\mathcal{L}) = 2$ である．補題
\ref{curves-lemma-genus-zero-positive-degree} により，基底 $s_0, s_1, s_2$ を，
$k$-ベクトル空間をなす $\mathcal{L}$ の大域切断全体から選ぶと，閉埋め込み
$$
\varphi_{(\mathcal{L}, (s_0, s_1, s_2))} :
X \longrightarrow \mathbf{P}^2_k
$$
を得る．従って $X$ はある次数 $d$ の平面曲線である．
$F \in k[T_0, T_1, T_2]_d$ をその方程式とする
（補題 \ref{curves-lemma-equation-plane-curve}）．$X$ の種数が $0$ なので，
$d$ は $1$ または $2$ である
（補題 \ref{curves-lemma-genus-plane-curve}）．$F$ は $\varphi(X)$ 上で
零切断に制限され，従って $F(s_0, s_1, s_2)$ は
$\mathcal{L}^{\otimes 2}$ の零切断であることに注意する．
$s_0, s_1, s_2$ は線形独立なので $F$ は一次式ではあり得ず，
すなわち $d = \deg(F) \geq 2$ である．従って $d = 2$ であり，
証明は完了する．
\end{proof}

\begin{proposition}[射影直線の特徴づけ]
\label{curves-proposition-projective-line}
$k$ を体とし，$X$ を $k$ 上の固有曲線とする．次は同値である：
\begin{enumerate}
\item $X \cong \mathbf{P}^1_k$,
\item $X$ は $k$ 上滑らかかつ幾何学的既約であり，$X$ の種数は $0$ で，
$X$ は奇数次数の可逆加群をもつ，
\item $X$ は $k$ 上幾何学的整であり，$X$ の種数は $0$ で，
$X$ は Gorenstein であり，$X$ は奇数次数の可逆層をもつ，
\item $H^0(X, \mathcal{O}_X) = k$ で，$X$ の種数は $0$ であり，
$X$ は Gorenstein で，$X$ は奇数次数の可逆層をもつ，
\item $X$ は $k$ 上幾何学的整であり，$X$ の種数は $0$ で，
$X$ は可逆 $\mathcal{O}_X$-加群で次数 $1$ のものをもつ，
\item $H^0(X, \mathcal{O}_X) = k$ で，$X$ の種数は $0$ であり，
$X$ は可逆 $\mathcal{O}_X$-加群で次数 $1$ のものをもつ，
\item $H^1(X, \mathcal{O}_X) = 0$ で，$X$ は可逆
$\mathcal{O}_X$-加群で次数 $1$ のものをもつ，
\item $H^1(X, \mathcal{O}_X) = 0$ であり，$X$ は閉点
$x_1, \ldots, x_n$ をもち，$\mathcal{O}_{X, x_i}$ は正規で，
$\gcd([\kappa(x_i) : k]) = 1$ である，および
\item ここにさらに追加する．
\end{enumerate}
\end{proposition}

\begin{proof}
各条件 (2) -- (8) が (1) を含意することを証明し，
(1) が (2) -- (8) を含意することの確認は省略する．

\medskip\noindent
(2) を仮定する．$k$ 上滑らかな概形は幾何学的簡約であり
（Varieties，補題 \ref{varieties-lemma-smooth-geometrically-normal}），
正則である（Varieties，補題 \ref{varieties-lemma-smooth-regular}）．
従って $X$ は Gorenstein である（Duality for Schemes，補題
\ref{duality-lemma-regular-gorenstein}）．これで (3) に帰着する．

\medskip\noindent
(3) を仮定する．$X$ は $k$ 上幾何学的整なので，Varieties，補題
\ref{varieties-lemma-regular-functions-proper-variety} により
$H^0(X, \mathcal{O}_X) = k$ である．従って (4) に帰着する．

\medskip\noindent
(4) を仮定する．$X$ は Gorenstein なので，補題
\ref{curves-lemma-duality-dim-1} の双対化加群 $\omega_X$ は，補題
\ref{curves-lemma-genus-gorenstein} により次数
$\deg(\omega_X) = -2$ をもつ．奇数次数の可逆加群が存在するという
仮定と合わせると，次数 $1$ の可逆加群が存在すると結論できる．
こうして (6) に帰着する．

\medskip\noindent
(5) を仮定する．$X$ は $k$ 上幾何学的整なので，Varieties，補題
\ref{varieties-lemma-regular-functions-proper-variety} により
$H^0(X, \mathcal{O}_X) = k$ である．従って (6) に帰着する．

\medskip\noindent
(6) を仮定する．このとき補題
\ref{curves-lemma-genus-zero-positive-degree} により
$X \cong \mathbf{P}^1_k$ である．

\medskip\noindent
(7) を仮定する．Varieties，補題
\ref{varieties-lemma-regular-functions-proper-variety} により，
$\kappa = H^0(X, \mathcal{O}_X)$ は $k$ 上有限な体であることに注意する．
$d = [\kappa : k] > 1$ ならば，任意の可逆層の次数は $d$ で割り切れ，
次数 $1$ の可逆層は存在し得ない．従って $d = 1$ であり，(6) に帰着する．

\medskip\noindent
(8) を仮定する．Varieties，補題
\ref{varieties-lemma-regular-functions-proper-variety} により，
$\kappa = H^0(X, \mathcal{O}_X)$ は $k$ 上有限な体であることに注意する．
$\kappa \subset \kappa(x_i)$ なので，次数の最大公約数に関する仮定から
$k = \kappa$ が分かる．同じ条件により，整数 $a_i$ で
$1 = \sum a_i[\kappa(x_i) : k]$ を満たすものを取れる．
Varieties，補題 \ref{varieties-lemma-regular-point-on-curve} により
$x_i$ は $X$ 上の有効 Cartier 因子を定めるので，可逆加群
$\mathcal{O}_X(\sum a_i x_i)$ を考えられる．$a_i$ の選び方により
$\mathcal{L}$ の次数は $1$ である．従って補題
\ref{curves-lemma-genus-zero-positive-degree} により
$X \cong \mathbf{P}^1_k$ である．
\end{proof}

\begin{lemma}
\label{curves-lemma-genus-zero-singular}
$X$ を体 $k$ 上の固有曲線で，$H^0(X, \mathcal{O}_X) = k$ を
満たすものとする．$X$ は特異で種数 $0$ であると仮定する．
このとき，図式
$$
\xymatrix{
x' \ar[d] \ar[r] & X' \ar[d]^\nu \ar[r] & \Spec(k') \ar[d] \\
x \ar[r] & X \ar[r] & \Spec(k)
}
$$
であって，次を満たすものが存在する：
\begin{enumerate}
\item $k'/k$ は非自明な有限拡大である，
\item $X' \cong \mathbf{P}^1_{k'}$,
\item $x'$ は $k'$-有理な $X'$ の点である，
\item $x$ は $k$-有理な $X$ の点である，
\item $X' \setminus \{x'\} \to X \setminus \{x\}$ は同型である，
\item $0 \to \mathcal{O}_X \to \nu_*\mathcal{O}_{X'} \to k'/k \to 0$
は短完全列である．ここで
$k'/k = \kappa(x')/\kappa(x)$ は点 $x$ 上の摩天楼層を表す．
\end{enumerate}
\end{lemma}

\begin{proof}
$\nu : X' \to X$ を $X$ の正規化とする．Varieties，節
\ref{varieties-section-normalization} および
\ref{varieties-section-normalization-one-dimensional} を参照されたい．
$X$ は特異なので，$\nu$ は同型でない．このとき
$k' = H^0(X', \mathcal{O}_{X'})$ は $k$ の有限拡大である
（Varieties，補題
\ref{varieties-lemma-regular-functions-proper-variety}）．短完全列
$$
0 \to \mathcal{O}_X \to \nu_*\mathcal{O}_{X'} \to \mathcal{Q} \to 0
$$
と，$\mathcal{Q}$ の台が有限個の閉点からなることにより，次を得る：
\begin{enumerate}
\item $H^1(X', \mathcal{O}_{X'}) = 0$，すなわち $X'$ は
種数 $0$ の $k'$ 上の曲線である，
\item 短完全列
$0 \to k \to k' \to H^0(X, \mathcal{Q}) \to 0$.
がある．
\end{enumerate}
特に $k'/k$ は非自明な拡大である．

\medskip\noindent
次に，しばしば {\it 導手イデアル}と呼ばれるもの
$$
\mathcal{I} = \SheafHom_{\mathcal{O}_X}(\nu_*\mathcal{O}_{X'}, \mathcal{O}_X)
$$
を考える．これは準連接 $\mathcal{O}_X$-加群である．
$\mathcal{I}$ を $\mathcal{O}_X$ のイデアルとみなすために，写像
$\varphi \mapsto \varphi(1)$ を用いる．従って $\mathcal{I}(U)$ は，
$f \in \mathcal{O}_X(U)$ であって
$f \left(\nu_*\mathcal{O}_{X'}(U)\right) \subset \mathcal{O}_X(U)$
を満たすものの集合である．言い換えれば，
条件は $f$ が $\mathcal{Q}$ を零化することである．同値に，
定義完全列
$$
0 \to \mathcal{I} \to \mathcal{O}_X \to
\SheafHom_{\mathcal{O}_X}(\mathcal{Q}, \mathcal{Q})
$$
がある．$U \subset X$ を $\mathcal{Q}$ の台を含むアフィン開とする．
このとき $V = \mathcal{Q}(U) = H^0(X, \mathcal{Q})$ は
$k$-ベクトル空間であり，その次元は $n - 1$ である．
$\mathcal{O}_X(U) \to \Hom_k(V, V)$ の像は可換部分代数なので，
次元 $\leq n - 1$ を $k$ 上にもつ
（これは行列代数の可換部分代数の性質であり，詳細は省略する）．
従って短完全列
$$
0 \to \mathcal{I} \to \mathcal{O}_X \to \mathcal{A} \to 0
$$
であって，
$\text{Supp}(\mathcal{A}) = \text{Supp}(\mathcal{Q})$ かつ
$\dim_k H^0(X, \mathcal{A}) \leq n - 1$ となるものを得る．
他方，表示
$\mathcal{I} = \SheafHom_{\mathcal{O}_X}(\nu_*\mathcal{O}_{X'}, \mathcal{O}_X)$
は $\mathcal{I}$ に $\nu_*\mathcal{O}_{X'}$-加群構造を与え，
包含写像 $\mathcal{I} \to \nu_*\mathcal{O}_{X'}$ は
$\nu_*\mathcal{O}_{X'}$-加群写像になる．従って
$\mathcal{I} = \nu_*\mathcal{I}'$ である．ここで
$\mathcal{I}' \subset \mathcal{O}_{X'}$ はある準連接イデアル層である．
Morphisms，補題
\ref{morphisms-lemma-affine-equivalence-modules} を参照されたい．
$\mathcal{A}'$ を余核として定める：
$$
0 \to \mathcal{I}' \to \mathcal{O}_{X'} \to \mathcal{A}' \to 0
$$
ここまでの完全列を組み合わせると，短完全列
$0 \to \mathcal{A} \to \nu_*\mathcal{A}' \to \mathcal{Q} \to 0$
を得る．上の評価と $\dim_k H^0(X, \mathcal{Q}) = n - 1$ を合わせると，
$$
\dim_k H^0(X', \mathcal{A}') =
\dim_k H^0(X, \mathcal{A}) + \dim_k H^0(X, \mathcal{Q}) \leq 2 n - 2
$$
を得る．しかし $X'$ は $k'$ 上の曲線なので，左辺は $n$ で割り切れる
（Varieties，補題 \ref{varieties-lemma-divisible}）．
$\mathcal{A}$ と $\mathcal{A}'$ は零ではあり得ないので，
$\dim_k H^0(X', \mathcal{A}') = n$ と結論できる．これは
$\mathcal{I}'$ が $k'$-有理点 $x'$ のイデアル層であることを意味する．
命題 \ref{curves-proposition-projective-line} により
$X' \cong \mathbf{P}^1_{k'}$ を得る．上の等式に戻ると，
$\dim_k H^0(X, \mathcal{A}) = 1$ と結論できる．従って
$\mathcal{I}$ は $k$-有理点 $x$ のイデアル層である．このとき
摩天楼層として $\mathcal{A} = \kappa(x) = k$ かつ
$\mathcal{A}' = \kappa(x') = k'$ である．上で与えた完全列を比較すれば，
補題で述べた構造層に関する結果が直ちに従う．
\end{proof}

\begin{example}
\label{curves-example-squish-on-P1}
実際，補題 \ref{curves-lemma-genus-zero-singular} で述べた状況は，任意の
非自明な有限拡大 $k'/k$ に対して生じる．すなわち，
$$
A = \{f \in k'[x] \mid f(0) \in k \}
$$
を考えられる．$A$ のスペクトルはアフィン曲線であり，
$B = k'[y]$ のスペクトルと，同型 $A_x \cong B_y$ で
$x^{-1}$ を $y$ に送るものを用いて貼り合わせられる．得られるものは
固有曲線 $X$ であり，$H^0(X, \mathcal{O}_X) = k$ を満たす．
その特異点 $x$ は極大イデアル $A \cap (x)$ に対応する．
$X$ の正規化は，補題の場合とちょうど同じく $\mathbf{P}^1_{k'}$ である．
\end{example}








\section{幾何種数}
\label{curves-section-geometric-genus}

\noindent
$X$ が $k$ 上の固有かつ {\bf 滑らかな}曲線で，
$H^0(X, \mathcal{O}_X) = k$ を満たすとき，
$$
p_g(X) = \dim_k H^0(X, \Omega_{X/k})
$$
を $X$ の {\it 幾何種数}という．補題 \ref{curves-lemma-genus-smooth} により，
$X$ の幾何種数は（算術）種数と一致する．しかし高次元では，幾何種数と
算術種数の間に相違がある．備考
\ref{curves-remark-genus-higher-dimension} を参照されたい．

\medskip\noindent
特異曲線に対しては，幾何種数を次のように定義する．

\begin{definition}
\label{curves-definition-geometric-genus}
$k$ を体とし，$X$ を $k$ 上の幾何学的既約曲線とする．$X$ の
{\it 幾何種数}とは，補題 \ref{curves-lemma-smooth-models} におけるような，
$X$ の滑らかな射影的モデルで，$k$ の拡大体上で定義されているかもしれない
ものの種数である．
\end{definition}

\noindent
$k$ が完全ならば，非特異射影的モデル $Y$（$X$ のもの）は滑らかであり
（補題 \ref{curves-lemma-nonsingular-model-smooth}），$X$ の幾何種数は
$Y$ の種数にほかならない．しかし $k$ が完全でないとき，これは
成り立たないことがある．この場合，拡大 $K/k$ を，非特異射影的モデル
$Y_K$（$(X_K)_{red}$ のもの）が滑らかな射影曲線となるように選び，
$X$ の幾何種数を $Y_K$ の種数と定める．これは補題
\ref{curves-lemma-smooth-models} および \ref{curves-lemma-genus-base-change}
により良定義である．

\begin{remark}
\label{curves-remark-genus-higher-dimension}
$X$ を $d$ 次元固有滑らか多様体で，代数閉体 $k$ 上のものとする．このとき
{\it 算術種数}はしばしば
$p_a(X) = (-1)^d(\chi(X, \mathcal{O}_X) - 1)$ と定義され，
{\it 幾何種数}は $p_g(X) = \dim_k H^0(X, \Omega^d_{X/k})$ と定義される．
この状況では，$\omega_X \cong \Omega_{X/k}^d$ はなお成り立つにもかかわらず，
算術種数と幾何種数はもはや一致しない．例えば $d = 2$ ならば，
\begin{align*}
p_a(X) - p_g(X) & =
h^0(X, \mathcal{O}_X) - h^1(X, \mathcal{O}_X) + h^2(X, \mathcal{O}_X) - 1
- h^0(X, \Omega^2_{X/k}) \\
& =
- h^1(X, \mathcal{O}_X) + h^2(X, \mathcal{O}_X) - h^0(X, \omega_X) \\
& =
- h^1(X, \mathcal{O}_X)
\end{align*}
である．ここで $h^i(X, \mathcal{F}) = \dim_k H^i(X, \mathcal{F})$ とし，
最後の等式は双対性から従う．従って曲面に対して，差
$p_g(X) - p_a(X)$ は常に非負である；これは曲面の不正則度と
呼ばれることがある．$X = C_1 \times C_2$ が，種数 $g_1$ と
$g_2$ の滑らかな射影曲線の積ならば，不正則度は $g_1 + g_2$ である．
\end{remark}






\section{Riemann--Hurwitz}
\label{curves-section-riemann-hurewitz}

\noindent
$k$ を体とし，$f : X \to Y$ を $k$ 上の滑らかな曲線の射とする．
このとき標準的完全列
$$
f^*\Omega_{Y/k} \xrightarrow{\text{d}f} \Omega_{X/k}
\longrightarrow \Omega_{X/Y} \longrightarrow 0
$$
を得る．これは Morphisms，補題
\ref{morphisms-lemma-triangle-differentials} による．$X$ と $Y$ は
滑らかなので，層 $\Omega_{X/k}$ と $\Omega_{Y/k}$ は可逆加群である．
Morphisms，補題
\ref{morphisms-lemma-smooth-omega-finite-locally-free} を参照されたい．
最初の写像は非零である，すなわち $f$ は生成点で \'etale であると仮定する．
補題 \ref{curves-lemma-generically-etale} を参照されたい．$R \subset X$ を，
different $\mathfrak{D}_f$（$f$ のもの）が切り出す閉部分概形とする．
Discriminants，補題
\ref{discriminant-lemma-discriminant-quasi-finite-morphism-smooth}
により，これは $\text{d}f$ の零点軌跡と同じで，有効 Cartier 因子であり，
$$
f^*\Omega_{Y/k} \otimes_{\mathcal{O}_X} \mathcal{O}_X(R) = \Omega_{X/k}
$$
を得る．特に，$X$, $Y$ が射影的で，
$k = H^0(Y, \mathcal{O}_Y) = H^0(X, \mathcal{O}_X)$
を満たし，$X$, $Y$ の種数がそれぞれ $g_X$, $g_Y$ ならば，
Riemann--Hurwitz 公式
\begin{align*}
2g_X - 2 & =
\deg(\Omega_{X/k}) \\
& =
\deg(f^*\Omega_{Y/k} \otimes_{\mathcal{O}_X} \mathcal{O}_X(R)) \\
& =
\deg(f) \deg(\Omega_{Y/k}) + \deg(R) \\
& =
\deg(f) (2g_Y - 2) + \deg(R)
\end{align*}
を得る．最初と最後の等式は補題 \ref{curves-lemma-genus-smooth} による．
第二の等式は上で与えた可逆層の同型による．第三の等式は次数の加法性
（Varieties，補題 \ref{varieties-lemma-degree-tensor-product}），
引戻しの次数に関する公式
（Varieties，補題
\ref{varieties-lemma-degree-pullback-map-proper-curves}），そして最後に
$\mathcal{O}_X(R)$ の次数に関する公式
（Varieties，補題
\ref{varieties-lemma-degree-effective-Cartier-divisor}）による．

\medskip\noindent
Riemann--Hurwitz 公式を用いるには，
$\deg(R) = \dim_k \Gamma(R, \mathcal{O}_R)$ を計算する必要がある．
$k$ 上の零次元概形の構造により（例えば Varieties，補題
\ref{varieties-lemma-algebraic-scheme-dim-0} を参照），$R$ は
Artin 局所環のスペクトルの有限非交和
$R = \coprod_{x \in R} \Spec(\mathcal{O}_{R, x})$ であり，各
$\mathcal{O}_{R, x}$ は $k$ 上有限次元である．従って
$$
\deg(R) = \sum\nolimits_{x \in R} \dim_k \mathcal{O}_{R, x} =
\sum\nolimits_{x \in R} d_x [\kappa(x) : k]
$$
ここで
$$
d_x = \text{length}_{\mathcal{O}_{R, x}} \mathcal{O}_{R, x} =
\text{length}_{\mathcal{O}_{X, x}} \mathcal{O}_{R, x}
$$
は $x$ の $R$ における重複度である
（Algebra，補題 \ref{algebra-lemma-pushdown-module} を参照）．
$x \in X$ を閉点とし，その像を $y \in Y$ とする．茎を見ると完全列
$$
\Omega_{Y/k, y} \to \Omega_{X/k, x} \to \Omega_{X/Y, x} \to 0
$$
を得る．局所生成元 $\eta_x$ と $\eta_y$ を，（階数 $1$ の自由）加群
$\Omega_{X/k, x}$ と $\Omega_{Y/k, y}$ にそれぞれ選ぶと，
$
\eta_y \mapsto h \eta_x
$
となる非零元 $h \in \mathcal{O}_{X, x}$ がある．完全列により，
$\Omega_{X/Y, x} \cong \mathcal{O}_{X, x}/h\mathcal{O}_{X, x}$ が
$\mathcal{O}_{X, x}$-加群として成り立つ．因子 $R$ は $h$ により
切り出される（上を参照）ので，
$\mathcal{O}_{R, x} = \mathcal{O}_{X, x}/h\mathcal{O}_{X, x}$ である．
従って次の等式を得る：
\begin{align*}
d_x
& =
\text{length}_{\mathcal{O}_{X, x}}(\mathcal{O}_{R, x}) \\
& =
\text{length}_{\mathcal{O}_{X, x}}(\mathcal{O}_{X, x}/h\mathcal{O}_{X, x}) \\
& =
\text{length}_{\mathcal{O}_{X, x}}(\Omega_{X/Y, x}) \\
& =
\text{ord}_{\mathcal{O}_{X, x}}(h) \\
& =
\text{ord}_{\mathcal{O}_{X, x}}(``\eta_y/\eta_x")
\end{align*}
最初の等式は $d_x$ の定義による．第二と第三の等式は上で見た．
第四の等式は $\text{ord}$ の定義である．Algebra，定義
\ref{algebra-definition-ord} を参照されたい．
$\mathcal{O}_{X, x}$ は離散付値環なので，整数
$\text{ord}_{\mathcal{O}_{X, x}}(h)$ は $h$ の付値にほかならない
ことに注意する．第五の等式は記憶のための表記である．

\medskip\noindent
ここで，重複度 $d_x$ を他の不変量によって「計算」できる場合を述べる．
すなわち，$\kappa(x)$ が $k$ 上分離的ならば，
$\eta_x = \text{d}s$ および $\eta_y = \text{d}t$ と選べる．ここで $s$ と
$t$ はそれぞれ $\mathcal{O}_{X, x}$ と $\mathcal{O}_{Y, y}$ の
素元である（補題 \ref{curves-lemma-uniformizer-works}）．このとき
$t \mapsto u s^{e_x}$ となる単元 $u \in \mathcal{O}_{X, x}$ がある．ここで
$e_x$ は拡大 $\mathcal{O}_{Y, y} \subset \mathcal{O}_{X, x}$ の
分岐指数である．従って
$$
\eta_y = \text{d}t = \text{d}(u s^{e_x}) =
e s^{e_x - 1} u \text{d}s + s^{e_x} \text{d}u
$$
と書ける．$\text{d}u = w \text{d}s$ と書く．ここで
$w \in \mathcal{O}_{X, x}$ である．すると
$$
``\eta_y/\eta_x" = e s^{e_x - 1} u + s^{e_x} w = (e_x u + s w)s^{e_x - 1}
$$
である．この零点の位数は $e_x - 1$ であると結論できる．ただし，
$\kappa(x)$ の標数が $p > 0$ で，$p$ が $e_x$ を割る場合は例外であり，
この場合の零点の位数は $> e_x - 1$ である．

\medskip\noindent
以上をすべて合わせると，$k$ の標数が零ならば，
$$
2g_X - 2 = (2g_Y - 2)\deg(f) +
\sum\nolimits_{x \in X} (e_x - 1)[\kappa(x) : k]
$$
となることが分かる．ここで $e_x$ は
$\mathcal{O}_{X, x}$ の $\mathcal{O}_{Y, f(x)}$ 上の分岐指数である．
この精密な公式が成り立つための必要十分条件は，すべての分岐が順分岐，
すなわち剰余体拡大 $\kappa(x)/\kappa(y)$ が分離的で，$e_x$ が
$k$ の標数と互いに素なことである．ただし，上の議論だけではこれを
証明するには不十分である．補題 \ref{curves-lemma-rhe} とその証明を参照されたい．

\begin{lemma}
\label{curves-lemma-generically-etale}
\begin{slogan}
滑らかな曲線の射が分離的であることと，ほとんど至る所で \'etale
であることは同値である
\end{slogan}
$k$ を体とし，$f : X \to Y$ を $k$ 上の滑らかな曲線の射とする．
次は同値である：
\begin{enumerate}
\item $\text{d}f : f^*\Omega_{Y/k} \to \Omega_{X/k}$ は非零である，
\item $\Omega_{X/Y}$ の台は $X$ の真の閉部分集合に含まれる，
\item 非空な開 $U \subset X$ であって
$f|_U : U \to Y$ が不分岐となるものが存在する，
\item 非空な開 $U \subset X$ であって
$f|_U : U \to Y$ が \'etale となるものが存在する，
\item 関数体の拡大 $k(X)/k(Y)$ は有限分離拡大である．
\end{enumerate}
\end{lemma}

\begin{proof}
$X$ と $Y$ は滑らかなので，層 $\Omega_{X/k}$ と $\Omega_{Y/k}$ は
可逆加群である．Morphisms，補題
\ref{morphisms-lemma-smooth-omega-finite-locally-free} を参照されたい．
完全列
$$
f^*\Omega_{Y/k} \longrightarrow \Omega_{X/k}
\longrightarrow \Omega_{X/Y} \longrightarrow 0
$$
（Morphisms，補題 \ref{morphisms-lemma-triangle-differentials}）を用いると，
(1) と (2) は同値であり，さらに
$f^*\Omega_{Y/k} \to \Omega_{X/k}$ が生成点で非零であるという条件とも
同値であることが分かる．(2) と (3) の同値性は Morphisms，補題
\ref{morphisms-lemma-unramified-omega-zero} から従う．(3) と (4) の
同値性は Morphisms，補題
\ref{morphisms-lemma-flat-unramified-etale} と，平坦性が自動的であること
（補題 \ref{curves-lemma-flat}）から従う．(5) と (4) の同値性を見るには，
Algebra，補題 \ref{algebra-lemma-smooth-at-generic-point} を用いる．
細部は省略する．
\end{proof}

\begin{lemma}
\label{curves-lemma-rh}
$f : X \to Y$ を体 $k$ 上の滑らかな固有曲線の射で，補題
\ref{curves-lemma-generically-etale} の同値な条件を満たすものとする．もし
$k = H^0(Y, \mathcal{O}_Y) = H^0(X, \mathcal{O}_X)$
で，$X$ と $Y$ の種数が $g_X$ と $g_Y$ ならば，
$$
2g_X - 2 = (2g_Y - 2) \deg(f) + \deg(R)
$$
である．ここで $R \subset X$ は $f$ の different が切り出す
有効 Cartier 因子である．
\end{lemma}

\begin{proof}
上の議論を参照されたい；そこでは Discriminants，補題
\ref{discriminant-lemma-discriminant-quasi-finite-morphism-smooth},
補題 \ref{curves-lemma-genus-smooth}，および Varieties，補題
\ref{varieties-lemma-degree-tensor-product} と
\ref{varieties-lemma-degree-pullback-map-proper-curves} を用いた．
\end{proof}

\begin{lemma}
\label{curves-lemma-uniformizer-works}
$X \to \Spec(k)$ が相対次元 $1$ で滑らかであるとする．ここで考える
閉点を $x \in X$ とする．$\kappa(x)$ が $k$ 上分離的ならば，任意の素元
$s$（離散付値環 $\mathcal{O}_{X, x}$ のもの）に対して，元 $\text{d}s$ は
$\Omega_{X/k, x}$ を $\mathcal{O}_{X, x}$ 上自由に生成する．
\end{lemma}

\begin{proof}
Algebra，補題 \ref{algebra-lemma-characterize-smooth-over-field} により，
環 $\mathcal{O}_{X, x}$ は離散付値環である．$x$ は閉じているので，
$\kappa(x)$ は $k$ 上有限である．従って $\kappa(x)/k$ が分離的ならば，
任意の素元 $s$ は Algebra，補題
\ref{algebra-lemma-computation-differential} により
$\Omega_{X/k, x} \otimes_{\mathcal{O}_{X, x}} \kappa(x)$ の非零元へ写る．
$\Omega_{X/k, x}$ は階数 $1$ の自由 $\mathcal{O}_{X, x}$-加群なので，
結果が従う．
\end{proof}

\begin{lemma}
\label{curves-lemma-rhe}
記法と仮定は補題 \ref{curves-lemma-rh} と同じとする．閉点
$x \in X$ に対し，$d_x$ を $x$ の $R$ における重複度とする．このとき
$$
2g_X - 2 = (2g_Y - 2) \deg(f) + \sum\nolimits d_x [\kappa(x) : k]
$$
である．さらに次が成り立つ：
\begin{enumerate}
\item $d_x = \text{length}_{\mathcal{O}_{X, x}}(\Omega_{X/Y, x})$,
\item $d_x \geq e_x - 1$ である．ここで $e_x$ は
$\mathcal{O}_{X, x}$ の $\mathcal{O}_{Y, y}$ 上の分岐指数である，
\item $d_x = e_x - 1$ であるための必要十分条件は，
$\mathcal{O}_{X, x}$ が $\mathcal{O}_{Y, y}$ 上順分岐することである．
\end{enumerate}
\end{lemma}

\begin{proof}
補題 \ref{curves-lemma-rh} と上の議論
（そこでは Varieties，補題
\ref{varieties-lemma-algebraic-scheme-dim-0} および Algebra，補題
\ref{algebra-lemma-pushdown-module} を用いた）により，重複度
$d_x$（$x$ の $R$ におけるもの）に関する結果を証明すれば十分である．(1) は上の
議論で証明した．上の議論では，$\kappa(x)$ が $k$ 上分離的な場合に
限って (2) と (3) を証明した．証明の残りでは，準有限 Gorenstein 射の
different に関する材料を用いて (2) と (3) を一様に扱う．

\medskip\noindent
まず，$f$ が準有限 Gorenstein 射であることに注意する．例えば，
$f$ は平坦準有限射で $X$ は Gorenstein であるため
（Duality for Schemes，補題
\ref{duality-lemma-flat-morphism-from-gorenstein-scheme} を参照），
または Discriminants，補題
\ref{discriminant-lemma-discriminant-quasi-finite-morphism-smooth}
（上で用いた）の証明で示されているため，これは成り立つ．従って
Discriminants，補題 \ref{discriminant-lemma-gorenstein-quasi-finite}
により $\omega_{X/Y}$ は可逆であり，$X$ を開部分で置き換え，さらに
ある $Y' \to Y$ で基底変換した後も同じことが成り立つ．以下ではこれを
断らずに用いる．

\medskip\noindent
アフィン開 $U \subset X$ と $V \subset Y$ を，
$x \in U$, $y \in V$, $f(U) \subset V$ であり，$x$ が $U$ の点で
$y$ 上にある唯一のものとなるように選ぶ．$U = \Spec(A)$ および
$V = \Spec(B)$ と書く．このとき $R \cap U$ は
$f|_U : U \to V$ の different である．Discriminants，補題
\ref{discriminant-lemma-base-change-different} により，この場合
different の形成は任意の基底変換と可換である．$U$ と $V$ の選び方により，
$$
A \otimes_B \kappa(y) =
\mathcal{O}_{X, x} \otimes_{\mathcal{O}_{Y, y}} \kappa(y) =
\mathcal{O}_{X, x}/(s^{e_x})
$$
である．ここで $e_x$ は補題の主張における分岐指数である．
$C = \mathcal{O}_{X, x}/(s^{e_x})$ を $\kappa(y)$ 上の有限代数とみなす．
$\mathfrak{D}_{C/\kappa(y)}$ を Discriminants，定義
\ref{discriminant-definition-different} の意味での，$C$ の
$\kappa(y)$ 上の different とする．次を示せば十分である：
$\mathfrak{D}_{C/\kappa(y)}$ が非零であるための必要十分条件は，拡大
$\mathcal{O}_{Y, y} \subset \mathcal{O}_{X, x}$ が順分岐することであり，
順分岐の場合には $\mathfrak{D}_{C/\kappa(y)}$ は
$s^{e_x - 1}$ が生成する $C$ のイデアルに等しい．順分岐とはちょうど，
$\kappa(x)/\kappa(y)$ が分離的で，$\kappa(y)$ の標数が $e_x$ を
割らないことだった．他方，$C/\kappa(y)$ の different が非零である
ための必要十分条件は
$\tau_{C/\kappa(y)} \in \omega_{C/\kappa(y)}$ が非零なことである．
実際，$\omega_{C/\kappa(y)}$ は可逆 $C$-加群
（$\omega_{A/B}$ の基底変換）なので，階数 $1$ の自由加群である．
その生成元を $\lambda$ とする．
$\tau_{C/\kappa(y)} = h\lambda$ と書く．ここで $h \in C$ である．このとき
$\mathfrak{D}_{C/\kappa(y)} = (h) \subset C$ であり，主張が従う．
Discriminants，補題 \ref{discriminant-lemma-tau-nonzero} により，
$\tau_{C/\kappa(y)} \not = 0$ であるための必要十分条件は，
$\kappa(x)/\kappa(y)$ が分離的で，$e_x$ が標数と互いに素なことである．
最後に，$\tau_{C/\kappa(y)}$ が非零であっても，なお
$s \tau_{C/\kappa(y)} = 0$ である．というのも，
$s\tau_{C/\kappa(y)} : C \to \kappa(y)$ は $c$ を冪零作用素 $sc$ の
トレースへ送り，これは零だからである．従って $sh = 0$ であり，
$h \in (s^{e_x - 1})$ である．ゆえに常に
$\mathfrak{D}_{C/\kappa(y)} \subset (s^{e_x - 1})$ である．
$(s^{e_x - 1}) \subset C$ は最小の非零イデアルなので，最後の主張を
証明した．
\end{proof}






\section{非分離射}
\label{curves-section-inseparable}

\noindent
非分離射のもとでの種数の振る舞いについて，いくつか注意を述べる．

\begin{lemma}
\label{curves-lemma-dominated-by-smooth}
$k$ を体とする．$f : X \to Y$ を $k$ 上の曲線の全射射とする．
$X$ が $k$ 上滑らかで $Y$ が正規ならば，$Y$ は $k$ 上滑らかである．
\end{lemma}

\begin{proof}
$y \in Y$ とする．$x \in X$ を，$y$ に写るように選ぶ．
Varieties，補題 \ref{varieties-lemma-flat-under-smooth} により，
$f$ が $x$ で平坦であることを示せば十分である．
これは補題 \ref{curves-lemma-flat} から従う．
\end{proof}

\begin{lemma}
\label{curves-lemma-purely-inseparable}
$k$ を標数 $p > 0$ の体とする．$f : X \to Y$ を $k$ 上の
固有非特異曲線の非定数射とする．関数体の拡大 $k(X)/k(Y)$ が
純非分離ならば，分解
$$
X = X_0 \to X_1 \to \ldots \to X_n = Y
$$
が存在し，各 $X_i$ は固有非特異曲線で，$X_i \to X_{i + 1}$ は
次数 $p$ の射であり，$k(X_{i + 1}) \subset k(X_i)$ は非分離である．
\end{lemma}

\begin{proof}
これは定理 \ref{curves-theorem-curves-rational-maps} と，体の有限純非分離拡大は
常に次数 $p$ の（非分離）拡大の列として得られるという事実から従う．
Fields，補題 \ref{fields-lemma-finite-purely-inseparable} を参照せよ．
\end{proof}

\begin{lemma}
\label{curves-lemma-inseparable-deg-p-smooth}
$k$ を標数 $p > 0$ の体とする．$f : X \to Y$ を $k$ 上の
固有非特異曲線の非定数射とする．$X$ が滑らかで，
$k(Y) \subset k(X)$ が次数 $p$ の非分離拡大ならば，同型
$Y = X^{(p)}$ が一意に存在し，そのもとで $f$ は $F_{X/k}$ である．
\end{lemma}

\begin{proof}
相対 Frobenius 射 $F_{X/k} : X \to X^{(p)}$ は Varieties，節
\ref{varieties-section-frobenius} で構成されている．$X^{(p)}$ は $k$ 上の
滑らかな固有曲線であり，$X$ の基底変換であることに注意する．
Varieties，補題 \ref{varieties-lemma-inseparable-deg-p-smooth} により，
射 $F_{X/k}$ の次数は $p$ である．したがって $k(X^{(p)})$ と $k(Y)$ は
いずれも $k(X)$ の部分体で，
$[k(X) : k(Y)] = [k(X) : k(X^{(p)})] = p$ を満たす．補題を証明するには，
$k(Y) = k(X^{(p)})$ が $k(X)$ の中で成り立つことを示せば十分である．
定理 \ref{curves-theorem-curves-rational-maps} を参照せよ．

\medskip\noindent
$K = k(X)$ と書く．写像 $\text{d} : K \to \Omega_{K/k}$ を考える．
補題 \ref{curves-lemma-generically-etale} から，$k(Y)$ は $\text{d}$ の核に
含まれる．Varieties，補題
\ref{varieties-lemma-relative-frobenius-omega} により，
$k(X^{(p)})$ も $\text{d}$ の核に含まれることが分かる．$X$ は滑らかな
曲線なので，$\Omega_{K/k}$ は次元 $1$ の $K$-ベクトル空間である．
そこで More on Algebra，補題 \ref{more-algebra-lemma-p-basis} により，
$\Ker(\text{d}) = kK^p$ かつ $[K : kK^p] = p$ である．
したがって次数を考えれば $k(Y) = kK^p = k(X^{(p)})$ となる．
\end{proof}

\begin{lemma}
\label{curves-lemma-purely-inseparable-smooth}
$k$ を標数 $p > 0$ の体とする．$f : X \to Y$ を $k$ 上の
固有非特異曲線の非定数射とする．$X$ が滑らかで，
$k(Y) \subset k(X)$ が純非分離ならば，一意な $n \geq 0$ と一意な同型
$Y = X^{(p^n)}$ が存在し，$f$ は $n$ 重の $X/k$ の相対 Frobenius である．
\end{lemma}

\begin{proof}
この $n$ 重の $X/k$ の相対 Frobenius は Varieties，注意
\ref{varieties-remark-n-fold-relative-frobenius} で定義されている．
補題 \ref{curves-lemma-inseparable-deg-p-smooth} と
\ref{curves-lemma-purely-inseparable} を組み合わせれば本補題が従う．
\end{proof}

\begin{lemma}
\label{curves-lemma-purely-inseparable-smooth-genus}
$k$ を標数 $p > 0$ の体とする．$f : X \to Y$ を $k$ 上の
固有非特異曲線の非定数射とする．次を仮定する．
\begin{enumerate}
\item $X$ は滑らかである；
\item $H^0(X, \mathcal{O}_X) = k$,
\item $k(X)/k(Y)$ は純非分離である．
\end{enumerate}
このとき $Y$ は滑らかで，$H^0(Y, \mathcal{O}_Y) = k$ であり，
$Y$ の種数は $X$ の種数に等しい．
\end{lemma}

\begin{proof}
補題 \ref{curves-lemma-purely-inseparable-smooth} により，
$Y = X^{(p^n)}$ は $X$ の $F_{\Spec(k)}^n$ による基底変換である．
したがって $Y$ は滑らかであり，コホモロジーと種数に関する主張は
補題 \ref{curves-lemma-genus-base-change} から従う．
\end{proof}

\begin{example}
\label{curves-example-inseparable}
この例では，非特異射影曲線の純非分離射のもとで種数が変わりうることを示す．
$k$ を標数 $3$ の体とする．元 $a \in k$ であって $3$ 乗でないものが
存在すると仮定する．
例えば $k = \mathbf{F}_3(a)$ とすればよい．$X$ を，節
\ref{curves-section-plane-curves} における斉次方程式
$$
F = T_1^2T_0 - T_2^3 + aT_0^3
$$
で定まる平面曲線とする．アフィン部分 $D_+(T_0)$ 上で座標
$x = T_1/T_0$ と $y = T_2/T_0$ を用いると，非特異アフィン曲線を定める
$x^2 - y^3 + a = 0$ を得る．さらに，無限遠点 $(0 : 1: 0)$ は滑らかな点である．
したがって $X$ は種数 $1$ の非特異射影曲線である
（補題 \ref{curves-lemma-genus-plane-curve}）．一方，射
$f : X \to \mathbf{P}^1_k$ であって，$D_+(T_0)$ 上で $(x, y)$ を
$x \in \mathbf{A}^1_k \subset \mathbf{P}^1_k$ に写すものを考える．このとき $f$ は $k$ 上の
固有非特異曲線の射であり，次数 $p = 3$ の非分離関数体拡大を誘導するが，
$X$ の種数は $1$，$\mathbf{P}^1_k$ の種数は $0$ である．
\end{example}

\begin{proposition}
\label{curves-proposition-unwind-morphism-smooth}
$k$ を標数 $p > 0$ の体とする．$f : X \to Y$ を $k$ 上の
滑らかな固有曲線の非定数射とする．このとき $f$ は
$$
X \longrightarrow X^{(p^n)} \longrightarrow Y
$$
と分解できる．ここで $X^{(p^n)} \to Y$ は滑らかな固有曲線の非定数射で，
分離体拡大 $k(X^{(p^n)})/k(Y)$ を誘導し，
$$
X^{(p^n)} = X \times_{\Spec(k), F_{\Spec(k)}^n} \Spec(k),
$$
であり，$X \to X^{(p^n)}$ は $n$ 重の $X$ の相対 Frobenius である．
\end{proposition}

\begin{proof}
Fields，補題 \ref{fields-lemma-separable-first} により，中間拡大
$k(X)/E/k(Y)$ であって，$k(X)/E$ が純非分離で $E/k(Y)$ が分離となるものが
存在する．定理 \ref{curves-theorem-curves-rational-maps} により，これは
$X \to Z \to Y$ という $f$ の分解に対応し，ここで $Z$ は非特異固有曲線である．
射 $X \to Z$ に補題 \ref{curves-lemma-purely-inseparable-smooth} を適用すればよい．
\end{proof}

\begin{lemma}
\label{curves-lemma-inseparable-linear-system}
$k$ を標数 $p > 0$ の体とする．$X$ を $k$ 上の滑らかな固有曲線とする．
$(\mathcal{L}, V)$ を $\mathfrak g^r_d$ とし，$r \geq 1$ とする．
このとき次の二つのいずれかが成り立つ．
\begin{enumerate}
\item $\mathfrak g^1_d$ が存在し，それに対応する射
$X \to \mathbf{P}^1_k$（補題 \ref{curves-lemma-linear-series}）は生成点で \'etale
（すなわち補題 \ref{curves-lemma-generically-etale} の意味で）である；または
\item $\mathfrak g^r_{d'}$ が $X^{(p)}$ 上に存在し，$d' \leq d/p$ である．
\end{enumerate}
\end{lemma}

\begin{proof}
$k$ 上線形独立な二元 $s, t \in V$ を選ぶ．このとき $f = s/t$ は射
$X \to \mathbf{P}^1_k$ を定める有理関数であり，この射に対応する線形系列は
$(\mathcal{L}, ks + kt)$ である．この射が生成点で \'etale でなければ，命題
\ref{curves-proposition-unwind-morphism-smooth} により $f \in k(X^{(p)})$ である．
次に基底 $s_0, \ldots, s_r$ を $V$ から選び，$\mathcal{L}' \subset \mathcal{L}$ を
$s_0, \ldots, s_r$ が生成する可逆層とする．$f_i = s_i/s_0$ を $k(X)$ の元として
定める．各組 $(s_0, s_i)$ に対して
$f_i \in k(X^{(p)})$ ならば，射
$$
\varphi = \varphi_{(\mathcal{L}', (s_0, \ldots, s_r)} :
X
\longrightarrow
\mathbf{P}^r_k = \text{Proj}(k[T_0, \ldots, T_r])
$$
は $X^{(p)}$ を経由する．実際，これはアフィン開 $D_+(T_0)$ 上で成り立ち，
補題 \ref{curves-lemma-extend-over-normal-curve} により，アフィン部分上の射を
滑らかな曲線 $X^{(p)}$ 全体へ延長できる．記号を導入し，分解
$$
X \xrightarrow{F_{X/k}} X^{(p)} \xrightarrow{\psi} \mathbf{P}^r_k
$$
を $\varphi$ の分解として書く．このとき
$\mathcal{N} = \psi^*\mathcal{O}_{\mathbf{P}^1_k}(1)$ は可逆
$\mathcal{O}_{X^{(p)}}$-加群で，$\mathcal{L}' = F_{X/k}^*\mathcal{N}$ を満たし，
$\psi^*T_0, \ldots, \psi^*T_r$ は $k$ 上線形独立である（それらの引き戻しは
$s_0, \ldots, s_r$ に $X$ 上で一致するからである）．
最後に，
$$
d = \deg(\mathcal{L}) \geq \deg(\mathcal{L}') =
\deg(F_{X/k}) \deg(\mathcal{N}) = p \deg(\mathcal{N})
$$
を得るから，望む結論が従う．ここでは Varieties，補題
\ref{varieties-lemma-check-invertible-sheaf-trivial}，
\ref{varieties-lemma-degree-pullback-map-proper-curves}，および
\ref{varieties-lemma-inseparable-deg-p-smooth} を用いた．
\end{proof}

\begin{lemma}
\label{curves-lemma-point-over-separable-extension}
$k$ を体とする．$X$ を $k$ 上の滑らかな固有曲線で，
$H^0(X, \mathcal{O}_X) = k$ かつ種数 $g \geq 2$ なるものとする．このとき，
閉点 $x \in X$ であって，$\kappa(x)/k$ が次数 $\leq 2g - 2$ の分離拡大と
なるものが存在する．
\end{lemma}

\begin{proof}
$\omega = \Omega_{X/k}$ とおく．補題 \ref{curves-lemma-genus-smooth} により，
これは次数 $2g - 2$ で，$g$ 個の大域切断をもつ．したがって
$\mathfrak g^{g - 1}_{2g - 2}$ が得られる．自明な補題
\ref{curves-lemma-linear-series-trivial-existence} により $\mathfrak g^1_{2g - 2}$ が
存在し，補題 \ref{curves-lemma-g1d} により射
$$
\varphi : X \longrightarrow \mathbf{P}^1_k
$$
を得る．その次数を $d \leq 2g - 2$ とする．$\varphi$ は平坦
（補題 \ref{curves-lemma-flat}）かつ有限（補題 \ref{curves-lemma-finite}）なので，
次数 $d$ の有限局所自由射である
（Morphisms，補題 \ref{morphisms-lemma-finite-flat}）．任意の有理点
$t \in \mathbf{P}^1_k$ と，任意の点 $x \in X$ であって
$\varphi(x) = t$ を満たすものを選ぶ．
すると
$$
d \geq [\kappa(x) : \kappa(t)] = [\kappa(x) : k]
$$
である．例えば Morphisms，補題
\ref{morphisms-lemma-finite-locally-free-universally-bounded} および
\ref{morphisms-lemma-characterize-universally-bounded} による．したがって
$k$ が完全（例えば標数零または有限）ならば補題は証明された．よって次の段落で
扱う場合に帰着する．

\medskip\noindent
$k$ を標数 $p > 0$ の無限体と仮定する．上と同様に，$X$ が
$\mathfrak g^{g - 1}_{2g - 2}$ をもつことを用いる．滑らかな固有曲線
$X^{(p)}$ の種数は $X$ と同じであり，したがって $> 0$ である．補題
\ref{curves-lemma-grd-inequalities} により，$X^{(p)}$ は $\mathfrak g^{g - 1}_d$ を
どの $d \leq g - 1$ に対してももたない．われわれの
$\mathfrak g^{g - 1}_{2g - 2}$ に補題
\ref{curves-lemma-inseparable-linear-system} を適用すると（$2g - 2/p \leq g - 1$
にも注意して），可能性 (2) は起こらないと分かる．したがって射
$$
\varphi : X \longrightarrow \mathbf{P}^1_k
$$
を得る．これは（補題の意味で）生成点で \'etale であり，次数は
$\leq 2g - 2$ である．$U \subset X$ を $\varphi$ が \'etale となる空でない
開部分概形とする．すると $\varphi(U) \subset \mathbf{P}^1_k$ は空でない
Zariski 開集合であり，$k$-有理点 $t \in \varphi(U)$ を選べる．実際，$k$ は
無限である．点 $u \in U$ であって $\varphi(u) = t$ を満たすものをとる．
このとき $\kappa(u)/\kappa(t)$ は分離であり
（Morphisms，補題 \ref{morphisms-lemma-etale-over-field}），
$\kappa(t) = k$ で，前と同様に $[\kappa(u) : k] \leq 2g - 2$ である．
\end{proof}

\noindent
次の補題は実のところ本節に属するものではないが，ほかに適切な置き場所が
見当たらない．

\begin{lemma}
\label{curves-lemma-ramification-to-algebraic-closure}
$X$ を体 $k$ 上の滑らかな曲線とする．
$\overline{x} \in X_{\overline{k}}$ を，像が $x \in X$ である閉点とする．
$\mathcal{O}_{X, x} \subset \mathcal{O}_{X_{\overline{k}}, \overline{x}}$ の
分岐指数は $\kappa(x)/k$ の非分離次数である．
\end{lemma}

\begin{proof}
$X$ を縮小すれば，\'etale 射 $\pi : X \to \mathbf{A}^1_k$ が存在すると
仮定してよい．Morphisms，補題
\ref{morphisms-lemma-smooth-etale-over-affine-space} を参照せよ．
このとき局所環の図式
$$
\xymatrix{
\mathcal{O}_{X_{\overline{k}}, \overline{x}} &
\mathcal{O}_{\mathbf{A}^1_{\overline{k}}, \pi(\overline{x})} \ar[l] \\
\mathcal{O}_{X, x} \ar[u] &
\mathcal{O}_{\mathbf{A}^1_k, \pi(x)} \ar[l] \ar[u]
}
$$
を考えられる．水平射は \'etale 射に対応するため，その分岐指数は $1$ である．
さらに，拡大 $\kappa(x)/\kappa(\pi(x))$ は分離だから，$\kappa(x)$ と
$\kappa(\pi(x))$ は $k$ 上同じ非分離次数をもつ．分岐指数の乗法性により，
$x$ がアフィン直線の点である場合に結果を証明すれば十分である．

\medskip\noindent
$X = \mathbf{A}^1_k$ と仮定する．この場合，$X$ の $x$ における局所環は
$$
\mathcal{O}_{X, x} = k[t]_{(P)}
$$
という形をしている．ここで $P$ は $k$ 上の既約モニック多項式である．
このとき $P(t) = Q(t^q)$ と書け，ここで $Q \in k[t]$ は分離多項式である．
Fields，補題 \ref{fields-lemma-irreducible-polynomials} を参照せよ．
$\kappa(x) = k[t]/(P)$ の非分離次数は $q$ であり，これは $k$ 上の次数である．
一方，$\overline{k}$ 上で $Q(t) = \prod (t - \alpha_i)$ と因数分解でき，
$\alpha_i$ たちは相異なる．$\alpha_i = \beta_i^q$ となる一意な
$\beta_i \in \overline{k}$ をとる．すると点 $\overline{x}$ はいずれかの
$\beta_i$ に対応する．そして
$$
k[t]_{(P)} \longrightarrow \overline{k}[t]_{(t - \beta_i)}
$$
の分岐指数は実際 $q$ に等しい．なぜなら素元 $P$ は
$(t - \beta_i)^q$ と単元との積に写るからである．
\end{proof}





\section{押し出し}
\label{curves-section-pushouts}

\noindent
$k$ を体とする．実線図式
$$
\xymatrix{
Z' \ar[d] \ar[r]_{i'} & X' \ar@{..>}[d]^a \\
Z \ar@{..>}[r]^i & X
}
$$
を考える．これは $k$ 上の概形の図式であって，次を満たす．
\begin{enumerate}
\item[(a)] $X'$ は $k$ 上分離かつ有限型で，次元は $\leq 1$ である；
\item[(b)] $i : Z' \to X'$ は閉埋め込みである；
\item[(c)] $Z'$ と $Z$ は $\Spec(k)$ 上有限である；
\item[(d)] $Z' \to Z$ は全射である．
\end{enumerate}
この状況では，$X'$ の点からなる任意の有限集合はアフィン開集合に含まれる．
Varieties，命題
\ref{varieties-proposition-finite-set-of-points-of-codim-1-in-affine} を参照せよ．
したがって More on Morphisms，命題
\ref{more-morphisms-proposition-pushout-along-closed-immersion-and-integral} の
仮定が満たされ，次を得る．
\begin{enumerate}
\item 押し出し $X = Z \amalg_{Z'} X'$ が概形の圏で存在する；
\item $i : Z \to X$ は閉埋め込みである；
\item $a : X' \to X$ は整かつ全射である；
\item More on Morphisms，補題 \ref{more-morphisms-lemma-pushout-separated} により
$X \to \Spec(k)$ は分離である；
\item More on Morphisms，補題 \ref{more-morphisms-lemma-pushout-finite-type} により
$X \to \Spec(k)$ は有限型である；
\item したがって Morphisms，補題 \ref{morphisms-lemma-finite-integral} および
\ref{morphisms-lemma-permanence-finite-type} により $a : X' \to X$ は有限である；
\item $X' \to \Spec(k)$ が固有ならば，Morphisms，補題
\ref{morphisms-lemma-image-proper-is-proper} により $X \to \Spec(k)$ は固有である．
\end{enumerate}

\noindent
次の補題は大幅に一般化できる．

\begin{lemma}
\label{curves-lemma-complete-local-ring-pushout}
上の状況で，$Z = \Spec(k')$ とし，ここで $k'$ は体であるとする．また
$Z' = \Spec(k'_1 \times \ldots \times k'_n)$ とし，$k'_i/k'$ は体の有限拡大とする．
$x \in X$ を $Z \to X$ の像，$x'_i \in X'$ を
$\Spec(k'_i) \to X'$ の像とする．このときファイバー積図式
$$
\xymatrix{
\prod\nolimits_{i = 1, \ldots, n} k'_i &
\prod\nolimits_{i = 1, \ldots, n} \mathcal{O}_{X', x'_i}^\wedge \ar[l] \\
k' \ar[u] &
\mathcal{O}_{X, x}^\wedge \ar[u] \ar[l]
}
$$
を得る．ここで水平射は剰余体への写像によって与えられる．
\end{lemma}

\begin{proof}
$\Spec(A)$ を $x$ の $X$ におけるアフィン開近傍として選ぶ．逆像を
$\Spec(A') \subset X'$ とする．構成によりファイバー積図式
$$
\xymatrix{
\prod\nolimits_{i = 1, \ldots, n} k'_i &
A' \ar[l] \\
k' \ar[u] &
A \ar[u] \ar[l]
}
$$
を得る．すべてが $A$ 上有限なので，極大イデアル $\mathfrak m \subset A$ であって
$x$ に対応するものに関して完備化した後も，この図式はファイバー積図式の
ままである（Algebra，補題 \ref{algebra-lemma-completion-flat}）．最後に
Algebra，補題 \ref{algebra-lemma-completion-finite-extension} を適用して，
$A'$ の完備化を同定する．
\end{proof}




\section{貼り合わせと押し潰し}
\label{curves-section-glueing-squishing}

\noindent
以下では $k[\epsilon]$ により，Varieties，定義
\ref{varieties-definition-dual-numbers} で定義された $k$ 上の双対数の代数を表す．

\begin{lemma}
\label{curves-lemma-no-in-between-over-k}
$k$ を代数閉体とする．$k \subset A$ を環拡大とし，$A$ がちょうど二つの
$k$-部分代数をもつとする．このとき $A = k \times k$ または
$A = k[\epsilon]$ である．
\end{lemma}

\begin{proof}
仮定は $k \not = A$ であり，任意の部分環 $k \subset C \subset A$ が
$k$ または $A$ に等しいことを意味する．$t \in A$，$t \not \in k$ とする．
すると $A$ は $t$ によって $k$ 上生成される．したがって $A = k[x]/I$ である．
ここで $I$ はあるイデアルである．$I = (0)$ ならば，許されない部分代数
$k[x^2]$ が存在する．そうでなければ $I$ はモニック多項式 $P$ によって
生成される．$P = \prod_{i = 1}^d (t - a_i)$ と書く．$d > 2$ ならば，
$(t - a_1)(t - a_2)$ が生成する部分代数から矛盾を得る．したがって $d = 2$ である．
$a_1 \not = a_2$ ならば $A = k \times k$，$a_1 = a_2$ ならば
$A = k[\epsilon]$ である．
\end{proof}

\begin{example}[点の貼り合わせ]
\label{curves-example-glue-points}
$k$ を代数閉体とする．$f : X' \to X$ を代数的 $k$-概形の射とする．
次が成り立つとき，$X$ は $a$ と $b$ を $X'$ において貼り合わせて得られるという．
\begin{enumerate}
\item $a, b \in X'(k)$ は相異なる点で，同じ点 $x \in X(k)$ に写る；
\item $f$ は有限で，$f^{-1}(X \setminus \{x\}) \to X \setminus \{x\}$ は同型である；
\item 短完全列
$$
0 \to \mathcal{O}_X \to f_*\mathcal{O}_{X'} \xrightarrow{a - b} x_*k \to 0
$$
が存在する．ここで右側の射は局所切断 $h$ を $f_*\mathcal{O}_{X'}$ からとり，
差 $h(a) - h(b) \in k$ に写す．
\end{enumerate}
この場合，短完全列
$$
0 \to \mathcal{O}_X^* \to f_*\mathcal{O}_{X'}^*
\xrightarrow{ab^{-1}} x_*k^* \to 0
$$
も存在する．ここで右側の射は局所切断 $h$ を $f_*\mathcal{O}_{X'}^*$ からとり，
乗法的な差 $h(a)h(b)^{-1} \in k^*$ に写す．
\end{example}

\begin{example}[接ベクトルの押し潰し]
\label{curves-example-squish-tangent-vector}
$k$ を代数閉体とする．$f : X' \to X$ を代数的 $k$-概形の射とする．
次が成り立つとき，$X$ は接ベクトル $\vartheta$ を $X'$ において
押し潰して得られるという．
\begin{enumerate}
\item $\vartheta : \Spec(k[\epsilon]) \to X'$ は $k$ 上の閉埋め込みで，
$f \circ \vartheta$ は点 $x \in X(k)$ を経由する；
\item $f$ は有限で，$f^{-1}(X \setminus \{x\}) \to X \setminus \{x\}$ は同型である；
\item 短完全列
$$
0 \to \mathcal{O}_X \to f_*\mathcal{O}_{X'} \xrightarrow{\vartheta} x_*k \to 0
$$
が存在する．ここで右側の射は局所切断 $h$ を $f_*\mathcal{O}_{X'}$ からとり，
$\epsilon$ の，$\vartheta^\sharp(h) \in k[\epsilon]$ における係数へ写す．
\end{enumerate}
この場合，短完全列
$$
0 \to \mathcal{O}_X^* \to f_*\mathcal{O}_{X'}^*
\xrightarrow{\vartheta} x_*k \to 0
$$
も存在する．ここで右側の射は局所切断 $h$ を $f_*\mathcal{O}_{X'}^*$ からとり，
$\text{d}\log(\vartheta^\sharp(h))$ に写す．ただし
$\text{d}\log : k[\epsilon]^* \to k$ は $a + b\epsilon$ を $b/a \in k$ に写す
アーベル群の準同型である．
\end{example}

\begin{lemma}
\label{curves-lemma-factor-almost-isomorphism}
$k$ を代数閉体とする．$f : X' \to X$ を代数的 $k$-概形の有限射とし，
$\mathcal{O}_X \subset f_*\mathcal{O}_{X'}$ で，$f$ は有限個の点を除いて
同型であるとする．このとき分解
$$
X' = X_n \to X_{n - 1} \to \ldots \to X_1 \to X_0 = X
$$
が存在し，各 $X_i \to X_{i - 1}$ は二点の貼り合わせまたは接ベクトルの
押し潰しのいずれかである（例 \ref{curves-example-glue-points} および
\ref{curves-example-squish-tangent-vector} を参照せよ）．
\end{lemma}

\begin{proof}
$U \subset X$ を，その上で $f$ が同型となる最大の開集合とする．すると
$X \setminus U = \{x_1, \ldots, x_n\}$ で，$x_i \in X(k)$ である．
分解 $X' \to Y \to X$ であって，$f$ の分解であり，両射が有限で，
$$
\mathcal{O}_X \subset g_*\mathcal{O}_Y \subset f_*\mathcal{O}_{X'}
$$
を満たすものを考える．ここで $g : Y \to X$ は与えられた射である．仮定により，
$\mathcal{O}_{X, x} \to (f_*\mathcal{O}_{X'})_x$ は，$x = x_i$ となる $i$ が
存在しない限り同型である．したがって余核
$$
f_*\mathcal{O}_{X'}/\mathcal{O}_X = \bigoplus \mathcal{Q}_i
$$
は，摩天楼層 $\mathcal{Q}_i$ であって $x_1, \ldots, x_n$ に台をもつものの直和である．
表示された商は連接 $\mathcal{O}_X$-加群なので，$\mathcal{Q}_i$ は
$\mathcal{O}_{X, x_i}$ 上有限長である．したがって，これらの長さの和，すなわち
余核全体の長さに関する帰納法を用いることができる．

\medskip\noindent
$n > 1$ ならば，$\mathcal{O}_X$-部分代数
$\mathcal{A} \subset f_*\mathcal{O}_{X'}$ を，$\mathcal{Q}_1$ の逆像をとることにより定義できる．
これにより非自明な分解が得られ，帰納法から結論が従う．

\medskip\noindent
$n = 1$ と仮定する．$x = x_1$ と略記する．有限 $k$-代数拡大
$$
A = \mathcal{O}_{X, x} \subset (f_*\mathcal{O}_{X'})_x = B
$$
を考える．$\mathcal{Q} = \mathcal{Q}_1$ は値 $B/A$ をもつ摩天楼層であることに
注意する．$k$-部分代数 $A \subset A + \mathfrak m_A B \subset B$ がある．
両包含が真ならば非自明な分解が得られ，上と同様に帰納法から結論が従う．
$A + \mathfrak m_A B = B$ ならば，Nakayama により $A = B$ であり，したがって
$f$ は同型で，証明すべきことはない．よって $B = A + \mathfrak m_A B$ と
仮定してよい．$C = B/\mathfrak m_A B$ とおく．$C$ が $2$ 個より多くの
$k$-部分代数をもつならば，$A$ と $B$ の間の部分代数を，$B$ で逆像を
とることにより得る．したがって $C$ はちょうど $2$ 個の $k$-部分代数をもつと
仮定してよい．補題 \ref{curves-lemma-no-in-between-over-k} により
$C = k \times k$ または $C = k[\epsilon]$ である．この場合，対応して $f$ は
二点の貼り合わせまたは接ベクトルの押し潰しである．
\end{proof}

\begin{lemma}
\label{curves-lemma-glue-points}
$k$ を代数閉体とする．$f : X' \to X$ が例 \ref{curves-example-glue-points} のような
二点 $a, b$ の貼り合わせならば，完全列
$$
k^* \to \Pic(X) \to \Pic(X') \to 0
$$
が存在する．最初の写像は $a$ と $b$ が $X'$ の異なる連結成分上にあれば零であり，
$X'$ が固有で $a$ と $b$ が $X'$ の同じ連結成分上にあれば単射である．
\end{lemma}

\begin{proof}
Varieties，補題 \ref{varieties-lemma-surjective-pic-birational-finite} により，
写像 $\Pic(X) \to \Pic(X')$ は全射である．短完全列
$$
0 \to \mathcal{O}_X^* \to f_*\mathcal{O}_{X'}^*
\xrightarrow{ab^{-1}} x_*k^* \to 0
$$
を用いると，
$$
H^0(X', \mathcal{O}_{X'}^*) \xrightarrow{ab^{-1}} k^* \to
H^1(X, \mathcal{O}_X^*) \to H^1(X, f_*\mathcal{O}_{X'}^*)
$$
を得る．$H^1(X, f_*\mathcal{O}_{X'}^*) \subset H^1(X', \mathcal{O}_{X'}^*)$ である
（例えば Leray スペクトル系列による；Cohomology，補題
\ref{cohomology-lemma-Leray} を参照せよ）．したがって
$\Pic(X) \to \Pic(X')$ の核は
$ab^{-1} : H^0(X', \mathcal{O}_{X'}^*) \to k^*$ の余核である．
$a$ と $b$ が $X'$ の異なる連結成分上にあれば，$ab^{-1}$ は全射である．
$k$ は代数閉なので，$k$ 上の被約連結固有概形上の任意の正則関数は $k$ の
元から来る．Varieties，補題
\ref{varieties-lemma-proper-geometrically-reduced-global-sections} を参照せよ．
したがって $ab^{-1}$ は，$X'$ が固有で，$a$ と $b$ が同じ連結成分上にあれば
零である．
\end{proof}

\begin{lemma}
\label{curves-lemma-squish-tangent-vector}
$k$ を代数閉体とする．$f : X' \to X$ が例
\ref{curves-example-squish-tangent-vector} のような接ベクトル $\vartheta$ の
押し潰しならば，完全列
$$
(k, +) \to \Pic(X) \to \Pic(X') \to 0
$$
が存在し，$X'$ が固有かつ被約ならば最初の写像は単射である．
\end{lemma}

\begin{proof}
Varieties，補題 \ref{varieties-lemma-surjective-pic-birational-finite} により，
写像 $\Pic(X) \to \Pic(X')$ は全射である．例
\ref{curves-example-squish-tangent-vector} の短完全列
$$
0 \to \mathcal{O}_X^* \to f_*\mathcal{O}_{X'}^*
\xrightarrow{\vartheta} x_*k \to 0
$$
を用いると，
$$
H^0(X', \mathcal{O}_{X'}^*) \xrightarrow{\vartheta} k \to
H^1(X, \mathcal{O}_X^*) \to H^1(X, f_*\mathcal{O}_{X'}^*)
$$
を得る．$H^1(X, f_*\mathcal{O}_{X'}^*) \subset H^1(X', \mathcal{O}_{X'}^*)$ である
（例えば Leray スペクトル系列による；Cohomology，補題
\ref{cohomology-lemma-Leray} を参照せよ）．したがって
$\Pic(X) \to \Pic(X')$ の核は写像
$\vartheta : H^0(X', \mathcal{O}_{X'}^*) \to k$ の余核である．$k$ は代数閉なので，
$k$ 上の被約連結固有概形上の任意の正則関数は $k$ の元から来る．
Varieties，補題
\ref{varieties-lemma-proper-geometrically-reduced-global-sections} を参照せよ．
これで補題の最後の主張が従う．
\end{proof}


\section{多重交差特異点と節点特異点}
\label{curves-section-multicross}

\noindent
本節では，曲線に現れうる最も単純な特異点について論じる．

\medskip\noindent
$k$ を体とする．完備局所 $k$-代数
\begin{equation}
\label{curves-equation-multicross}
A = \{(f_1, \ldots, f_n) \in k[[t]] \times \ldots \times k[[t]] \mid
f_1(0) = \ldots = f_n(0)\}
\end{equation}
を考える．Varieties，定義 \ref{varieties-definition-wedge} で導入した言葉では，
$A$ は $n$ 個の，$1$ 変数の形式的冪級数環を $k$ 上でくさび和にしたものである．
$k[[t]] \times \ldots \times k[[t]]$
は $A$ の全商環における整閉包であることに注意する．したがって
$\delta$-不変量の $A$ における値は $n - 1$ である．同型
$$
k[[x_1, \ldots, x_n]]/(\{x_ix_j\}_{i \not = j}) \longrightarrow A
$$
があり，これは $x_i$ を $(0, \ldots, 0, t, 0, \ldots, 0)$ という $A$ の元に写すことで
得られる．したがって $\dim(A) = 1$ かつ
$\dim_k \mathfrak m/\mathfrak m^2 = n$ である．特に，$A$ が正則であることと
$n = 1$ は同値である．

\begin{lemma}
\label{curves-lemma-multicross-algebra}
$k$ を分離閉体とする．$A$ を $1$ 次元被約 Nagata 局所 $k$-代数で，剰余体が
$k$ であるものとする．このとき
$$
\delta\text{-不変量 }A \geq \text{枝数 }A - 1
$$
等号が成り立つならば，$A^\wedge$ は (\ref{curves-equation-multicross}) の形である．
\end{lemma}

\begin{proof}
$A$ の剰余体は分離閉なので，$A$ の枝数は $A$ の幾何学的枝数に等しい．
More on Algebra，定義 \ref{more-algebra-definition-number-of-branches} を参照せよ．
不等式は Varieties，補題
\ref{varieties-lemma-delta-number-branches-inequality} から従う．等号が成り立つと
仮定する．$A$ を $A$ の完備化で置き換えてよい．これは枝数も
$\delta$-不変量も変えない．More on Algebra，補題
\ref{more-algebra-lemma-one-dimensional-number-of-branches} および Varieties，補題
\ref{varieties-lemma-delta-same-after-completion} を参照せよ．すると $A$ は強 Hensel
環である．Algebra，補題 \ref{algebra-lemma-complete-henselian} を参照せよ．
Varieties，補題 \ref{varieties-lemma-delta-number-branches-inequality-sh} により，
$A$ は完備離散付値環のくさび和である．その各々は Algebra，補題
\ref{algebra-lemma-regular-complete-containing-coefficient-field} により
$k[[t]]$ と同型である．したがって $A$ は (\ref{curves-equation-multicross}) の形である．
\end{proof}

\begin{definition}
\label{curves-definition-multicross}
$k$ を代数閉体とする．$X$ を代数的 $1$ 次元 $k$-概形とする．
$x \in X$ を閉点とする．$x$ は，完備化 $\mathcal{O}_{X, x}^\wedge$ が
(\ref{curves-equation-multicross}) と同型であるような $n \geq 2$ が存在するとき，
{\it 多重交差特異点を定める}という．$x$ は，$n = 2$ のとき，
{\it 節点}，{\it 通常二重点}，または {\it 節点特異点を定める}という．
\end{definition}

\noindent
これらの特異点は，ある意味で，代数閉体上の曲線に現れうる最も単純な種類の
特異点である．

\begin{lemma}
\label{curves-lemma-multicross}
$k$ を代数閉体とする．$X$ を被約な代数的 $1$ 次元 $k$-概形とする．
$x \in X$ とする．次は同値である．
\begin{enumerate}
\item $x$ は多重交差特異点を定める；
\item $\delta$-不変量の $X$ の $x$ における値は，$X$ の $x$ における枝数から
$1$ を引いたものに等しい；
\item 射の列
$U_n \to U_{n - 1} \to \ldots \to U_0 = U \subset X$
が存在する．ここで $U$ は $x$ の開近傍，$U_n$ は非特異であり，各
$U_i \to U_{i - 1}$ は例 \ref{curves-example-glue-points} のような二点の貼り合わせである．
\end{enumerate}
\end{lemma}

\begin{proof}
(1) と (2) の同値性は補題 \ref{curves-lemma-multicross-algebra} である．

\medskip\noindent
(3) を仮定する．$i$ に関する降下帰納法により，$U_i$ のすべての特異点が
多重交差であることを示す．これは $U_n$ に対して成り立つ．実際，$U_n$ は
特異点をもたない．$U_i$ が $U_{i + 1}$ の $a, b \in U_{i + 1}$ を一点
$c \in U_i$ に貼り合わせて得られるならば，
$$
\mathcal{O}_{U_i, c}^\wedge \subset
\mathcal{O}_{U_{i + 1}, a}^\wedge \times \mathcal{O}_{U_{i + 1}, b}^\wedge
$$
は $k$ における剰余類が同じである元全体の集合である．したがって $c$ における
枝数は $a$ と $b$ における枝数の和であり，$\delta$-不変量の $c$ における値は
$\delta$-不変量の $a$ と $b$ における値の和に $1$ を加えたものである
（表示された包含の余次元が $1$ だからである）．これで望みどおり (2) が成り立つ．

\medskip\noindent
同値な条件 (1) と (2) を仮定する．$U \subset X$ を，$x$ が $U$ の唯一の
特異点となる開集合として選んでよい．次に正規化射
$$
U^\nu = U_n \to U_{n - 1} \to \ldots \to U_1 \to U_0 = U
$$
に補題 \ref{curves-lemma-factor-almost-isomorphism} を適用する．示すべきことは，
どの段階でも接ベクトルを押し潰していないことである．
$U_{i + 1} \to U_i$ が接ベクトルを押し潰す段階となる最小の $i$ であると
仮定する．その接ベクトルを $\theta$ とし，$u' \in U_{i + 1}$ にあり，
$u \in U_i$ の上にあるとする．上と同様に論じると，$u_i$ は多重交差特異点で
あることが分かる（写像 $U_i \to \ldots \to U_0$ は点の対の貼り合わせだからで
ある）．しかし今，$u'$ と $u$ における枝数は等しく，$\delta$-不変量の
$U_i$ の $u$ における値は $1$ だけ，$\delta$-不変量の $U_{i + 1}$ の
$u'$ における値より大きい．
補題 \ref{curves-lemma-multicross-algebra} により，これは $u$ が多重交差特異点では
ありえないことを意味し，矛盾である．
\end{proof}

\begin{lemma}
\label{curves-lemma-multicross-gorenstein-is-nodal}
$k$ を代数閉体とする．$X$ を被約な代数的 $1$ 次元 $k$-概形とする．
$x \in X$ を多重交差特異点とする（定義 \ref{curves-definition-multicross}）．
$X$ が Gorenstein ならば，$x$ は節点である．
\end{lemma}

\begin{proof}
写像 $\mathcal{O}_{X, x} \to \mathcal{O}_{X, x}^\wedge$ は平坦であり，
$\kappa(x) = \mathcal{O}_{X, x}^\wedge/\mathfrak m_x \mathcal{O}_{X, x}^\wedge$
という意味で不分岐である（More on Algebra，節
\ref{more-algebra-section-permanence-completion} を参照せよ）．したがって，
$X$ が Gorenstein ならば $\mathcal{O}_{X, x}$ は Gorenstein であり，さらに
Dualizing Complexes，補題 \ref{dualizing-lemma-flat-under-gorenstein} により
$\mathcal{O}_{X, x}^\wedge$ は Gorenstein である．よって
(\ref{curves-equation-multicross}) の環 $A$ で $n \geq 2$ なるものが Gorenstein である
ことと $n = 2$ が同値であることを示せば十分である．

\medskip\noindent
$n = 2$ ならば，$A = k[[x, y]]/(xy)$ は完全交差であり，したがって Gorenstein
である．例えばこれは，$k[[x, y]] \to A$ に Duality for Schemes，補題
\ref{duality-lemma-gorenstein-lci} を適用することと，正則局所環 $k[[x, y]]$ が
Dualizing Complexes，補題 \ref{dualizing-lemma-regular-gorenstein} により
Gorenstein であることから従う．

\medskip\noindent
$n > 2$ と仮定する．$A$ が Gorenstein ならば，$A$ は $A$ 上の双対化複体と
なるはずである（Duality for Schemes，定義
\ref{duality-definition-gorenstein}）．すると $R\Hom(k, A)$ は $k[n]$ に
等しくなるはずであり，ここで $n \in \mathbf{Z}$ である．Dualizing Complexes，補題
\ref{dualizing-lemma-find-function} を参照せよ．したがって
$\Ext^1_A(k, A) \cong k$ または $\Ext^1_A(k, A) = 0$ となるはずである
（$n$ の値に応じる．実際には $n$ は $-1$ でなければならないが，ここでは
重要でない）．完全列
$$
0 \to \mathfrak m_A \to A \to k \to 0
$$
を用いると，
$$
\Ext^1_A(k, A) = \Hom_A(\mathfrak m_A, A)/A
$$
を得る．ここで $A \to \Hom_A(\mathfrak m_A, A)$ は
$a \mapsto (a' \mapsto aa')$ によって与えられる．
$e_i \in \Hom_A(\mathfrak m_A, A)$ を，
$(f_1, \ldots, f_n) \in \mathfrak m_A$ を
$(0, \ldots, 0, f_i, 0, \ldots, 0)$ に写す元とする．読者は容易に，
$e_1, \ldots, e_{n - 1}$ が $k$ 上線形独立であり，その独立性が
$\Hom_A(\mathfrak m_A, A)/A$ の中で成り立つことを確かめられる．したがって
$\dim_k \Ext^1_A(k, A) \geq n - 1 \geq 2$ であり，証明が完了する．
（$e_1 + \ldots + e_n$ は $1$ の，写像 $A \to \Hom_A(\mathfrak m_A, A)$ による
像であることに注意せよ．）
\end{proof}




\section{Picard 群の捩れ}
\label{curves-section-torsion-in-pic}

\noindent
本節では，体上の $1$ 次元固有概形の Picard 群における捩れを評価する．
これは曲線の半安定還元を調べる際に用いる．

\medskip\noindent
補題 \ref{curves-lemma-torsion-picard-smooth-projective} の結果を初等的に得る方法は
ないように見える．証明を分析すると，二つの主要な材料がある．
(1) 滑らかな射影曲線上の次数零の可逆層を分類するアーベル多様体が存在すること，
(2) アーベル多様体上の捩れ点の構造を決定できることである．

\begin{lemma}
\label{curves-lemma-torsion-picard-smooth-projective}
$k$ を代数閉体とする．$X$ を種数 $g$ の $k$ 上の滑らかな射影曲線とする．
\begin{enumerate}
\item $n \geq 1$ が $k$ で可逆ならば，
$\Pic(X)[n] \cong (\mathbf{Z}/n\mathbf{Z})^{\oplus 2g}$.
\item $k$ の標数が $p > 0$ ならば，整数 $0 \leq f \leq g$ が存在し，
$\Pic(X)[p^m] \cong (\mathbf{Z}/p^m\mathbf{Z})^{\oplus f}$ であり，
これはすべての $m \geq 1$ に対して成り立つ．
\end{enumerate}
\end{lemma}

\begin{proof}
$\Pic^0(X) \subset \Pic(X)$ を次数 $0$ の可逆層からなる部分群とする．
言い換えれば，短完全列
$$
0 \to \Pic^0(X) \to \Pic(X) \xrightarrow{\deg} \mathbf{Z} \to 0.
$$
が存在する．群 $\Pic^0(X)$ は $k$-点の群であり，その群概形は
$\underline{\Picardfunctor}^0_{X/k}$ である．Picard Schemes of Curves，補題
\ref{pic-lemma-picard-pieces} を参照せよ．同じ補題により，
$\underline{\Picardfunctor}^0_{X/k}$ は Groupoids，定義
\ref{groupoids-definition-abelian-variety} の意味で $g$ 次元の $k$ 上の
アーベル多様体である．したがって Groupoids，命題
\ref{groupoids-proposition-review-abelian-varieties} の結果から結論が従う．
\end{proof}

\begin{lemma}
\label{curves-lemma-torsion-picard-becomes-visible}
$k$ を体とする．$n$ は $k$ の標数と互いに素であるとする．
$X$ を $k$ 上の滑らかな固有曲線で，$H^0(X, \mathcal{O}_X) = k$ かつ
種数 $g$ のものとする．
\begin{enumerate}
\item $g = 1$ ならば，有限分離拡大 $k'/k$ が存在し，$X_{k'}$ は
$k'$-有理点をもち，
$\Pic(X_{k'})[n] \cong (\mathbf{Z}/n\mathbf{Z})^{\oplus 2}$.
\item $g \geq 2$ ならば，有限分離拡大 $k'/k$ で
$[k' : k] \leq (2g - 2)(n^{2g})!$ を満たすものが存在し，$X_{k'}$ は
$k'$-有理点をもち，
$\Pic(X_{k'})[n] \cong (\mathbf{Z}/n\mathbf{Z})^{\oplus 2g}$.
\end{enumerate}
\end{lemma}

\begin{proof}
$g \geq 2$ と仮定する．まず補題 \ref{curves-lemma-point-over-separable-extension} により，
次数が高々 $2g - 2$ の有限分離拡大を選び，その上で $X$ が有理点をもつように
できる．したがって $X$ が $k$-有理点 $x \in X(k)$ をもつと仮定してよいが，
今度は $[k' : k] \leq (n^{2g})!$ という評価のもとで補題を証明しなければならない．
$k \subset k^{sep} \subset \overline{k}$ を，代数閉包の中の分離代数閉包とする．
補題 \ref{curves-lemma-torsion-picard-smooth-projective} により，
$$
\Pic(X_{\overline{k}})[n] \cong (\mathbf{Z}/n\mathbf{Z})^{\oplus 2g}
$$
Picard Schemes of Curves，補題 \ref{pic-lemma-torsion-descends} により，
$$
\Pic(X_{k^{sep}})[n] \cong (\mathbf{Z}/n\mathbf{Z})^{\oplus 2g}
$$
を得る．Picard Schemes of Curves，補題 \ref{pic-lemma-torsion-descends} により，
連続作用
$$
\text{Gal}(k^{sep}/k)
\longrightarrow
\text{Aut}(\Pic(X_{k^{sep}})[n]
$$
がある．この写像の核の固定体 $k'$ に対して補題が成り立つ．作用は連続なので
核は開であり，したがって $k'/k$ は有限である．Galois 理論により
$\text{Gal}(k'/k)$ は表示された射の像である．濃度 $n^{2g}$ の集合の置換群は
濃度 $(n^{2g})!$ をもつので，Galois 理論により
$[k' : k] \leq (n^{2g})!$ と結論する．（もちろん，これで
$|\text{GL}_{2g}(\mathbf{Z}/n\mathbf{Z})|$ という評価をもつ補題が証明されるが，
ここで必要なのは何らかの評価が存在することだけである．）

\medskip\noindent
種数が $1$ ならば，$X$ が有理点をもつようになる有限分離体拡大の次数には
上界がない（詳細は省略する）．それでも，例えば Varieties，補題
\ref{varieties-lemma-smooth-separable-closed-points-dense} により，そのような拡大は
存在する．証明の残りは $g \geq 2$ の場合と同じである．
\end{proof}

\begin{proposition}
\label{curves-proposition-torsion-picard-reduced-proper}
$k$ を代数閉体とする．$X$ を $k$ 上の固有概形で，被約，連結，かつ次元 $1$ の
ものとする．$g$ を $X$ の種数，$g_{geom}$ を $X$ の既約成分の幾何種数の和とする．
$\ell$ を $k$ の標数と異なる任意の素数とすると，
$$
\dim_{\mathbf{F}_\ell} \Pic(X)[\ell]
\leq g + g_{geom}
$$
であり，等号が成り立つことと $X$ のすべての特異点が多重交差であることは
同値である．
\end{proposition}

\begin{proof}
$\nu : X^\nu \to X$ を正規化とする
（Varieties，補題 \ref{varieties-lemma-prepare-delta-invariant}）．分解
$$
X^\nu = X_n \to X_{n - 1} \to \ldots \to X_1 \to X_0 = X
$$
を補題 \ref{curves-lemma-factor-almost-isomorphism} のように選ぶ．
$h^0_i = \dim_k H^0(X_i, \mathcal{O}_{X_i})$ および
$h^1_i = \dim_k H^1(X_i, \mathcal{O}_{X_i})$ と書く．補題
\ref{curves-lemma-glue-points} および \ref{curves-lemma-squish-tangent-vector} により，各
$n > i \geq 0$ に対して次の三つの可能性のいずれかが成り立つ．
\begin{enumerate}
\item $X_i$ は異なる連結成分上にある $a, b \in X_{i + 1}$ を貼り合わせて得られる．
この場合，
$\Pic(X_i) = \Pic(X_{i + 1})$,
$h^0_{i + 1} = h^0_i + 1$, $h^1_{i + 1} = h^1_i$,
\item $X_i$ は同じ連結成分上にある $a, b \in X_{i + 1}$ を貼り合わせて得られる．
この場合，短完全列
$$
0 \to k^* \to \Pic(X_i) \to \Pic(X_{i + 1}) \to 0,
$$
があり，$h^0_{i + 1} = h^0_i$，$h^1_{i + 1} = h^1_i - 1$ である；
\item $X_i$ は $X_{i + 1}$ における接ベクトルを押し潰して得られる．
この場合，短完全列
$$
0 \to (k, +) \to \Pic(X_i) \to \Pic(X_{i + 1}) \to 0,
$$
があり，$h^0_{i + 1} = h^0_i$，$h^1_{i + 1} = h^1_i - 1$ である．
\end{enumerate}
構造層のコホモロジー群の次元に関する主張を証明するには，例
\ref{curves-example-glue-points} および \ref{curves-example-squish-tangent-vector} の完全列を用いる．
$k$ は代数閉で，その標数は $\ell$ と互いに素なので，$(k, +)$ と $k^*$ は
$\ell$-可除であり，その $\ell$-捩れは $(k, +)[\ell] = 0$ および
$k^*[\ell] \cong \mathbf{F}_\ell$ である．したがって
$$
\dim_{\mathbf{F}_\ell} \Pic(X_{i + 1})[\ell] -
\dim_{\mathbf{F}_\ell}\Pic(X_i)[\ell]
$$
は，場合 (2) を除いて零であり，場合 (2) では $-1$ に等しい．この過程の最後に，
正規化 $X^\nu = X_n$ を得る．これは滑らかな射影曲線の $k$ 上の非交和である．したがって
\begin{enumerate}
\item $h^1_n = g_{geom}$ であり，
\item $\dim_{\mathbf{F}_\ell} \Pic(X_n)[\ell] = 2g_{geom}$.
\end{enumerate}
最後の等式は補題 \ref{curves-lemma-torsion-picard-smooth-projective} による．
$g = h^1_0$ なので，型 (2) と (3) の段階の数は高々
$h^1_0 - h^1_n = g - g_{geom}$ である．階数の差についての計算から，
$$
\dim_{\mathbf{F}_\ell} \Pic(X)[\ell] \leq
g - g_{geom} + 2g_{geom} = g + g_{geom}
$$
と結論する．等号が成り立つことと，型 (3) の段階が一つも現れないことは同値である．
補題 \ref{curves-lemma-multicross} により，これはまさに $X$ のすべての特異点が
多重交差であることを意味する．逆に，すべての特異点が多重交差ならば，補題
\ref{curves-lemma-multicross} により，上のような列
$X^\nu = X_n \to \ldots \to X_0 = X$ であって型 (3) の段階を含まないものが
存在する．すると補題で等号が成り立つ（各点に対する局所的な列を貼り合わせて，
$x$ のすべての特異点に対して働く一つの列を得ればよい；詳細は省略する）．
\end{proof}





\section{種数と幾何種数}
\label{curves-section-genus-geometric-genus}

\noindent
$k$ を体とし，その代数閉包を $\overline{k}$ とする．
$X$ を次元 $\leq 1$ の $k$ 上の固有概形とする．
$g_{geom}(X/k)$ を，$X_{\overline{k}}$ の次元 $1$ の既約成分の
幾何種数の和と定義する．

\begin{lemma}
\label{curves-lemma-bound-geometric-genus}
$k$ を体とする．$X$ を次元 $\leq 1$ の $k$ 上の固有概形とする．このとき
$$
g_{geom}(X/k) = \sum\nolimits_{C \subset X} g_{geom}(C/k)
$$
である．ここで和は，既約成分 $C \subset X$ のうち次元が $1$ のもの全体にわたる．
\end{lemma}

\begin{proof}
これは定義と次の事実から直ちに従う．すなわち，既約成分 $\overline{Z}$ が
$X_{\overline{k}}$ に含まれるなら，これは既約成分 $Z$ の上へ写り，その像は $X$ の既約成分である
（Varieties，補題 \ref{varieties-lemma-image-irreducible}），しかも同じ次元をもつ
（Morphisms，補題 \ref{morphisms-lemma-dimension-fibre-after-base-change} を
$\overline{Z}$ の生成点に適用する）．
\end{proof}

\begin{lemma}
\label{curves-lemma-geometric-genus-normalization}
$k$ を体とする．$X$ を次元 $\leq 1$ の $k$ 上の固有概形とする．このとき
\begin{enumerate}
\item $g_{geom}(X/k) = g_{geom}(X_{red}/k)$ である．
\item $X' \to X$ が双有理固有射ならば，
$g_{geom}(X'/k) = g_{geom}(X/k)$ である．
\item $X^\nu \to X$ が正規化射ならば，
$g_{geom}(X^\nu/k) = g_{geom}(X/k)$ である．
\end{enumerate}
\end{lemma}

\begin{proof}
(1) は補題 \ref{curves-lemma-bound-geometric-genus} から直ちに従う．
$X' \to X$ が固有双有理ならば，これは有限であり，ある稠密開集合上で
同型である（Varieties，補題 \ref{varieties-lemma-finite-in-codim-1} および
\ref{varieties-lemma-modification-normal-iso-over-codimension-1} を参照）．
したがって $X'_{\overline{k}} \to X_{\overline{k}}$ はある稠密開集合上で同型である．
ゆえに $X'_{\overline{k}}$ と $X_{\overline{k}}$ の既約成分は全単射に対応し，
対応する成分の函数体は同型である．特に，これらの成分の非特異射影モデルは
同型であり，したがって幾何種数も等しい．これで (2) が示された．
(3) は (1)，(2)，および $X^\nu \to X_{red}$ が双有理であること
（Morphisms，補題 \ref{morphisms-lemma-normalization-birational}）から従う．
\end{proof}

\begin{lemma}
\label{curves-lemma-genus-goes-down}
$k$ を体とする．$X$ を次元 $\leq 1$ の $k$ 上の固有概形とする．
$f : Y \to X$ を有限射とし，ある稠密開集合 $U \subset X$ 上で
$f$ が閉浸入になると仮定する．このとき
$$
\dim_k H^1(X, \mathcal{O}_X) \geq \dim_k H^1(Y, \mathcal{O}_Y)
$$
である．
\end{lemma}

\begin{proof}
完全列
$$
0 \to \mathcal{G} \to \mathcal{O}_X \to f_*\mathcal{O}_Y \to \mathcal{F} \to 0
$$
を考える．これは $X$ 上の連接層の完全列である．仮定により $\mathcal{F}$ の台は有限個の閉点に含まれ，したがって
その高次コホモロジーは消える
（Varieties，補題 \ref{varieties-lemma-chi-tensor-finite}）．
一方，Cohomology，命題 \ref{cohomology-proposition-vanishing-Noetherian} により
$H^2(X, \mathcal{G}) = 0$ である．したがって形式的に，誘導される写像
$H^1(X, \mathcal{O}_X) \to H^1(X, f_*\mathcal{O}_Y)$
は全射である．さらに
$H^1(X, f_*\mathcal{O}_Y) = H^1(Y, \mathcal{O}_Y)$ なので
（Cohomology of Schemes，補題 \ref{coherent-lemma-relative-affine-cohomology}），
補題の結論を得る．
\end{proof}

\begin{lemma}
\label{curves-lemma-genus-normalization}
$k$ を体とする．$X$ を次元 $\leq 1$ の $k$ 上の固有概形とする．
$X' \to X$ が双有理固有射ならば，
$$
\dim_k H^1(X, \mathcal{O}_X) \geq \dim_k H^1(X', \mathcal{O}_{X'})
$$
である．さらに，$X$ が被約で，
$H^0(X, \mathcal{O}_X) \to H^0(X', \mathcal{O}_{X'})$ が全射であり，
上の不等式で等号が成り立つならば，$X' = X$ である．
\end{lemma}

\begin{proof}
$f : X' \to X$ が固有双有理ならば，これは有限であり，ある稠密開集合上で
同型である（Varieties，補題 \ref{varieties-lemma-finite-in-codim-1} および
\ref{varieties-lemma-modification-normal-iso-over-codimension-1} を参照）．
したがって補題 \ref{curves-lemma-genus-goes-down} により不等式を得る．
$X$ が被約であると仮定する．このとき
$\mathcal{O}_X \to f_*\mathcal{O}_{X'}$ は単射であり，短完全列
$$
0 \to \mathcal{O}_X \to f_*\mathcal{O}_{X'} \to \mathcal{F} \to 0
$$
を得る．第二の主張の仮定のもと，長完全コホモロジー列から
$H^0(X, \mathcal{F}) = 0$ を得る．すると $\mathcal{F} = 0$ である．実際，
$\mathcal{F}$ は大域切断で生成される
（Varieties，補題 \ref{varieties-lemma-chi-tensor-finite}）．
したがって $\mathcal{O}_X = f_*\mathcal{O}_{X'}$ である．$f$ はアフィンだから，
これは $X = X'$ を意味する．
\end{proof}

\begin{lemma}
\label{curves-lemma-bound-geometric-genus-curve}
$k$ を体とする．$C$ を $k$ 上の固有曲線とする．
$\kappa = H^0(C, \mathcal{O}_C)$ とおく．このとき
$$
[\kappa : k]_s \dim_\kappa H^1(C, \mathcal{O}_C) \geq g_{geom}(C/k)
$$
である．
\end{lemma}

\begin{proof}
Varieties，補題 \ref{varieties-lemma-regular-functions-proper-variety} により，
$\kappa$ は体であり，$k$ の有限次拡大である．Fields，補題
\ref{fields-lemma-separable-degree} により
$[\kappa : k]_s = |\Mor_k(\kappa, \overline{k})|$ であり，したがって
$\Spec(\kappa \otimes_k \overline{k})$ は $[\kappa : k]_s$ 個の点をもち，
それぞれの剰余体は $\overline{k}$ である．ゆえに
$$
C_{\overline{k}} =
\bigcup\nolimits_{t \in \Spec(\kappa \otimes_k \overline{k})} C_t
$$
である（集合論的な合併）．ここで
$C_t = C \times_{\Spec(\kappa), t} \Spec(\overline{k})$ であり，
$t$ は $k$-代数写像 $t : \kappa \to \overline{k}$ とみなす．したがって
$g_{geom}(C/k) = \sum_t g_{geom}(C_t/\overline{k})$ であり，この和の添字集合の
濃度は $[\kappa : k]_s$ である．また
$$
H^0(C_t, \mathcal{O}_{C_t}) = \overline{k}
\quad\text{かつ}\quad
\dim_{\overline{k}} H^1(C_t, \mathcal{O}_{C_t}) =
\dim_\kappa H^1(C, \mathcal{O}_C)
$$
である．これはコホモロジーと基底変換
（Cohomology of Schemes，補題 \ref{coherent-lemma-flat-base-change-cohomology}）による．
正規化 $C_t^\nu$ は，$C_t$ の既約成分の非特異射影モデルの非交和である
（Morphisms，補題 \ref{morphisms-lemma-normalization-in-terms-of-components}）．
したがって $\dim_{\overline{k}} H^1(C_t^\nu, \mathcal{O}_{C_t^\nu})$ は
$g_{geom}(C_t/\overline{k})$ に等しい．補題 \ref{curves-lemma-genus-goes-down} により
$$
\dim_{\overline{k}} H^1(C_t, \mathcal{O}_{C_t}) \geq
\dim_{\overline{k}} H^1(C_t^\nu, \mathcal{O}_{C_t^\nu})
$$
であり，これで証明は完了する．
\end{proof}

\begin{lemma}
\label{curves-lemma-bound-torsion-simple}
$k$ を体とする．$X$ を次元 $\leq 1$ の $k$ 上の固有概形とする．
$\ell$ を $k$ で可逆な素数とする．このとき
$$
\dim_{\mathbf{F}_\ell} \Pic(X)[\ell] \leq
\dim_k H^1(X, \mathcal{O}_X) + g_{geom}(X/k)
$$
である．ここで $g_{geom}(X/k)$ は上で定義したものである．
\end{lemma}

\begin{proof}
写像 $\Pic(X) \to \Pic(X_{\overline{k}})$ は，Varieties，補題
\ref{varieties-lemma-change-fields-pic} により単射である．
Cohomology of Schemes，補題 \ref{coherent-lemma-flat-base-change-cohomology} により，
$\dim_k H^1(X, \mathcal{O}_X)$ は
$\dim_{\overline{k}} H^1(X_{\overline{k}}, \mathcal{O}_{X_{\overline{k}}})$ に等しい．
したがって $k$ は代数閉であると仮定してよい．

\medskip\noindent
$X_{red}$ を $X$ の被約化とする．このとき全射
$\mathcal{O}_X \to \mathcal{O}_{X_{red}}$ は全射
$H^1(X, \mathcal{O}_X) \to H^1(X, \mathcal{O}_{X_{red}})$
を誘導する．実際，Cohomology，命題
\ref{cohomology-proposition-vanishing-Noetherian} により，擬連接層のコホモロジーは
次数 $\geq 2$ で消える．$X_{red} \to X$ は $\overline{k}$ 上の既約成分について
同型を誘導し，Picard 群の $\ell$-捩れについても同型を誘導するので
（Picard Schemes of Curves，補題 \ref{pic-lemma-torsion-descends}），
$X$ を $X_{red}$ で置き換えてよい．こうして命題
\ref{curves-proposition-torsion-picard-reduced-proper} に帰着する．
\end{proof}






\section{節点曲線}
\label{curves-section-nodal}

\noindent
代数閉体上の通常二重点はすでに定義した．定義 \ref{curves-definition-multicross} を参照せよ．
すなわち，$x \in X$ が次元 $1$ の代数概形の閉点であり，その概形が代数閉体 $k$ 上にあるならば，
$x$ が通常二重点であることと
$$
\mathcal{O}_{X, x}^\wedge \cong k[[x, y]]/(xy)
$$
であることは同値である．(\ref{curves-equation-multicross}) に続く議論
（第 \ref{curves-section-multicross} 節）を参照せよ．

\begin{definition}
\label{curves-definition-nodal}
$k$ を体とする．$X$ を次元 $1$ の局所代数的 $k$-概形とする．
\begin{enumerate}
\item 閉点 $x \in X$ へ写る通常二重点 $\overline{x} \in X_{\overline{k}}$ が
存在するとき，$x$ を{\it 節点}，または{\it 通常二重点}といい，あるいは
{\it 節点特異点を定める}という．
\item {\it $X$ の特異点が高々節点的である}とは，$X$ のすべての閉点が，
構造射 $X \to \Spec(k)$ の滑らかな軌跡に属するか，または通常二重点であることをいう．
\end{enumerate}
\end{definition}

\noindent
次元 $1$ の代数概形 $X$ の特異点が高々節点的であるとき，$X$ を
{\it 節点曲線}と呼ぶことが多い．節点曲線に固有性を要求することもある．
このように定義した節点曲線は既約とは限らないため，曲線を次元 $1$ の多様体とする
先の定義とは整合しない．

\begin{lemma}
\label{curves-lemma-reduced-quotient-regular-ring-dim-2}
$(A, \mathfrak m)$ を次元 $2$ の正則局所環とし，$I \subset \mathfrak m$ をイデアルとする．
\begin{enumerate}
\item $A/I$ が被約ならば，$I = (0)$，$I = \mathfrak m$，または
$I = (f)$ である．ここで $f \in \mathfrak m$ は零でない．
\item $A/I$ の深さが $1$ ならば，$I = (f)$ である．ここで
$f \in \mathfrak m$ は零でない．
\end{enumerate}
\end{lemma}

\begin{proof}
$I \not = 0$ と仮定し，$I = (f_1, \ldots, f_r)$ と書く．$A$ は UFD なので
（More on Algebra，補題 \ref{more-algebra-lemma-regular-local-UFD}），
$f_i = fg_i$ と書ける．ここで $f$ は $f_1, \ldots, f_r$ の最大公約元である．
したがって $g_1, \ldots, g_r$ の最大公約元は $1$ であり，これは高さ $1$ の
素イデアルで $g_1, \ldots, g_r$ を含むものがないことを意味する．
そこで $(g_1, \ldots, g_r) = A$ ならば $I = (f)$ であり，そうでなければ
$\dim(A) = 2$ から
$V(g_1, \ldots, g_r) = \{\mathfrak m\}$, i.e.,
$\mathfrak m = \sqrt{(g_1, \ldots, g_r)}$.
を得る．

\medskip\noindent
$A/I$ が被約，すなわち $I$ が根基的であると仮定する．$f$ が単元ならば，
$I$ が根基的であることから $I = \mathfrak m$ である．$f \in \mathfrak m$ ならば，
$f^n$ は $A/I$ で零に写る．よって被約性から $f \in I$ であり，$I = (f)$ を得る．

\medskip\noindent
$A/I$ の深さが $1$ であると仮定する．このとき $\mathfrak m$ は $A/I$ の
随伴素イデアルではない．$f$ の $I$ を法とする類は $g_1, \ldots, g_r$ で
零化されるから，$f$ の類は $A/I$ で零である．したがって所望のとおり
$I = (f)$ である．
\end{proof}

\noindent
$\kappa$ を体とし，$V$ を $\kappa$ 上のベクトル空間とする．
$q \in \text{Sym}^2_\kappa(V)$ について，誘導される $\kappa$-線形写像
$V^\vee \to V$ が同型であるとき，これを{\it 非退化}であるという．
$q = \sum_{i \leq j} a_{ij} x_i x_j$ であり，ある $\kappa$-基底
$x_1, \ldots, x_n$ が $V$ に選ばれているならば，これは行列
$$
\left(
\begin{matrix}
2a_{11} & a_{12} & \ldots \\
a_{12} & 2a_{22} & \ldots \\
\ldots & \ldots & \ldots
\end{matrix}
\right)
$$
の行列式が零でないことを意味する．これは，$q$ の $x_i$ に関する偏導函数が
$0$ を概形論的に切り出すという条件と同値である．

\begin{lemma}
\label{curves-lemma-nodal-algebraic}
$k$ を体とする．$(A, \mathfrak m, \kappa)$ を Noether 局所 $k$-代数とする．
次の条件は同値である．
\begin{enumerate}
\item $\kappa/k$ は分離的で，$A$ は被約であり，
$\dim_\kappa(\mathfrak m/\mathfrak m^2) = 2$ で，非退化元
$q \in \text{Sym}^2_\kappa(\mathfrak m/\mathfrak m^2)$ が存在し，これは
$\mathfrak m^2/\mathfrak m^3$ で零に写る．
\item $\kappa/k$ は分離的で，$\text{depth}(A) = 1$ であり，
$\dim_\kappa(\mathfrak m/\mathfrak m^2) = 2$ で，非退化元
$q \in \text{Sym}^2_\kappa(\mathfrak m/\mathfrak m^2)$ が存在し，これは
$\mathfrak m^2/\mathfrak m^3$ で零に写る．
\item $\kappa/k$ は分離的で，
$A^\wedge \cong \kappa[[x, y]]/(ax^2 + bxy + cy^2)$
は $k$-代数としての同型であり，$ax^2 + bxy + cy^2$ は $\kappa$ 上の
非退化二次形式である．
\end{enumerate}
\end{lemma}

\begin{proof}
(3) を仮定する．$A^\wedge$ は被約である．実際，$ax^2 + bxy + cy^2$ は
既約であるか，互いに同伴でない二つの素元の積である．したがって
$A \subset A^\wedge$ は被約であり，(1) が成り立つ．

\medskip\noindent
(1) を仮定する．$A$ は Artin 環ではありえない．実際，
$\mathfrak m \not = (0)$ なので，Artin 環ならば被約ではない．よって $\dim(A) \geq 1$ であり，Algebra，補題
\ref{algebra-lemma-criterion-reduced} により $\text{depth}(A) \geq 1$ である．
一方，$\dim(A) = 2$ ならば $A$ は正則であり，これは Algebra，補題
\ref{algebra-lemma-regular-graded} により $q$ の存在と矛盾する．
したがって $\dim(A) \leq 1$ であり，Algebra，補題
\ref{algebra-lemma-bound-depth} により $\text{depth}(A) = 1$ を得る．ゆえに (2) が成り立つ．

\medskip\noindent
(2) を仮定する．$A$ の深さは $A^\wedge$ の深さと等しく
（More on Algebra，補題 \ref{more-algebra-lemma-completion-depth}），他の条件も完備化で
変わらないから，$A$ は完備であると仮定してよい．More on Algebra，補題
\ref{more-algebra-lemma-lift-residue-field} のように $\kappa \to A$ を選ぶ．
$\dim_\kappa(\mathfrak m/\mathfrak m^2) = 2$ なので，基底へ写る
$x_0, y_0 \in \mathfrak m$ を選べる．連続な $\kappa$-代数写像
$$
\kappa[[x, y]] \longrightarrow A
$$
を得る．これは $x \mapsto x_0$ および $y \mapsto y_0$ で定められる．
$q$ を $ax_0^2 + bx_0y_0 + cy_0^2$ の
$\text{Sym}^2_\kappa(\mathfrak m/\mathfrak m^2)$ における類とする．
$Q(x, y) = ax^2 + bxy + cy^2$ と書き，これを二変数多項式とみなす．すると
$$
Q(x_0, y_0) = ax_0^2 + bx_0y_0 + cy_0^2 =
\sum\nolimits_{i + j = 3} a_{ij} x_0^iy_0^j
$$
となる $a_{ij}$ が存在し，これは $A$ の元である．$x_0$, $y_0$ の選択を変えることにより，
消失次数を増やせることを示したい．$x_1, y_1 \in \mathfrak m^2$ と仮定する．すると
$$
Q(x_0 + x_1, y_0 + y_1) = Q(x_0, y_0) +
(2ax_0 + by_0)x_1 + (bx_0 + 2cy_0)y_1 \bmod \mathfrak m^4
$$
$Q$ の非退化性は，$2ax_0 + by_0$ と $bx_0 + 2cy_0$ が
$\kappa$-基底として $\mathfrak m/\mathfrak m^2$ をなすことにほかならない．補題の前の議論を
参照せよ．したがって $x_1, y_1 \in \mathfrak m^2$ を，
$Q(x_0 + x_1, y_0 + y_1) \in \mathfrak m^4$ となるように選べる．
同様に帰納を続ければ，次を満たす $x_i, y_i \in \mathfrak m^{i + 1}$ が得られる：
$$
Q(x_0 + x_1 + \ldots + x_n, y_0 + y_1 + \ldots + y_n) \in \mathfrak m^{n + 3}
$$
$A$ は完備なので，$x_\infty = \sum x_i$ および $y_\infty = \sum y_i$ とおける．
$\kappa[[x, y]] \longrightarrow A$ を，$x$ を $x_\infty$ へ，$y$ を $y_\infty$ へ
送る写像とする．Algebra，補題 \ref{algebra-lemma-completion-generalities} により，これは全射
$\kappa[[x, y]]/(Q) \longrightarrow A$ を誘導する．補題
\ref{curves-lemma-reduced-quotient-regular-ring-dim-2} により，$k[[x, y]] \to A$ の核は
単項である．しかしこの核は $Q$ の真の約元を含みえない．そのような約元は
次数 $1$ で，$x, y$ に関するものであり，$\dim(\mathfrak m/\mathfrak m^2) = 2$ と矛盾するからである．
したがって所望のとおり $Q$ が核を生成する．
\end{proof}

\begin{lemma}
\label{curves-lemma-2-branches-delta-1}
$k$ を体とする．$(A, \mathfrak m, \kappa)$ を Nagata 局所 $k$-代数とする．
次の条件は同値である．
\begin{enumerate}
\item $k \to A$ は補題 \ref{curves-lemma-nodal-algebraic} のとおりである．
\item $\kappa/k$ は分離的で，$A$ は次元 $1$ の被約環であり，
$\delta$-不変量について $A$ の値は $1$ で，$A$ は $2$ 本の幾何学的枝をもつ．
\end{enumerate}
これらが成り立つならば，$A'$，すなわち $A$ の全商環における整閉包は，
$1$ 個または $2$ 個の極大イデアル $\mathfrak m'$ をもち，拡大
$\kappa(\mathfrak m')/k$ は分離的である．
\end{lemma}

\begin{proof}
いずれの場合も $A$ と $A^\wedge$ は被約である．場合 (2) では，被約な Nagata
局所環の完備化が被約であることによる
（More on Algebra，補題 \ref{more-algebra-lemma-completion-reduced}）．
いずれの場合も $A$ と $A^\wedge$ の次元は $1$ である
（More on Algebra，補題 \ref{more-algebra-lemma-completion-dimension}）．
Varieties，補題 \ref{varieties-lemma-delta-same-after-completion} および
More on Algebra，補題
\ref{more-algebra-lemma-one-dimensional-number-of-branches} により，$\delta$-不変量および
幾何学的枝の数について，$A$ と $A^\wedge$ の値は一致する．Varieties，補題
\ref{varieties-lemma-pre-delta-invariant} のように，$A'$ を $A$ の全商環における
整閉包とする．Varieties，補題
\ref{varieties-lemma-normalization-same-after-completion} により，
$A' \otimes_A A^\wedge$ は $A^\wedge$ に対して同じ役割を果たす．したがって
$A$ を $A^\wedge$ で置き換え，$A$ は完備であると仮定してよい．

\medskip\noindent
(1) を仮定する．$A$ が二本の幾何学的枝をもち，$\delta$-不変量が $1$ であることを
示せば十分である．$A = \kappa[[x, y]]/(ax^2 + bxy + cy^2)$ で，
$q = ax^2 + bxy + cy^2$ は非退化であると仮定してよい．二つの場合がある．

\medskip\noindent
場合 I：$q$ が $\kappa$ 上で分解する．この場合，座標を変えることにより
$q = xy$ と仮定してよい．すると
$$
A' = \kappa[[x, y]]/(x) \times \kappa[[x, y]]/(y)
$$
である．

\medskip\noindent
場合 II：$q$ が分解しない．この場合 $c \not = 0$ であり，非退化であるとは
$b^2 - 4ac \not = 0$ ということである．したがって
$\kappa' = \kappa[t]/(a + bt + ct^2)$ は分離的な次数 $2$ の拡大であり，その基礎体は $\kappa$ である．
このとき $t = y/x$ は $A$ 上整であり，
$$
A' = \kappa'[[x]]
$$
を得る．右辺では $y$ は $tx$ に写る．

\medskip\noindent
いずれの場合も，$\delta$-不変量が $1$ で幾何学的枝の数が $2$ であることを
直接確かめられる．これで (1) が (2) を含意することが分かる．さらに補題の
最後の主張も従う．

\medskip\noindent
(2) を仮定する．More on Algebra，補題
\ref{more-algebra-lemma-number-of-branches-1} により，$A'$ は二つの極大イデアルを
もつか，または $A'$ は一つの極大イデアルをもち，
$[\kappa(\mathfrak m') : \kappa]_s = 2$ である．

\medskip\noindent
場合 I：$A'$ は二つの極大イデアル $\mathfrak m'_1$, $\mathfrak m'_2$ をもち，
その剰余体は $\kappa_1$, $\kappa_2$ である．$\delta$-不変量は $A'/A$ の長さであり，
全射 $A'/A \to (\kappa_1 \times \kappa_2)/\kappa$ が存在するから，
$\kappa = \kappa_1 = \kappa_2$ である．$A$ は完備であり
（Algebra，補題 \ref{algebra-lemma-complete-henselian} により Hensel 的でもある），
$A'$ は $A$ 上有限なので，$A' = A_1 \times A_2$ である
（Algebra，補題 \ref{algebra-lemma-finite-over-henselian}）．
$A'$ は正規環だから，$A_1$ と $A_2$ は離散付値環である．したがって
$A_1$ と $A_2$ は $\kappa[[t]]$ と同型である（$k$-代数として）．これは
More on Algebra，補題 \ref{more-algebra-lemma-power-series-over-residue-field} による．
$\delta$-不変量は $1$ なので，$A$ は $A_1$ と $A_2$ のくさび和である
（Varieties，定義 \ref{varieties-definition-wedge}）．したがって容易に
$A \cong \kappa[[x, y]]/(xy)$ を得る．

\medskip\noindent
場合 II：$A'$ はただ一つの極大イデアル $\mathfrak m'$ をもち，その剰余体は $\kappa'$ であり，
$[\kappa' : \kappa]_s = 2$ である．場合 I とまったく同様に論じると，
$[\kappa' : \kappa] = 2$ であり，$\kappa'$ は $\kappa$ 上分離的である．
$A'$ は正規なので，$A'$ は $\kappa'[[t]]$ と同型である（上の参照を見よ）．
$A'/A$ の長さは $1$ だから，
$$
A = \{f \in \kappa'[[t]] \mid f(0) \in \kappa\}
$$
を得る．ここから簡単な計算により，$A$ は場合 (1) の形であることが分かる．
\end{proof}

\begin{lemma}
\label{curves-lemma-fitting-ideal-well-defined}
$k$ を体とする．$A = k[[x_1, \ldots, x_n]]$ とする．
$I = (f_1, \ldots, f_m) \subset A$ をイデアルとする．任意の $r \geq 0$ に対し，
$A/I$ において $r \times r$ 小行列式で生成されるイデアルを考える．ここで小行列式は
行列 $(\partial f_j/\partial x_i)$ のものである．このイデアルは $I$ の生成元の選択にも，
正則パラメータ系 $x_1, \ldots, x_n$ の選択にも依存せず，後者は $A$ のものである．
\end{lemma}

\begin{proof}
この補題の「正しい」証明は，このイデアルが第 $(n - r)$ Fitting イデアルであることを
示すものである．対象の連続微分加群は $A/I$ のもので，基礎体は $k$ である．
ここでは直接証明する．
$g_1, \ldots g_l$ が $I$ の別の生成元系ならば，$g_s = \sum a_{sj}f_j$ と書け，
行列の等式
$$
(\partial g_s/\partial x_i) = (a_{sj}) (\partial f_j/\partial x_i)
+ (\partial a_{sj}/\partial x_i f_j)
$$
を得る．最後の項は $A/I$ で零である．Cauchy--Binet の公式により，$g_s$ に対する
小行列式のイデアルは $f_j$ に対するものに含まれる．対称性から両者は等しい．
$y_1, \ldots, y_n \in \mathfrak m_A$ が別の正則パラメータ系ならば，行列
$(\partial y_j/\partial x_i)$ は可逆であり，微分の連鎖律を用いればよい．
細部は省略する．
\end{proof}

\begin{lemma}
\label{curves-lemma-fitting-ideal}
$k$ を体とする．$A = k[[x_1, \ldots, x_n]]$ とする．
$I = (f_1, \ldots, f_m) \subset \mathfrak m_A$ をイデアルとする．次の条件は同値である．
\begin{enumerate}
\item $k \to A/I$ は補題 \ref{curves-lemma-nodal-algebraic} のとおりである．
\item $A/I$ は被約であり，$(n - 1) \times (n - 1)$ 小行列式，すなわち行列
$(\partial f_j/\partial x_i)$ のものが $I + \mathfrak m_A$ を生成する．
\item $\text{depth}(A/I) = 1$ であり，$(n - 1) \times (n - 1)$ 小行列式，すなわち行列
$(\partial f_j/\partial x_i)$ のものが $I + \mathfrak m_A$ を生成する．
\end{enumerate}
\end{lemma}

\begin{proof}
補題 \ref{curves-lemma-fitting-ideal-well-defined} により，証明中に座標系と生成元の選択を
変えてよい．

\medskip\noindent
(1) が成り立つならば，座標を変えて $x_1, \ldots, x_{n - 2}$ が $A/I$ で零に写り，
$A/I = k[[x_{n - 1}, x_n]]/(a x_{n - 1}^2 + b x_{n - 1}x_n + c x_n^2)$
となるようにできる．ここで $a x_{n - 1}^2 + b x_{n - 1}x_n + c x_n^2$ は
ある非退化二次形式である．このとき直接計算により (2) と (3) の両方が成り立つ．

\medskip\noindent
$(n - 1) \times (n - 1)$ 小行列式，すなわち行列
$(\partial f_j/\partial x_i)$ のものが $I + \mathfrak m_A$ を生成すると仮定する．
ある $i$ と $j$ に対して偏導函数
$\partial f_j/\partial x_i$ が $A$ の単元であると仮定する．このときパラメータ系
$f_j, x_1, \ldots, x_{i - 1}, \hat x_i, x_{i + 1}, \ldots, x_n$
と生成元系
$f_j, f_1, \ldots, f_{j - 1}, \hat f_j, f_{j + 1}, \ldots, f_m$
を用いてよい．後者は $I$ の生成元系である．こうして正則パラメータ系
$x_1, \ldots, x_n$ と生成元系 $x_1, f_2, \ldots, f_m$ を得る．この生成元系は $I$ のものである．
次に，$i \geq 2$ と $j \geq 2$ で，$\partial f_j/\partial x_i$ が $A$ の単元となるものを
探す．そのような組があれば，上と同じ置換により，正則パラメータ系
$x_1, \ldots, x_n$ と生成元系 $x_1, x_2, f_3, \ldots, f_m$ があると仮定できる．
後者は $I$ の生成元系である．
これを続けると，有限回で次の状況に達する．正則パラメータ系
$x_1, \ldots, x_n$ と生成元系 $x_1, \ldots, x_t, f_{t + 1}, \ldots, f_m$ があり，
後者は $I$ の生成元系である．さらに $\partial f_j/\partial x_i \in \mathfrak m_A$ であり，
これはすべての $i, j \geq t + 1$ に対して成り立つ．

\medskip\noindent
この場合，偏導函数の行列は次のブロック形をもつ：
$$
\left(
\begin{matrix}
I_{t \times t} & * \\
0 & \mathfrak m_A
\end{matrix}
\right)
$$
したがって，どの $(n - 1) \times (n - 1)$ 小行列式も
$\mathfrak m_A^{n - 1 - t}$ に属する．$I \not = \mathfrak m_A$ であることに注意せよ．
さもなければ小行列式のイデアルは $1$ を含む．ゆえに
$n - 1 - t \leq 1$ である．実際，$\mathfrak m_A \setminus \mathfrak m_A^2 + I$ の元が
存在する（さもなければ Nakayama により $I = \mathfrak m_A$）．したがって
$t \geq n - 2$ である．上で $t \not = n$ を見た．同様に，$t = n - 1$ ならば
可逆な $(n - 1) \times (n - 1)$ 小行列式が存在し，これも許されない．したがって
$t = n - 2$ である．このとき $A/I$ は $k[[x_{n - 1}, x_n]]$ の商であり，補題
\ref{curves-lemma-reduced-quotient-regular-ring-dim-2} により，場合 (2)，(3) のいずれでも
$I$ は $x_1, \ldots, x_{n - 2}, f$ で生成される．ここで
$f = f(x_{n - 1}, x_n)$ である．この場合，小行列式に関する条件は，$f$ の二次項が
非退化であること，すなわち $A/I$ が補題 \ref{curves-lemma-nodal-algebraic} の形であることに
ほかならない．
\end{proof}

\begin{lemma}
\label{curves-lemma-nodal}
$k$ を体とする．$X$ を次元 $1$ の代数的 $k$-概形とする．
$x \in X$ を閉点とする．次の条件は同値である．
\begin{enumerate}
\item $x$ は節点である．
\item $k \to \mathcal{O}_{X, x}$ は補題 \ref{curves-lemma-nodal-algebraic} のとおりである．
\item 任意の $\overline{x} \in X_{\overline{k}}$ で $x$ へ写るものは節点特異点を定める．
\item $\kappa(x)/k$ は分離的で，$\mathcal{O}_{X, x}$ は被約であり，
$\Omega_{X/k}$ の第一 Fitting イデアルは $\mathfrak m_x$ を
$\mathcal{O}_{X, x}$ において生成する．
\item $\kappa(x)/k$ は分離的で，$\text{depth}(\mathcal{O}_{X, x}) = 1$ であり，
$\Omega_{X/k}$ の第一 Fitting イデアルは $\mathfrak m_x$ を
$\mathcal{O}_{X, x}$ において生成する．
\item $\kappa(x)/k$ は分離的で，$\mathcal{O}_{X, x}$ は被約であり，
$\delta$-不変量は $1$ で，幾何学的枝を $2$ 本もつ．
\end{enumerate}
\end{lemma}

\begin{proof}
まず $k$ は代数閉であると仮定する．この場合，(1) と (3) の同値性は自明である．
(1)，(3) と (2) の同値性は，二変数非退化二次形式が座標変換を除けば $xy$ だけで
あることから従う．(1) と (6) の同値性は補題 \ref{curves-lemma-multicross-algebra} である．
$X$ を $x$ のアフィン近傍で置き換えると，閉浸入
$X \to \mathbf{A}^n_k$ で $x$ を $0$ へ写すものが存在すると仮定してよい．
$f_1, \ldots, f_m \in k[x_1, \ldots, x_n]$ をイデアル $I$ の生成元とする．
ここでこのイデアルは $X$ の $\mathbf{A}^n_k$ における定義イデアルである．このとき $\Omega_{X/k}$ は
$R = k[x_1, \ldots, x_n]/I$-加群 $\Omega_{R/k}$ に対応し，後者は表示
$$
R^{\oplus m} \xrightarrow{(\partial f_j/\partial x_i)}
R^{\oplus n} \to \Omega_{R/k} \to 0
$$
をもつ（Algebra，第 \ref{algebra-section-differentials} 節および
第 \ref{algebra-section-netherlander} 節を参照）．したがって $\Omega_{R/k}$ の
第一 Fitting イデアルは，$(n - 1) \times (n - 1)$ 小行列式，すなわち行列
$(\partial f_j/\partial x_i)$ のものによって生成されるイデアルである．ゆえに，
$k[x_1, \ldots, x_n] \to R$ を極大イデアル $(x_1, \ldots, x_n)$ で完備化したものに
補題 \ref{curves-lemma-fitting-ideal} を適用すれば，(2)，(4)，(5) は同値である．

\medskip\noindent
次に $k$ を任意の体とする．場合 (2)，(4)，(5)，(6) では剰余体 $\kappa(x)$ は
$k$ 上分離的である．場合 (1)，(3) でも同様であることを示そう．
$Z \subset X$ を，$X$ 上で $\Omega_{X/k}$ の第一 Fitting イデアルにより定義される
閉部分概形とする．$Z$ の形成は体拡大と可換である
（Divisors，補題 \ref{divisors-lemma-base-change-and-fitting-ideal-omega}）．
(1) または (3) が成り立つならば，点 $\overline{x}$ で $X_{\overline{k}}$ に属し，
$\overline{x}$ が重複度 $1$ の $Z_{\overline{k}}$ の孤立点となるものが存在する
（$\overline{k}$ 上では補題の条件が同値だからである）．特に
$Z_{\overline{x}}$ は $\overline{x}$ において幾何学的に被約である
（$\overline{k}$ は代数閉である）．したがって $Z$ は $x$ において幾何学的に被約である
（Varieties，補題 \ref{varieties-lemma-geometrically-reduced-upstairs}）．特に $Z$ は
$x$ において被約であり，したがって $Z = \Spec(\kappa(x))$ が $x$ の近傍で成り立ち，
$\kappa(x)$ は $k$ 上幾何学的に被約である．
これは $\kappa(x)/k$ が分離的であることを意味する
（Algebra，補題 \ref{algebra-lemma-characterize-separable-field-extensions}）．

\medskip\noindent
前段落の議論から，(1) または (3) が成り立つならば，$\Omega_{X/k}$ の第一 Fitting
イデアルは $\mathfrak m_x$ を生成する．写像
$\mathcal{O}_{X, x} \to \mathcal{O}_{X_{\overline{k}}, \overline{x}}$ は平坦であり，
$\mathcal{O}_{X_{\overline{k}}, \overline{x}}$ は被約で深さ $1$ をもつから，(4) と (5) が
成り立つ（Algebra，補題 \ref{algebra-lemma-descent-reduced} および
\ref{algebra-lemma-apply-grothendieck} を用いる）．逆に，(4) は Algebra，補題
\ref{algebra-lemma-criterion-reduced} により (5) を含意する．(5) が成り立つならば，
$Z$ は $x$ において幾何学的に被約である（$\kappa(x)/k$ は分離的であり，$Z$ は
ある近傍で $x$ だからである）．したがって $Z_{\overline{k}}$ は，任意の点
$\overline{x}$ で被約である；ここでその点は $X_{\overline{k}}$ の点で $x$ の上にある．
言い換えると，
$\Omega_{X_{\overline{k}}/\overline{k}}$ の第一 Fitting イデアルは
$\mathfrak m_{\overline{x}}$ を $\mathcal{O}_{X_{\overline{k}, \overline{x}}}$ において生成する．
さらに，$\mathcal{O}_{X, x} \to \mathcal{O}_{X_{\overline{k}}, \overline{x}}$ は平坦なので，
$\text{depth}(\mathcal{O}_{X_{\overline{k}}, \overline{x}}) = 1$ である
（上の参照を見よ）．したがって $\overline{x} \in X_{\overline{k}}$ に対して (5) が成り立ち，
代数閉体上での同値性から (3) が従う．こうして (1)，(3)，(4)，(5) は同値である．

\medskip\noindent
(2) と (6) の同値性は補題 \ref{curves-lemma-2-branches-delta-1} から従う．

\medskip\noindent
最後に，(2) $=$ (6) と (1) $=$ (3) $=$ (4) $=$ (5) の同値性を示す．まず
$X$ の $x$ における幾何学的枝の数と $X_{\overline{k}}$ の $\overline{x}$ における
幾何学的枝の数は，Varieties，補題
\ref{varieties-lemma-geometric-branches-and-change-of-fields} により等しい．
ここまでに得た情報から，いずれの場合もこの数は $2$ である．
一方，場合 (1) では $X$ が $x$ において幾何学的に被約であることは明らかであり，したがって
$$
\delta\text{-不変量：}X\text{ の点 }x \leq
\delta\text{-不変量：}X_{\overline{k}}\text{ の点 }\overline{x}
$$
である．これは Varieties，補題 \ref{varieties-lemma-delta-invariant-and-change-of-fields}
による．場合 (1) では右辺が $1$ なので，$\delta$-不変量は $X$ の $x$ において
$1$ でなければならない（もし零ならば，Varieties，補題
\ref{varieties-lemma-delta-invariant-is-zero} により $\mathcal{O}_{X, x}$ は
離散付値環となるが，これは単枝であり矛盾する）．したがって (5) が成り立つ．
逆に (2) $=$ (5) が成り立つならば，Varieties，補題
\ref{varieties-lemma-geometrically-normal-in-codim-1} の仮定 (a)，(b)，(c) は，
$x \in X$ に対して補題 \ref{curves-lemma-2-branches-delta-1} により成り立つ．よって Varieties，補題
\ref{varieties-lemma-delta-invariant-and-change-of-fields-better} を適用でき，上の不等式で
等号が成り立つ．したがって $\overline{x} \in X_{\overline{k}}$ に対して (5) が成り立ち，
代数閉体上の条件の同値性により場合 (1) に戻る．
\end{proof}

\begin{remark}[節点に付随する二次拡大]
\label{curves-remark-quadratic-extension}
$k$ を体とする．$(A, \mathfrak m, \kappa)$ を Noether 局所 $k$-代数とする．
$(A, \mathfrak m, \kappa)$ が補題 \ref{curves-lemma-nodal-algebraic} のとおりであるか，
$A$ が補題 \ref{curves-lemma-2-branches-delta-1} のとおり Nagata であるか，または
$A$ が完備で補題 \ref{curves-lemma-fitting-ideal} のとおりであると仮定する．このとき $A$ は
次数 $2$ の分離的 $\kappa$-代数 $\kappa'$ を次のように標準的に定める．
\begin{enumerate}
\item $q = ax^2 + bxy + cy^2$ を補題 \ref{curves-lemma-nodal-algebraic} のような非退化二次形式とし，
座標 $x, y$ を $a \not = 0$ となるように選び，
$\kappa' = \kappa[x]/(ax^2 + bx + c)$ とおく．
\item $A' \supset A$ を $A$ の全商環における整閉包とし，
$\kappa' = A'/\mathfrak m A'$ とおく．または
\item $\kappa'$ を，次を満たす $\kappa$-代数とする：
$\text{Proj}(\bigoplus_{n \geq 0} \mathfrak m^n/\mathfrak m^{n + 1}) =
\Spec(\kappa')$.
\end{enumerate}
(1) と (2) の同値性は補題 \ref{curves-lemma-2-branches-delta-1} の証明で示した．
これらと (3) の同値性は省略する．$X$ が局所 Noether $k$-概形で，$x \in X$ が
$\mathcal{O}_{X, x} = A$ を満たす点ならば，(3) から
$\Spec(\kappa') = X^\nu \times_X \Spec(\kappa)$ を得る．ここで
$\nu : X^\nu \to X$ は正規化射である．
\end{remark}

\begin{remark}[自明な二次拡大]
\label{curves-remark-trivial-quadratic-extension}
$k$ を体とする．$(A, \mathfrak m, \kappa)$ は注意
\ref{curves-remark-quadratic-extension} のとおりとし，$\kappa'/\kappa$ を付随する
次数 $2$ の分離的代数とする．次の条件は同値である．
\begin{enumerate}
\item $\kappa' \cong \kappa \times \kappa$ は $\kappa$-代数としての同型である．
\item 形式 $q$ を補題 \ref{curves-lemma-nodal-algebraic} において $xy$ と選べる．
\item $A$ は二本の枝をもつ．
\item 拡大 $A'/A$ は補題 \ref{curves-lemma-2-branches-delta-1} のものであり，二つの極大イデアルをもつ．
\item $A^\wedge \cong \kappa[[x, y]]/(xy)$ は $k$-代数としての同型である．
\end{enumerate}
これらの条件の同値性は補題 \ref{curves-lemma-2-branches-delta-1} の証明で示した．
$X$ が局所 Noether $k$-概形で，$x \in X$ が $\mathcal{O}_{X, x} = A$ を満たす点ならば，
これは二点 $x_1, x_2$ が正規化 $X^\nu$ にあり，$x$ の上に位置し，
$\kappa(x) = \kappa(x_1) = \kappa(x_2)$ となることにほかならない．
\end{remark}

\begin{definition}
\label{curves-definition-split-node}
$k$ を体とする．$X$ を次元 $1$ の代数的 $k$-概形とする．
$x \in X$ を閉点とする．$x$ が {\it 分裂節点}であるとは，
$x$ が節点で，$\kappa(x) = k$ であり，注意
\ref{curves-remark-trivial-quadratic-extension} の同値な条件が
$A = \mathcal{O}_{X, x}$ に対して成り立つことをいう．
\end{definition}

\noindent
この概念についてすでに分かっていることを述べる，当然の補題を定式化する．

\begin{lemma}
\label{curves-lemma-split-node}
$k$ を体とする．$X$ を次元 $1$ の代数的 $k$-概形とする．
$x \in X$ を閉点とする．次の条件は同値である．
\begin{enumerate}
\item $x$ は分裂節点である．
\item $x$ は節点であり，ちょうど二点 $x_1, x_2$ が正規化 $X^\nu$ 上で
$x$ の上にあって，$k = \kappa(x_1) = \kappa(x_2)$ を満たす．
\item $\mathcal{O}_{X, x}^\wedge \cong k[[x, y]]/(xy)$ は $k$-代数としての同型である．
\item ここにさらに追加する．
\end{enumerate}
\end{lemma}

\begin{proof}
注意 \ref{curves-remark-trivial-quadratic-extension} の議論と
補題 \ref{curves-lemma-nodal} から従う．
\end{proof}

\begin{lemma}
\label{curves-lemma-node-field-extension}
$K/k$ を体の拡大とする．$X$ を局所代数的 $k$-概形で次元 $1$ のものとする．
$y \in X_K$ を像 $x \in X$ をもつ点とする．次の条件は同値である．
\begin{enumerate}
\item $x$ は $X$ の閉点かつ節点である．
\item $y$ は $Y$ の閉点かつ節点である．
\end{enumerate}
\end{lemma}

\begin{proof}
$x$ が $X$ の閉点ならば，$y$ も閉点である（剰余体を見よ）．
しかし逆は必ずしも成り立たない；すなわち，$Y$ の閉点が $X$ の非閉点へ
写ることがある．ただしこの場合，$y$ は節点ではありえない．実際，そのとき
$X$ は $x$ において幾何学的に単枝となる（$x$ は $X$ の生成点であり，
$\mathcal{O}_{X, x}$ は Artin 環となり，任意の Artin 局所環は幾何学的に単枝だからである）．
したがって $Y$ は $y$ において幾何学的に単枝である
（Varieties，補題
\ref{varieties-lemma-geometrically-unibranch-and-change-of-fields}),
これは補題 \ref{curves-lemma-nodal} により $y$ が節点ではありえないことを意味する．
したがって，$x$ と $y$ はともに閉点であると仮定してよいし，そう仮定する．

\medskip\noindent
代数閉包 $\overline{k}$，$\overline{K}$ と，写像 $\overline{k} \to \overline{K}$ で
与えられた写像 $k \to K$ を延長するものを選ぶ．補題 \ref{curves-lemma-nodal} における (1) と (3) の
同値性を用いれば，$k$ と $K$ が代数閉である場合に帰着する．この場合，補題
\ref{curves-lemma-nodal} の証明と同様に議論して，幾何学的枝の数と $\delta$-不変量は，
$X$ の $x$ におけるものと $Y$ の $y$ におけるものとで同じであることが分かる．
別の議論として，次の同型を選んでもよい：
$k[[x_1, \ldots, x_n]]/(g_1, \ldots, g_m) \to \mathcal{O}_{X, x}^\wedge$
これは Varieties，補題
\ref{varieties-lemma-complete-local-ring-structure-as-algebra}.
のような $k$-代数の同型である．Varieties，補題
\ref{varieties-lemma-base-change-complete-local-ring} により，これは同型
$K[[x_1, \ldots, x_n]]/(g_1, \ldots, g_m) \to \mathcal{O}_{Y, y}^\wedge$
を $K$-代数について与える．定義により，示すべきことは
$$
k[[x_1, \ldots, x_n]]/(g_1, \ldots, g_m) \cong k[[s, t]]/(st)
$$
と
$$
K[[x_1, \ldots, x_n]]/(g_1, \ldots, g_m) \cong K[[s, t]]/(st)
$$
が同値であることである．読者自身でこれを証明されたい．
$k$ と $K$ は代数閉体なので，これはこれらの環が補題
\ref{curves-lemma-nodal-algebraic} のとおりであることと同じである．補題
\ref{curves-lemma-fitting-ideal} を介すると，これは行列の $(n - 1) \times (n - 1)$ 小行列式
$(\partial g_j/\partial x_i)$ に関する主張へ翻訳され，使用する体に依存しないことは明らかである．
詳細は省略する．
\end{proof}

\begin{lemma}
\label{curves-lemma-node-etale-local}
$k$ を体とする．$X$ を局所代数的 $k$-概形で次元 $1$ のものとする．
$Y \to X$ を \'etale 射とする．$y \in Y$ を像 $x \in X$ をもつ点とする．
次の条件は同値である．
\begin{enumerate}
\item $x$ は $X$ の閉点かつ節点である．
\item $y$ は $Y$ の閉点かつ節点である．
\end{enumerate}
\end{lemma}

\begin{proof}
補題 \ref{curves-lemma-node-field-extension} により，$k$ の代数閉包へ基底変換してよい．
すると $x$ と $y$ の剰余体は $k$ である．ゆえに写像
$\mathcal{O}_{X, x}^\wedge \to \mathcal{O}_{Y, y}^\wedge$ は同型である
（例えば \'Etale Morphisms，補題 \ref{etale-lemma-characterize-etale-completions}，または
More on Algebra，補題 \ref{more-algebra-lemma-flat-unramified} による）．
したがって補題は明らかである．
\end{proof}

\begin{lemma}
\label{curves-lemma-node-over-separable-extension}
$k'/k$ を有限分離体拡大とする．$X$ を局所代数的 $k'$-概形で次元 $1$ のものとする．
$x \in X$ を閉点とする．次の条件は同値である．
\begin{enumerate}
\item $x$ は節点である．
\item $x$ は，$X$ を局所代数的 $k$-概形とみなしたときにも節点である．
\end{enumerate}
\end{lemma}

\begin{proof}
補題 \ref{curves-lemma-nodal} における節点の特徴づけから直ちに従う．
\end{proof}

\begin{lemma}
\label{curves-lemma-nodal-lci}
$k$ を体とする．$X$ を局所代数的 $k$-概形で，次元 $1$ の等次元なものとする．
次の条件は同値である．
\begin{enumerate}
\item $X$ の特異点は高々節点的である．
\item $X$ は $k$ 上の局所完全交差であり，閉部分概形 $Z \subset X$ で
$\Omega_{X/k}$ の第一 Fitting イデアルにより切り出されるものは $k$ 上非分岐である．
\end{enumerate}
\end{lemma}

\begin{proof}
この補題の証明は読者自身で見つけられたい．以下は先の結果を組み合わせるだけであり，
実際に何が起きているかを覆い隠すかもしれない．

\medskip\noindent
(2) を仮定する．$Z \to \Spec(k)$ は準有限なので
（Morphisms，補題 \ref{morphisms-lemma-unramified-quasi-finite}），点 $x \in Z$ の
剰余体は $k$ 上有限（かつ分離的）であることが Morphisms，補題
\ref{morphisms-lemma-residue-field-quasi-finite} から分かる．したがって各 $x \in Z$ は
$X$ の閉点である．これは Morphisms，補題
\ref{morphisms-lemma-algebraic-residue-field-extension-closed-point-fibre}.
による．局所環 $\mathcal{O}_{X, x}$ は Algebra，補題
\ref{algebra-lemma-lci-CM} により Cohen--Macaulay である．次元論により
$\dim(\mathcal{O}_{X, x}) = 1$ なので（Varieties，第
\ref{varieties-section-algebraic-schemes} 節），$\text{depth}(\mathcal{O}_{X, x})) = 1$
と結論する．したがって $x$ は補題 \ref{curves-lemma-nodal} により節点である．
$x \in X$ かつ $x \not \in Z$ ならば，$X \to \Spec(k)$ は $x$ において滑らかである．
これは Divisors，補題 \ref{divisors-lemma-d-fitting-ideal-omega-smooth} による．

\medskip\noindent
(1) を仮定する．この仮定のもとで $X$ はすべての閉点において幾何学的に被約である
（Varieties，補題 \ref{varieties-lemma-geometrically-reduced-upstairs} を参照）．したがって
$X \to \Spec(k)$ は，Varieties，補題
\ref{varieties-lemma-geometrically-reduced-dense-smooth-open} により稠密開集合上で滑らかである．
ゆえに $Z$ は閉であり，閉点のみからなる．Divisors，補題
\ref{divisors-lemma-d-fitting-ideal-omega-smooth} により射
$X \setminus Z \to \Spec(k)$ は滑らかである．したがって $X \setminus Z$ は Morphisms，補題
\ref{morphisms-lemma-smooth-syntomic} と Morphisms，定義
\ref{morphisms-definition-syntomic} における局所完全交差の定義により，局所完全交差である．
補題 \ref{curves-lemma-nodal} により，すべての点 $x \in Z$ について局所環
$\mathcal{O}_{Z, x}$ は $\kappa(x)$ に等しく，$\kappa(x)$ は $k$ 上分離的である．
したがって $Z \to \Spec(k)$ は非分岐である
（Morphisms，補題 \ref{morphisms-lemma-unramified-over-field}）．最後に，補題
\ref{curves-lemma-nodal} と補題 \ref{curves-lemma-nodal-algebraic} の (3) により，
$\mathcal{O}_{X, x}$ は Divided Power Algebra，定義
\ref{dpa-definition-lci} の意味で完全交差である．しかし Divided Power Algebra，補題
\ref{dpa-lemma-check-lci-agrees} と Morphisms，補題
\ref{morphisms-lemma-local-complete-intersection} により，これは体上の局所完全交差概形を
定義するのに用いた概念と一致する．これで証明が完了する．
\end{proof}

\begin{lemma}
\label{curves-lemma-facts-about-nodal-curves}
$k$ を体とする．$X$ を局所代数的 $k$-概形で，次元 $1$ の等次元なものとし，
その特異点は高々節点的であるとする．このとき $X$ は Gorenstein かつ幾何学的に被約である．
\end{lemma}

\begin{proof}
Gorenstein であるという主張は，補題 \ref{curves-lemma-nodal-lci} と Duality for Schemes，補題
\ref{duality-lemma-gorenstein-lci} から従う．別の方法として，代数閉包へ移った後に
確かめれば十分であること（Duality for Schemes，補題
\ref{duality-lemma-gorenstein-base-change}），Noether 局所環が Gorenstein であることと
その完備化がそうであることが同値であること（Dualizing Complexes，補題
\ref{dualizing-lemma-flat-under-gorenstein}）を用い，局所環 $k[[t]]$ と
$k[[x, y]]/(xy)$ が Gorenstein であることを直接証明してもよい．

\medskip\noindent
次に $X$ が幾何学的に被約であることを見るには，$X_{\overline{k}}$ が被約であることを
示せば十分である（Varieties，補題
\ref{varieties-lemma-perfect-reduced} および
\ref{varieties-lemma-geometrically-reduced}）．しかし $X_{\overline{k}}$ は代数閉体上の
節点曲線である．したがって $X_{\overline{k}}$ の完備局所環は
$\overline{k}[[t]]$ または $\overline{k}[[x, y]]/(xy)$ のいずれかと同型であり，
これらは所望のとおり被約である．
\end{proof}

\begin{lemma}
\label{curves-lemma-closed-subscheme-nodal-curve}
$k$ を体とする．$X$ を局所代数的 $k$-概形で，等次元かつ次元 $1$ であり，
その特異点は高々節点的であるものとする．$Y \subset X$ が被約閉部分概形で，
等次元かつ次元 $1$ ならば，次が成り立つ．
\begin{enumerate}
\item $Y$ の特異点は高々節点的である．
\item $Z \subset X$ が $X \setminus Y$ の概形論的閉包ならば，
\begin{enumerate}
\item 概形論的交叉 $Y \cap Z$ は $k$ の有限分離拡大のスペクトルの非交和である．
\item $Y \cap Z$ の各点は $X$ の節点である．
\item $Y \to \Spec(k)$ は $Y \cap Z$ の各点において滑らかである．
\end{enumerate}
\end{enumerate}
\end{lemma}

\begin{proof}
$X$ と $Y$ は被約で，等次元かつ次元 $1$ なので，$Y$ は $X$ の既約成分の
ある部分集合の概形論的和である（被約環では $(0)$ は極小素イデアルの交わりである）．
$y \in Y$ を閉点とする．$y$ が $X \to \Spec(k)$ の滑らかな軌跡にあれば，
$y$ は $X$ のただ一つの既約成分上にあり，$Y$ と $X$ は $y$ のある開近傍で一致する．
したがって $Y \to \Spec(k)$ は $y$ において滑らかである．
$y$ が $X$ の節点であるが，なお $X$ のただ一つの既約成分上にあるならば，
同じ議論により $y$ は $Y$ の節点である．$y$ が $1$ 個より多くの $X$ の既約成分上に
あると仮定する．$X$ の $y$ における幾何学的枝の数は，補題 \ref{curves-lemma-nodal} により
$2$ であるから，Properties，補題
\ref{properties-lemma-number-of-branches-irreducible-components}.
により高々 $2$ 個の既約成分が $y$ を通る．$Y$ がその双方を含むならば，
再び $Y = X$ が $y$ のある開近傍で成り立ち，$y$ は $Y$ の節点である．最後に，
$Y$ は既約成分の一つだけを含むと仮定する．$X$ を $x$ の開近傍で置き換えると，
$Y$ は二つの既約成分の一方で，$Z$ は他方であると仮定してよい．Properties，補題
\ref{properties-lemma-number-of-branches-irreducible-components}
により再び，$X$ は $y$ において二本の枝をもつことが分かる；すなわち局所環
$\mathcal{O}_{X, y}$ は二本の枝をもち，それらは $\mathcal{O}_{Y, y}$ と
$\mathcal{O}_{Z, y}$ から来る．次のように書く：
$\mathcal{O}_{X, y}^\wedge \cong \kappa(y)[[u, v]]/(uv)$
これは注意 \ref{curves-remark-trivial-quadratic-extension} のとおりである．例えば体
$\kappa(y)$ は $k$ 上有限分離的であることが補題 \ref{curves-lemma-nodal} から従う．
したがって，必要なら $u$ と $v$ の役割を交換すると，写像
$\mathcal{O}_{X, y} \to \mathcal{O}_{Y, Y}$
の完備化は $\kappa(y)[[u, v]]/(uv) \to \kappa(y)[[u]]$ に対応し，写像
$\mathcal{O}_{X, y} \to \mathcal{O}_{Y, Y}$
の完備化は $\kappa(y)[[u, v]]/(uv) \to \kappa(y)[[v]]$ に対応する．
$Y \cap Z$ の概形論的交叉は両者のイデアルの和で切り出され，その和は完備化において
$(u, v)$，すなわち極大イデアルである．したがって (2)(a) と (2)(b) は明らかである．
最後に (2)(c) が成り立つ：$\mathcal{O}_{Y, y}$ の完備化は正則であり，したがって
$\mathcal{O}_{Y, y}$ は正則である（More on Algebra，補題
\ref{more-algebra-lemma-completion-regular}）．また $\kappa(y)/k$ は分離的なので，
Algebra，補題 \ref{algebra-lemma-separable-smooth} により，ある開近傍で滑らかである．
\end{proof}




\section{節点曲線族}
\label{curves-section-families-nodal}

\noindent
Stacks プロジェクトでは曲線は次元 $1$ の既約多様体であるが，文献における
「半安定曲線」や「節点曲線」は通常既約ではなく，しばしば固有であると仮定される．
特に「半安定曲線の族」，「節点曲線の族」，または「節点族」という句で用いる場合が
そうである．したがって，文献と整合し，かつ内部でも整合する用語を選ぶのは少し難しい．
さらに，まず本当に調べたいのは次の補題で導入する概念である（これは始域上局所的である）．

\begin{lemma}
\label{curves-lemma-nodal-family}
$f : X \to S$ を概形の射とする．次の条件は同値である．
\begin{enumerate}
\item $f$ は平坦かつ局所有限表示で，すべての空でないファイバー $X_s$ は
等次元かつ次元 $1$ であり，$X_s$ の特異点は高々節点的である．
\item $f$ は相対次元 $1$ の syntomic 射であり，$\text{Sing}(f) \subset X$ という
$\Omega_{X/S}$ の第一 Fitting イデアルで定義される閉部分概形は $S$ 上非分岐である．
\end{enumerate}
\end{lemma}

\begin{proof}
形成 $\text{Sing}(f)$ は基底変換と可換であることを思い出そう；Divisors，補題
\ref{divisors-lemma-base-change-and-fitting-ideal-omega}.
を参照せよ．したがって補題は，補題 \ref{curves-lemma-nodal-lci}，Morphisms，補題
\ref{morphisms-lemma-syntomic-flat-fibres}，および Morphisms，補題
\ref{morphisms-lemma-unramified-etale-fibres} から従う．
（自明な Morphisms，補題 \ref{morphisms-lemma-syntomic-locally-finite-presentation}
および \ref{morphisms-lemma-syntomic-flat} も用いる．）
\end{proof}

\begin{definition}
\label{curves-definition-nodal-family}
$f : X \to S$ を概形の射とする．$f$ が
{\it 相対次元 $1$ で高々節点的}であるとは，$f$ が補題
\ref{curves-lemma-nodal-family} の同値な条件を満たすことをいう．
\end{definition}

\noindent
この回りくどい用語を用いる理由はいくつかある\footnote{よりよい提案があれば，
Stacks プロジェクトの保守者に電子メールを送ってほしい．}．第一に，この概念を
文献の他の概念と混同しないようにしたい（本節の導入を参照）．第二に，この概念には
相対次元のより高い射への一般化がいくつか考えられる（例えば，\'etale 局所的に
高々節点的な射の合成である射を考えたり，ファイバーは高次元だが特異点が高々通常
二重点である射を考えたりできる）．

\begin{lemma}
\label{curves-lemma-smooth-relative-dimension-1}
相対次元 $1$ の滑らかな射は，相対次元 $1$ で高々節点的である．
\end{lemma}

\begin{proof}
省略する．
\end{proof}

\begin{lemma}
\label{curves-lemma-base-change-nodal-family}
$f : X \to S$ を相対次元 $1$ で高々節点的とする．このとき，$f$ の任意の
基底変換についても同じことが成り立つ．
\end{lemma}

\begin{proof}
これは次の事実から従う：syntomic 射の基底変換は syntomic である
（Morphisms，補題 \ref{morphisms-lemma-base-change-syntomic}）；相対次元 $1$ の射の
基底変換は相対次元 $1$ である
（Morphisms，補題 \ref{morphisms-lemma-base-change-relative-dimension-d}）；形成
$\text{Sing}(f)$ は基底変換と可換である（Divisors，補題
\ref{divisors-lemma-base-change-and-fitting-ideal-omega}）；
）；非分岐射の基底変換は非分岐である
（Morphisms，補題 \ref{morphisms-lemma-base-change-unramified}）．
\end{proof}

\noindent
次の補題によれば，射が相対次元 $1$ で高々節点的かどうかはファイバー上で確認できる．

\begin{lemma}
\label{curves-lemma-locus-where-nodal}
$f : X \to S$ を平坦かつ局所有限表示な概形の射とする．このとき最大の開部分概形
$U \subset X$ で，$f|_U : U \to S$ が相対次元 $1$ で高々節点的となるものが存在する．
さらに，$U$ の形成は任意の基底変換と可換である．
\end{lemma}

\begin{proof}
Morphisms，補題 \ref{morphisms-lemma-set-points-where-fibres-lci} により，$f$ が
syntomic となるこのような開集合が存在する．したがって $f$ は syntomic 射であると
仮定してよい．特に $f$ は Cohen--Macaulay 射である
（Duality for Schemes，補題 \ref{duality-lemma-lci-gorenstein} および
\ref{duality-lemma-gorenstein-CM-morphism}）．したがって $X$ は，$f$ が所与の相対次元を
もつような開かつ閉な部分概形の非交和である；Morphisms，補題
\ref{morphisms-lemma-flat-finite-presentation-CM-fibres-relative-dimension}.
を参照せよ．この分解は任意の基底変換で保たれる；Morphisms，補題
\ref{morphisms-lemma-base-change-relative-dimension-d} を参照せよ．一つ以外の成分を
捨てることにより，$f$ は相対次元 $1$ の syntomic 射であると仮定してよい．
$\text{Sing}(f) \subset X$ を，$\Omega_{X/S}$ の第一 Fitting イデアルで定義される
閉部分概形とする．最大の開部分概形 $W \subset \text{Sing}(f)$ で，$W \to S$ が
非分岐となり，その形成が基底変換と可換となるものが存在する
（Morphisms，補題 \ref{morphisms-lemma-set-points-where-fibres-unramified}）．また形成
$\text{Sing}(f)$ は基底変換と可換なので（Divisors，補題
\ref{divisors-lemma-base-change-and-fitting-ideal-omega}),
），
$$
U = (X \setminus \text{Sing}(f)) \cup W
$$
は $X$ の最大の開部分概形であって，$f|_U : U \to S$ が相対次元 $1$ で高々節点的となり，
$U$ の形成が基底変換と可換となるものであることが分かる．
\end{proof}

\begin{lemma}
\label{curves-lemma-nodal-family-precompose-etale}
$f : X \to S$ を相対次元 $1$ で高々節点的とする．$Y \to X$ が \'etale 射ならば，
合成 $g : Y \to S$ は相対次元 $1$ で高々節点的である．
\end{lemma}

\begin{proof}
$g$ は平坦かつ局所有限表示な射の合成として，平坦かつ局所有限表示であることに注意する
（Morphisms，補題 \ref{morphisms-lemma-etale-locally-finite-presentation}，
\ref{morphisms-lemma-etale-flat},
\ref{morphisms-lemma-composition-finite-presentation}，および
\ref{morphisms-lemma-composition-flat} を用いる）．したがって，ファイバーの特異点が
高々節点的であることを示せば十分である．これは補題 \ref{curves-lemma-node-etale-local} から従う
（また，\'etale 射と滑らかな射の合成が滑らかであることには，Morphisms，補題
\ref{morphisms-lemma-etale-smooth-unramified} および
\ref{morphisms-lemma-composition-smooth} を用いる）．
\end{proof}

\begin{lemma}
\label{curves-lemma-nodal-family-postcompose-etale}
$S' \to S$ を概形の \'etale 射とする．$f : X \to S'$ を相対次元 $1$ で
高々節点的とする．このとき合成 $g : X \to S$ は相対次元 $1$ で高々節点的である．
\end{lemma}

\begin{proof}
$g$ は平坦かつ局所有限表示な射の合成として，平坦かつ局所有限表示であることに注意する
（Morphisms，補題 \ref{morphisms-lemma-etale-locally-finite-presentation}，
\ref{morphisms-lemma-etale-flat},
\ref{morphisms-lemma-composition-finite-presentation}，および
\ref{morphisms-lemma-composition-flat} を用いる）．したがって $g$ のファイバーの特異点が
高々節点的であることを示せば十分である．これは補題
\ref{curves-lemma-node-over-separable-extension} と，滑らかな点についての同様の結果から従う．
\end{proof}

\begin{lemma}
\label{curves-lemma-nodal-family-etale-local-source}
$f : X \to S$ を概形の射とし，$\{U_i \to X\}$ を \'etale 被覆とする．
次の条件は同値である．
\begin{enumerate}
\item $f$ は相対次元 $1$ で高々節点的である．
\item 各 $U_i \to S$ は相対次元 $1$ で高々節点的である．
\end{enumerate}
言い換えると，相対次元 $1$ で高々節点的であるという性質は始域上 \'etale 局所的である．
\end{lemma}

\begin{proof}
一方の向きは補題 \ref{curves-lemma-nodal-family-precompose-etale} で見た．他方の向きについて，
局所有限表示，平坦，または相対次元 $1$ であるという性質は始域上 \'etale 局所的である
ことに注意する（Descent，補題
\ref{descent-lemma-locally-finite-presentation-fppf-local-source},
\ref{descent-lemma-flat-fpqc-local-source}，および
\ref{descent-lemma-dimension-at-point}）．ファイバーを取れば，$S$ が体のスペクトルである
場合に帰着する．この場合，結果は補題 \ref{curves-lemma-node-etale-local} から従う
（滑らかであることが始域上 \'etale 局所的であることには Descent，補題
\ref{descent-lemma-smooth-smooth-local-source} を用いる）．
\end{proof}

\begin{lemma}
\label{curves-lemma-nodal-family-fpqc-local-target}
$f : X \to S$ を概形の射とし，$\{U_i \to S\}$ を fpqc 被覆とする．
次の条件は同値である．
\begin{enumerate}
\item $f$ は相対次元 $1$ で高々節点的である．
\item 各 $X \times_S U_i \to U_i$ は相対次元 $1$ で高々節点的である．
\end{enumerate}
言い換えると，相対次元 $1$ で高々節点的であるという性質は終域上 fpqc 局所的である．
\end{lemma}

\begin{proof}
一方の向きは補題 \ref{curves-lemma-base-change-nodal-family} で見た．他方の向きについて，
局所有限表示，平坦，または相対次元 $1$ であるという性質は終域上 fpqc 局所的である
ことに注意する（Descent，補題
\ref{descent-lemma-descending-property-locally-finite-presentation},
\ref{descent-lemma-descending-property-flat}，および Morphisms，補題
\ref{morphisms-lemma-dimension-fibre-after-base-change}）．ファイバーを取れば，$S$ が
体のスペクトルである場合に帰着する．この場合，結果は補題
\ref{curves-lemma-node-field-extension} から従う（滑らかであることが終域上 fpqc 局所的である
ことには Descent，補題 \ref{descent-lemma-descending-property-smooth} を用いる）．
\end{proof}

\begin{lemma}
\label{curves-lemma-descend-nodal-family}
$S = \lim S_i$ を，推移射がアフィンである概形の有向系の極限とする．
$0 \in I$ とし，$f_0 : X_0 \to Y_0$ を $S_0$ 上の概形の射とする．
$S_0$，$X_0$，$Y_0$ は準コンパクトかつ準分離的であると仮定する．
$f_i : X_i \to Y_i$ を $f_0$ の $S_i$ への基底変換とし，$f : X \to Y$ を
$f_0$ の $S$ への基底変換とする．もし
\begin{enumerate}
\item $f$ は相対次元 $1$ で高々節点的であり，
\item $f_0$ は局所有限表示である，
\end{enumerate}
ならば，$i \geq 0$ で $f_i$ が相対次元 $1$ で高々節点的となるものが存在する．
\end{lemma}

\begin{proof}
Limits，補題 \ref{limits-lemma-descend-syntomic} により，ある $i$ が存在し，その $f_i$ は
syntomic である．すると $X_i = \coprod_{d \geq 0} X_{i, d}$ は開かつ閉な部分概形の
非交和で，各 $X_{i, d} \to Y_i$ は相対次元 $d$ をもつ；Morphisms，補題
\ref{morphisms-lemma-syntomic-relative-dimension} を参照せよ．Morphisms，補題
\ref{morphisms-lemma-dimension-fibre-after-base-change} が与える基底変換によるファイバー次元の
挙動から，$X \to X_i$ の像は $X_{i, 1}$ に入ることが分かる．すると $i' \geq i$ で，
$X_{i'} \to X_i$ の像が $X_{i, 1}$ に入るものが存在する；Limits，補題
\ref{limits-lemma-limit-contained-in-constructible} を参照せよ．したがって
$f_{i'} : X_{i'} \to Y_{i'}$ は相対次元 $1$ の syntomic 射である
（再び Morphisms，補題 \ref{morphisms-lemma-dimension-fibre-after-base-change} による）．
射 $\text{Sing}(f_{i'}) \to Y_{i'}$ を考える．その $Y$ への基底変換は非分岐射である．
したがって Limits，補題 \ref{limits-lemma-descend-unramified} により，$i'$ を大きくすれば
射 $\text{Sing}(f_{i'}) \to Y_{i'}$ は非分岐となる．これで証明が完了する．
\end{proof}

\begin{lemma}
\label{curves-lemma-formal-local-structure-nodal-family}
$f : T \to S$ を概形の射とする．$t \in T$ を像 $s \in S$ をもつ点とする．
次を仮定する．
\begin{enumerate}
\item $f$ は $t$ において平坦である．
\item $\mathcal{O}_{S, s}$ は Noether 環である．
\item $f$ は局所有限型である．
\item $t$ はファイバー $T_s$ の分裂節点である．
\end{enumerate}
このとき，$h \in \mathfrak m_s^\wedge$ と同型
$$
\mathcal{O}_{T, t}^\wedge \cong
\mathcal{O}_{S, s}^\wedge[[x, y]]/(xy - h)
$$
で，$\mathcal{O}_{S, s}^\wedge$-代数のものが存在する．
\end{lemma}

\begin{proof}
$S$ を $\Spec(\mathcal{O}_{S, s})$ で，$T$ を $\Spec(\mathcal{O}_{S, s})$ への
基底変換で置き換える．すると $T$ は局所 Noether であり，したがって
$\mathcal{O}_{T, t}$ は Noether 環である．$A = \mathcal{O}_{S, s}^\wedge$，
$\mathfrak m = \mathfrak m_A$，$B = \mathcal{O}_{T, t}^\wedge$ とおく．
More on Algebra，補題 \ref{more-algebra-lemma-flat-completion} により $A \to B$ は平坦である．
さらに
$\mathcal{O}_{T, t}/\mathfrak m_s \mathcal{O}_{T, t} = \mathcal{O}_{T_s, t}$
なので，$B/\mathfrak m B = \mathcal{O}_{T_s, t}^\wedge$ である．仮定 (4) と補題
\ref{curves-lemma-split-node} により，$\overline{u}, \overline{v} \in B/\mathfrak m B$ が存在して，写像
$$
(A/\mathfrak m)[[x, y]] \longrightarrow B/\mathfrak m B,\quad
x \longmapsto \overline{u},
x \longmapsto \overline{v}
$$
は全射で，核は $(xy)$ である．

\medskip\noindent
$n \geq 1$ と $u, v \in B$ があり，これらが $\overline{u}, \overline{v}$ へ写り，
$$
u v = h + \delta
$$
が，ある $h \in A$ と $\delta \in \mathfrak m^nB$ に対して成り立つと仮定する．
$u', v' \in B$ で，$u - u', v - v' \in \mathfrak m^n B$ を満たし，さらに
$$
u' v' = h' + \delta'
$$
がある $h' \in A$ と $\delta' \in \mathfrak m^{n + 1}B$ に対して成り立つものが存在すると
主張する．これを見るため，$\delta = \sum f_i b_i$ と書く．ここで
$f_i \in \mathfrak m^n$ かつ $b_i \in B$ である．次に
$b_i = a_i + u b_{i, 1} + v b_{i, 2} + \delta_i$ と書く．ここで
$a_i \in A$，$b_{i, 1}, b_{i, 2} \in B$，および $\delta_i \in \mathfrak m B$ と書く．
これは $B$ の剰余体が $A$ の剰余体と一致し，$u$ と $v$ の $B/\mathfrak m B$ における像が
極大イデアルを生成するため可能である．そこで
$$
u' = u - \sum b_{i, 2}f_i,\quad
v' = v - \sum b_{i, 1}f_i
$$
とおくと，
$$
u'v' = h + \delta - \sum (b_{i, 1}u + b_{i, 2}v)f_i + \sum
c_{ij}f_if_j =
h + \sum a_if_i + \sum f_i \delta_i + \sum c_{ij}f_if_j
$$
を，ある $c_{i, j} \in B$ に対して得る．したがって上の形の式が得られ，
$h' = h + \sum a_if_i$ かつ
$\delta' = \sum f_i \delta_i + \sum c_{ij}f_if_j$ である．

\medskip\noindent
帰納法で議論し，$u_1, v_1 \in B$ という $\overline{u}, \overline{v}$ の任意の持ち上げから
始めると，前段落の結果により元の列 $u_n, v_n \in B$ と $h_n \in A$ が得られ，
$u_n - u_{n + 1} \in \mathfrak m^n B$,
$v_n - v_{n + 1} \in \mathfrak m^n B$,
$h_n - h_{n + 1} \in \mathfrak m^n$,
かつ $u_n v_n - h_n \in \mathfrak m^n B$ を満たす．$A$ と $B$ は完備なので，
$u_\infty = \lim u_n$，$v_\infty = \lim v_n$，および
$h_\infty = \lim h_n$ とおけ，すると $u_\infty v_\infty = h_\infty$ を $B$ において得る．
したがって $A$-代数写像
$$
A[[x, y]]/(xy - h_\infty) \longrightarrow B
$$
で，$x$ を $u_\infty$ へ，$v$ を $v_\infty$ へ送るものを得る．これは平坦な
$A$-代数の間の写像であり，$\mathfrak m$ で割ると同型である．$\mathfrak m$ を法として
全射なので，完備性と Algebra，補題 \ref{algebra-lemma-completion-generalities} により全射である．
そこで Algebra，補題 \ref{algebra-lemma-mod-injective} を適用すれば，同型であると結論できる．
\end{proof}

\begin{lemma}
\label{curves-lemma-genus-in-nodal-family-of-curves}
$f : X \to S$ を概形の射とする．次を仮定する．
\begin{enumerate}
\item $f$ は固有である．
\item $f$ は相対次元 $1$ で高々節点的である．
\item $f$ の幾何学的ファイバーは連結である．
\end{enumerate}
このとき，(a) $f_*\mathcal{O}_X = \mathcal{O}_S$ であり，これは任意の基底変換後にも
成り立つ；(b) $R^1f_*\mathcal{O}_X$ は有限局所自由な $\mathcal{O}_S$-加群で，その形成は
任意の基底変換と可換である；(c) $R^qf_*\mathcal{O}_X = 0$ が $q \geq 2$ に対して成り立つ．
\end{lemma}

\begin{proof}
(a) は Derived Categories of Schemes，補題
\ref{perfect-lemma-proper-flat-geom-red-connected}.
から従う．Derived Categories of Schemes，補題
\ref{perfect-lemma-proper-flat-h0} により，$S$ 上局所的に
$Rf_*\mathcal{O}_X = \mathcal{O}_S \oplus P$ と書ける．ここで $P$ は完全で，
tor 振幅は $[1, \infty)$ に含まれる．形成 $Rf_*\mathcal{O}_X$ は任意の基底変換と
可換であることを思い出そう（Derived Categories of Schemes，補題
\ref{perfect-lemma-flat-proper-perfect-direct-image-general}).
したがって $s \in S$ に対して
$$
H^i(P \otimes_{\mathcal{O}_S}^\mathbf{L} \kappa(s)) =
H^i(X_s, \mathcal{O}_{X_s})
\text{ ただし }i \geq 1
$$
を得る．これは $i = 1$ でない限り零である．なぜなら $X_s$ は次元 $1$ の Noether 概形だからである；
Cohomology，命題 \ref{cohomology-proposition-vanishing-Noetherian} を参照せよ．
すると $P = H^1(P)[-1]$ であり，$H^1(P)$ は有限局所自由である．これは例えば
More on Algebra，補題 \ref{more-algebra-lemma-lift-perfect-from-residue-field} による．
すべてが基底変換と両立するので結論を得る．
\end{proof}





\section{節点族の \'Etale 局所構造}
\label{curves-section-etale-local-nodal}


\noindent
概形の射
$$
\Spec(\mathbf{Z}[u, v, a]/(uv - a))
\longrightarrow
\Spec(\mathbf{Z}[a])
$$
を考える．次の補題は，この射が節点曲線族の
\'etale 局所構造の模型であることを示す．

\begin{lemma}
\label{curves-lemma-etale-local-structure-nodal-family}
$f : X \to S$ を概形の射とする．$f$ は相対次元 $1$ で
高々節点的であると仮定する．$x \in X$ を，ファイバー $X_s$ の
特異点である点とする．このとき概形の可換図式
$$
\xymatrix{
X \ar[d] &
U \ar[rr] \ar[l] \ar[rd] & &
W \ar[r] \ar[ld] &
\Spec(\mathbf{Z}[u, v, a]/(uv - a)) \ar[d] \\
S & &
V \ar[ll] \ar[rr] & & \Spec(\mathbf{Z}[a])
}
$$
が存在する．ここで $X \leftarrow U$，$S \leftarrow V$，および
$U \to W$ は \'etale 射であり，右側の正方形はファイバー積をなし，
さらに点 $u \in U$ が存在して，その像は点 $x$（$X$ の点）である．
\end{lemma}

\begin{proof}[Artin 近似を用いる第一の証明]
まず絶対 Noether 近似を用いて，$\mathbf{Z}$ 上有限型の概形の場合に帰着する．
問題は $X$ および $S$ 上局所的である．したがって $X$ と $S$ は
アフィンであると仮定してよい．このとき $S = \Spec(R)$ と書け，
$R$ をフィルター余極限 $R = \colim R_i$ と書ける．ここで各環は
有限型 $\mathbf{Z}$-代数である．
Limits，補題 \ref{limits-lemma-descend-finite-presentation} を用いると，
ある $i$ と射 $f_i : X_i \to \Spec(R_i)$ で，$S$ への基底変換が $f$ となるものを
見いだせる．$i$ を大きくした後，$f_i$ は相対次元 $1$ で高々節点的であると
仮定してよい；補題 \ref{curves-lemma-descend-nodal-family} を参照せよ．
$x_i \in X_i$ を $x$ の像とすると，これはそのファイバーの特異点である．
例えば形成 $\text{Sing}(f)$ が基底変換と可換だからである
（Divisors，補題 \ref{divisors-lemma-base-change-and-fitting-ideal-omega}）．
$f_i : X_i \to S_i$ と $x_i$ に対して補題を証明できれば，
$f : X \to S$ に対する補題は基底変換によって従う．
したがって次段落で扱う場合に帰着する．

\medskip\noindent
$S$ は $\mathbf{Z}$ 上有限型であると仮定する．$s \in S$ を $x$ の像とする．
$\kappa(x)$ は $\kappa(s)$ の有限分離拡大であることを思い出そう．
例えば $\text{Sing}(f) \to S$ が非分岐であることから，または
$x$ がファイバー $X_s$ の節点であって補題 \ref{curves-lemma-nodal} を
適用できることから従う．さらに，$\kappa'/\kappa(x)$ を，次数 $2$ の分離代数で，
備考 \ref{curves-remark-quadratic-extension} において $\mathcal{O}_{X_s, x}$ に
付随させたものとする．More on Morphisms，補題
\ref{more-morphisms-lemma-realize-prescribed-residue-field-extension-etale}
により，\'etale 近傍 $(V, v) \to (S, s)$ を，拡大
$\kappa(v)/\kappa(s)$ が次を実現するように選べる：
拡大 $\kappa(x)/\kappa(s)$ を，$\kappa' \cong \kappa(x) \times \kappa(x)$ の
場合に実現し，拡大 $\kappa'/\kappa(s)$ を，$\kappa'$ が体の場合に実現する．
$X$ を $X \times_S V$ で，$S$ を $V$ で置き換えた後，
次段落で述べる状況に帰着する．

\medskip\noindent
$S$ は $\mathbf{Z}$ 上有限型で，$x \in X_s$ は分裂節点であると仮定する；
定義 \ref{curves-definition-split-node} を参照せよ．補題
\ref{curves-lemma-formal-local-structure-nodal-family} により，
$\mathcal{O}_{S, s}$-代数同型
$$
\mathcal{O}_{X, x}^\wedge \cong
\mathcal{O}_{S, s}^\wedge[[s, t]]/(st - h)
$$
が，ある $h \in \mathfrak m_s^\wedge \subset \mathcal{O}_{S, s}^\wedge$ に対して
存在する．言い換えると，準同型
$$
\sigma : \mathbf{Z}[a] \longrightarrow \mathcal{O}_{S, s}^\wedge
$$
で $a$ を $h$ に送るものを考えると，$\mathcal{O}_{S, s}$-代数同型
$$
\mathcal{O}_{X, x}^\wedge
\longrightarrow
\mathcal{O}_{Y_\sigma, y_\sigma}^\wedge
$$
が存在する．ここで
$$
Y_\sigma = \Spec(\mathbf{Z}[u, v, t]/(uv - a))
\times_{\Spec(\mathbf{Z}[a]), \sigma} \Spec(\mathcal{O}_{S, s}^\wedge)
$$
であり，$y_\sigma$ は $Y_\sigma$ の点で，
$\Spec(\mathcal{O}_{S, s}^\wedge)$ の閉点の上にあり，座標 $u, v$ が
零となるものである．$\mathcal{O}_{S, s}$ は More on Algebra，命題
\ref{more-algebra-proposition-ubiquity-G-ring} により G 環であるから，
More on Morphisms，補題
\ref{more-morphisms-lemma-relative-map-approximation-pre}
を適用して結論を得る．
\end{proof}

\begin{proof}[Artin 近似を用いない証明]
この証明は概略だけを述べる．これは Mohan Swaminathan によって提供され，
\cite[命題 X.2.1]{ACG} の議論を模倣する．前の証明とは異なり，
この証明は実際の方程式を用いる．

\medskip\noindent
第一の証明とまったく同様に，$S$ が $\mathbf{Z}$ 上有限型であり，
$x \in X_s$ が分裂節点である場合に帰着する；定義
\ref{curves-definition-split-node} を参照せよ．補題 \ref{curves-lemma-split-node} により，
$\kappa(s)$-代数同型
$$
\mathcal{O}_{X_s, x}^\wedge \cong \kappa(s)[[u, v]]/(uv)
$$
が存在する．$X_s$ の $x$ における接空間は次元 $2$ であることに注意する．
したがって Varieties，補題 \ref{varieties-lemma-immersion-into-affine} により，
Zariski 縮小の後，閉埋め込み $X \to Y$ は $S$ 上の概形の間のものであり，
$Y \to S$ は相対次元 $2$ の滑らかな射であると仮定してよい．
$X$ は（定義により）相対次元 $1$ の syntomic 射であるから，
$X$ は $Y$ の有効 Cartier 因子である
（Divisors，補題
\ref{divisors-lemma-immersion-lci-into-smooth-regular-immersion} から従う）．
したがって $Y$ と $X$ を縮小した後，$Y = \Spec(B)$ はアフィンであり，
ある非零因子 $g \in B$ が存在して $X = \Spec(B/gB)$ となると仮定してよい．

\medskip\noindent
$S = \Spec(R)$ と書き，$\mathfrak p \subset R$ を $s$ に対応する素イデアルとする．
$\mathfrak q \subset B$ を，$x$（$Y$ の点と見る）に対応する素イデアルとする．
写像
$$
B_\mathfrak q / \mathfrak p B_\mathfrak q
\longrightarrow
B_\mathfrak q / g B_\mathfrak q + \mathfrak p B_\mathfrak q
\longrightarrow
\kappa(s)[[u, v]]/(uv)
$$
があり，第二の矢印は完備化後に同型となる．
$B$ を適当な主局所化で置き換えた後，元 $U, V \in B$ が存在し，
$u$ および $v$ に $(u, v)^2$ を法として写ると仮定してよい．
このとき $B_\mathfrak q/\mathfrak p B_\mathfrak q$ の完備化は
$\kappa(s)[[U, V]]$ である．さらに $Y$ を縮小し，$g$ に単元を掛けた後，
$g$ は $UV$ と高次の項との和に，$\kappa(s)[[U, V]]$ において写ると
仮定してよい．$Y$ をさらに縮小して，次を仮定してよい．
\begin{enumerate}
\item $\mathfrak q = \mathfrak pB + (U, V)$
\item $\Omega_{B/R}$ は $\text{d}U$ および $\text{d}V$ を基底とする自由加群である
\end{enumerate}
$\text{d}(g) = g_2 \text{d}U + g_1 \text{d}V$ と書く．ここで $g_1, g_2 \in B$ である．
イデアル $J = (g_1, g_2) \subset B$ は，「座標」$U$ と $V$ の
選び方によらないことに注意する（$g$ は固定する）．
$\mathfrak p B_\mathfrak q + (U, V)^3B_\mathfrak q$ を法として見ると，
$g_1$ と $g_2$ は $U$ と $V$ に，
$\mathfrak p B_\mathfrak q + (U, V)^2B_\mathfrak q$ を法として合同である．
したがって上の議論を，$U$ を $g_1$ で，
$V$ を $g_2$ で置き換えてやり直せる．すると $J = (U, V)$ を得るので，
$g_1, g_2 \in (U, V) = J$ が従う（もちろん，これらは以前の座標の選択に
対する $g_1$ および $g_2$ とは異なる）．

\medskip\noindent
$\Spec(B/J) \to S = \Spec(R)$ は $x$ において \'etale であることに注意する．
したがって $S$ を \'etale 近傍で置き換え，$Y$ を再び縮小した後，
$R \to B \to B/J$ は同型であると仮定してよい．
$r \in R$ を，$B/J$ における像が $-g$ の像となる元とする．このとき
$$
g + r \equiv a U + b V \bmod (U, V)^2
$$
と，ある $a, b \in R$ に対して書ける．$\text{d}g \in J\Omega_{B/R}$ を用いて
微分を取れば，$a = b = 0$ を得る．したがって
$$
g + r = \alpha U^2 + \beta UV + \gamma V^2
$$
と，ある $\alpha, \beta, \gamma \in B$ に対して書ける．すると
$\beta$ の像は $1$（$\kappa(x)$ において）であり，$\alpha$ と $\gamma$ の像は
$0$（$\kappa(x)$ において）である．したがって $\beta$ と
$\alpha + \beta + \gamma$ は $B$ で可逆であると仮定してよい．
$Y'$ を $Y$ の次で与えられる被覆とする：
$$
Y' = \Spec(B[t]/(t(1 - t) - \alpha \gamma / \beta^2))
$$
ここで $x'$ は $x$ の上にある，$t = 1$ を満たす一意な点である．
$Y' \to Y$ は $x'$ において \'etale であることに注意する．
元
$$
T_1 = (\alpha + t \beta)U + ((1 - t)\beta + \gamma)V
\quad\text{かつ}\quad
T_2 =
\frac{(\alpha + (1 - t) \beta)U + (t\beta + \gamma)V}{\alpha + \beta + \gamma}
$$
を $\Gamma(Y', \mathcal{O}_{Y'})$ の元として考える．このとき
$$
g = T_1 T_2 - r
$$
が成り立つ．$T_1$ と $T_2$ も上の意味で座標なので，射
$Y' \to \mathbf{A}^2_S$ は $T_1$ と $T_2$ によって与えられ，
$x'$ の近傍で \'etale であり，
引き戻し $X' \subset Y'$，すなわち $X$ の引き戻しは
$\mathbf{A}^2_S = \Spec(R[x_1, x_2])$ の閉部分概形で，方程式
$x_1x_2 - r = 0$ によって定義されるものの逆像である
（ここでも必要なら $Y'$ を縮小する）．所望の図式を得るために
すべてを組み合わせることは読者に委ねる．
\end{proof}




\section{さらなる消滅結果}
\label{curves-section-more-vanishing}

\noindent
節 \ref{curves-section-vanishing} の続きである．

\begin{lemma}
\label{curves-lemma-h1-nonzero-degree-leq-2g-2}
状況 \ref{curves-situation-Cohen-Macaulay-curve} において，$X$ は整で，
種数 $g$ をもつと仮定する．$\mathcal{L}$ を可逆 $\mathcal{O}_X$-加群とする．
$Z \subset X$ を $0$ 次元閉部分概形で，イデアル層
$\mathcal{I} \subset \mathcal{O}_X$ をもつものとする．
$H^1(X, \mathcal{I}\mathcal{L})$ が
零でないならば，
$$
\deg(\mathcal{L}) \leq 2g - 2 + \deg(Z)
$$
が成り立ち，$\mathcal{I}\mathcal{L} \cong \omega_X$ でない限り不等号は狭義である．
\end{lemma}

\begin{proof}
任意の曲線，例えば $X$ は Cohen--Macaulay である．
$H^1(X, \mathcal{I}\mathcal{L})$ が零でないならば，零でない写像
$\mathcal{I}\mathcal{L} \to \omega_X$ が存在する；補題
\ref{curves-lemma-duality-dim-1-CM} を参照せよ．
$\mathcal{I}\mathcal{L}$ は捩れをもたないので，この写像は単射である．
体は Gorenstein であり $X$ は被約なので，Gorenstein 軌跡
$U \subset X$（$X$ の Gorenstein 軌跡）は空でない；Duality for Schemes，補題
\ref{duality-lemma-gorenstein} を参照せよ．
この補題は $\omega_X|_U$ が可逆であることも示す．
このようにして短完全列
$$
0 \to \mathcal{I}\mathcal{L} \to \omega_X \to \mathcal{Q} \to 0
$$
を得る．ここで $\mathcal{Q}$ の台は零次元である．したがって
\begin{align*}
0 & \leq \dim \Gamma(X, \mathcal{Q})\\
& =
\chi(\mathcal{Q}) \\
& =
\chi(\omega_X) - \chi(\mathcal{I}\mathcal{L}) \\
& =
\chi(\omega_X) - \deg(\mathcal{L}) - \chi(\mathcal{I}) \\
& =
2g - 2 - \deg(\mathcal{L}) + \deg(Z)
\end{align*}
を得る．ここで補題 \ref{curves-lemma-euler} と \ref{curves-lemma-rr}，式
(\ref{curves-equation-genus})，および Varieties，補題
\ref{varieties-lemma-chi-tensor-finite} と
\ref{varieties-lemma-degree-on-proper-curve} を用いた．さらに
$\deg(Z) = \dim_k \Gamma(Z, \mathcal{O}_Z) = \chi(\mathcal{O}_Z)$
および短完全列
$0 \to \mathcal{I} \to \mathcal{O}_X \to \mathcal{O}_Z \to 0$
も用いた．補題が従う．
\end{proof}

\begin{lemma}
\label{curves-lemma-degree-more-than-2g-2}
\begin{reference}
\cite[補題 2]{Jongmin}
\end{reference}
状況 \ref{curves-situation-Cohen-Macaulay-curve} において，
$X$ は整で種数 $g$ をもつと仮定する．
$\mathcal{L}$ を可逆 $\mathcal{O}_X$-加群とする．
$Z \subset X$ を $0$ 次元閉部分概形で，イデアル層
$\mathcal{I} \subset \mathcal{O}_X$ をもつものとする．
$\deg(\mathcal{L}) > 2g - 2 + \deg(Z)$ ならば，
$H^1(X, \mathcal{I}\mathcal{L}) = 0$ であり，次のいずれかが起こる．
\begin{enumerate}
\item $H^0(X, \mathcal{I}\mathcal{L}) \not = 0$，または
\item $g = 0$ かつ $\deg(\mathcal{L}) = \deg(Z) - 1$．
\end{enumerate}
(2) の場合に $Z = \emptyset$ ならば，$X \cong \mathbf{P}^1_k$ であり，
$\mathcal{L}$ は $\mathcal{O}_{\mathbf{P}^1}(-1)$ に対応する．
\end{lemma}

\begin{proof}
$H^1(X, \mathcal{I}\mathcal{L})$ の消滅は補題
\ref{curves-lemma-h1-nonzero-degree-leq-2g-2} から従う．
$H^0(X, \mathcal{I}\mathcal{L}) = 0$ ならば，
$\chi(\mathcal{I}\mathcal{L}) = 0$ である．短完全列
$0 \to \mathcal{I}\mathcal{L} \to \mathcal{L} \to \mathcal{O}_Z \to 0$
から $\deg(\mathcal{L}) = g - 1 + \deg(Z)$ と結論する．
したがって $g - 1 + \deg(Z) > 2g - 2 + \deg(Z)$ であり，これは
$g = 0$ を含意するので (2) が成り立つ．(2) の場合に
$Z = \emptyset$ ならば，$\mathcal{L}^{-1}$ は次数 $1$ の可逆層である．
これは同型 $X \to \mathbf{P}^1_k$ が存在し，$\mathcal{L}^{-1}$ が
$\mathcal{O}_{\mathbf{P}^1}(1)$ の引き戻しであることを含意する；補題
\ref{curves-lemma-genus-zero-positive-degree} による．
\end{proof}

\begin{lemma}
\label{curves-lemma-degree-more-than-2g}
\begin{reference}
\cite[補題 3]{Jongmin}
\end{reference}
状況 \ref{curves-situation-Cohen-Macaulay-curve} において，
$X$ は整で種数 $g$ をもつと仮定する．
$\mathcal{L}$ を可逆 $\mathcal{O}_X$-加群とする．
$\deg(\mathcal{L}) \geq 2g$ ならば，$\mathcal{L}$ は大域生成である．
\end{lemma}

\begin{proof}
$Z \subset X$ を，$\mathcal{L}$ の大域切断が切り出す閉部分概形とする．
補題 \ref{curves-lemma-degree-more-than-2g-2} により $Z \not = X$ である．
$\mathcal{I} \subset \mathcal{O}_X$ を $Z$ を切り出すイデアル層とする．
短完全列
$$
0 \to \mathcal{I}\mathcal{L}
\to \mathcal{L} \to \mathcal{O}_Z \to 0
$$
を考える．$Z \not = \emptyset$ ならば，
$H^1(X, \mathcal{I}\mathcal{L})$ は零でない．これはコホモロジーの
長完全列から従う．補題 \ref{curves-lemma-duality-dim-1-CM} により，
零でなく，したがって単射である写像
$$
\mathcal{I}\mathcal{L}
\longrightarrow
\omega_X
$$
を得る．特に，単射
$H^0(X, \mathcal{L}) = H^0(X, \mathcal{I}\mathcal{L})
\to H^0(X, \omega_X)$ を得る．しかしこれは不可能である．実際，
$$
\dim_k H^0(X, \mathcal{L}) = \dim_k H^1(X, \mathcal{L}) +
\deg(\mathcal{L}) + 1 - g \geq g + 1
$$
であり，式 (\ref{curves-equation-genus}) により
$\dim H^0(X, \omega_X) = g$ だからである．
\end{proof}

\begin{lemma}
\label{curves-lemma-degree-more-than-2g-1-and-Z}
状況 \ref{curves-situation-Cohen-Macaulay-curve} において，
$X$ は整で種数 $g$ をもつと仮定する．
$\mathcal{L}$ を可逆 $\mathcal{O}_X$-加群とする．
$Z \subset X$ を空でない $0$ 次元閉部分概形とする．
$\deg(\mathcal{L}) \geq 2g - 1 + \deg(Z)$ ならば，$\mathcal{L}$ は
大域生成であり，$H^0(X, \mathcal{L}) \to H^0(X, \mathcal{L}|_Z)$ は全射である．
\end{lemma}

\begin{proof}
補題 \ref{curves-lemma-degree-more-than-2g} により大域生成である．
$\mathcal{I} \subset \mathcal{O}_X$ が $Z$ のイデアル層ならば，
$H^1(X, \mathcal{I}\mathcal{L}) = 0$ が補題
\ref{curves-lemma-h1-nonzero-degree-leq-2g-2} により成り立つ．したがって全射である．
\end{proof}

\begin{lemma}
\label{curves-lemma-degree-at-least-2g+1}
状況 \ref{curves-situation-Cohen-Macaulay-curve} において，
$X$ は $k$ 上幾何学的に整で，種数 $g$ をもつと仮定する．
$\mathcal{L}$ を可逆 $\mathcal{O}_X$-加群とする．
$\deg(\mathcal{L}) \geq 2g + 1$ ならば，$\mathcal{L}$ は非常に豊富である．
\end{lemma}

\begin{proof}
補題 \ref{curves-lemma-degree-more-than-2g} により，$\mathcal{L}$ は大域生成であるから，
射 $f : X \to \mathbf{P}^n_k$ を定める．ここで
$n = h^0(X,\mathcal{L}) - 1$ である．$\mathcal{L}$ が非常に豊富であることを
示すとは，$f$ が閉埋め込みであることを示すことである．
$f$ を代数閉包 $\overline{k}$（$k$ の代数閉包）へ基底変換したものが閉埋め込みであることを
確認すれば十分である（Descent，補題
\ref{descent-lemma-descending-property-closed-immersion}）．
したがって $k$ は代数閉であると仮定してよい；仮定により $X$ は依然として整である．
補題 \ref{curves-lemma-degree-more-than-2g-1-and-Z} により，任意の $0$ 次元閉部分概形
$Z\subset X$ で次数 2 のものに対して，制限写像
$H^0(X, \mathcal{L}) \to H^0(X, \mathcal{L}|_Z)$ は全射である．
Varieties，補題
\ref{varieties-lemma-variant-separate-points-tangent-vectors}
により，$\mathcal{L}$ は非常に豊富である．
\end{proof}


\begin{lemma}
\label{curves-lemma-vanishing-on-gorenstein}
\begin{reference}
\cite[補題 4]{Jongmin} の弱い形
\end{reference}
$k$ を体とする．$X$ を $k$ 上の固有概形で，被約かつ連結で，
次元 $1$ のものとする．
$\mathcal{L}$ を可逆 $\mathcal{O}_X$-加群とする．
$Z \subset X$ を $0$ 次元閉部分概形で，イデアル層
$\mathcal{I} \subset \mathcal{O}_X$ をもつものとする．
$H^1(X, \mathcal{I}\mathcal{L}) \not = 0$ ならば，
被約連結閉部分概形 $Y \subset X$ で，次元 $1$ かつ
$$
\deg(\mathcal{L}|_Y) \leq -2\chi(Y, \mathcal{O}_Y) + \deg(Z \cap Y)
$$
を満たすものが存在する．ここで $Z \cap Y$ は概形論的交わりである．
\end{lemma}

\begin{proof}
$H^1(X, \mathcal{I}\mathcal{L})$ が零でないならば，零でない写像
$\varphi : \mathcal{I}\mathcal{L} \to \omega_X$ が存在する；補題
\ref{curves-lemma-duality-dim-1-CM} を参照せよ．
$Y \subset X$ を，既約成分 $C$（$X$ の既約成分）のうち，
$\varphi$ が $C$ の一般点で零でないものすべての合併とする．
すると $Y$ は被約閉部分概形である．
$\mathcal{J} \subset \mathcal{O}_X$ を $Y$ のイデアル層とする．
$\mathcal{J}\mathcal{I}\mathcal{L}$ は（$\mathcal{L}$ の部分加群として）
埋め込まれた随伴点をもたず，$\varphi$ は $\mathcal{J}$ の台の一般点で零である
（$Y$ の選び方と $X$ が被約であることによる）ので，$\varphi$ は
$$
\mathcal{I}\mathcal{L} \to
\mathcal{I}\mathcal{L}/\mathcal{J}\mathcal{I}\mathcal{L} \to \omega_X
$$
と分解する．$\mathcal{I}\mathcal{L}/\mathcal{J}\mathcal{I}\mathcal{L}$ は
$Y$ 上の連接層の順像と見なせる．記号の濫用により，その連接層も同じ記号で表す．
補題 \ref{curves-lemma-closed-immersion-dim-1-CM} により
$\omega_Y = \SheafHom(\mathcal{O}_Y, \omega_X)$ なので，
$$
\mathcal{I}\mathcal{L}/
\mathcal{J}\mathcal{I}\mathcal{L} 
\to \omega_Y
$$
という $\mathcal{O}_Y$-加群の写像で，$Y$ の一般点で単射であるものを得る．
$\mathcal{I}' \subset \mathcal{O}_Y$ を $Z \cap Y$ のイデアル層とする．
写像
$\mathcal{I}\mathcal{L}/\mathcal{J}\mathcal{I}\mathcal{L} \to
\mathcal{I}'\mathcal{L}|_Y$ で，核が閉点に台をもつものが存在する．
$\omega_Y$ は Cohen--Macaulay 加群なので，上の写像は単射
$\mathcal{I}'\mathcal{L}|_Y \to \omega_Y$ を経由して分解する．
したがって完全列
$$
0 \to \mathcal{I}'\mathcal{L}|_Y \to \omega_Y \to \mathcal{Q} \to 0
$$
を得る．これは $Y$ 上の連接層の完全列で，$\mathcal{Q}$ は次元 $0$ に台をもつ
（これには $\omega_Y$ が $Y$ の一般点で可逆加群であることを用いる）．
したがって
$$
0 \leq \dim \Gamma(Y, \mathcal{Q}) =
\chi(\mathcal{Q}) = \chi(\omega_Y) - \chi(\mathcal{I}'\mathcal{L}) =
-2\chi(\mathcal{O}_Y) - \deg(\mathcal{L}|_Y) + \deg(Z \cap Y)
$$
と結論する．これは補題 \ref{curves-lemma-euler} および
Varieties，補題 \ref{varieties-lemma-chi-tensor-finite} による．
$Y$ が連結ならば，これで補題が証明される．連結でなければ，
$Y$ の連結成分の一つについて議論の最後の部分を繰り返す．
\end{proof}

\begin{lemma}
\label{curves-lemma-global-generation}
$k$ を体とする．$X$ を $k$ 上の固有概形で，被約かつ連結で，
次元 $1$ のものとする．
$\mathcal{L}$ を可逆 $\mathcal{O}_X$-加群とする．
任意の被約連結閉部分概形 $Y \subset X$ で次元 $1$ のものに対して
$$
\deg(\mathcal{L}|_Y) \geq 2\dim_k H^1(Y, \mathcal{O}_Y)
$$
が成り立つと仮定する．このとき $\mathcal{L}$ は大域生成である．
\end{lemma}

\begin{proof}
$X$ の既約成分の数に関する帰納法による．
$X$ が既約ならば，$X$ を体
$k' = H^0(X, \mathcal{O}_X)$ 上の概形と見て補題
\ref{curves-lemma-degree-more-than-2g} を適用すれば従う．
$X$ は既約でないと仮定する．先へ進む前に，$k$ が有限ならば，
$k$ を純超越拡大 $K$ で置き換える．これは Varieties，補題
\ref{varieties-lemma-globally-generated-base-change},
\ref{varieties-lemma-degree-base-change},
\ref{varieties-lemma-geometrically-reduced-any-base-change}，および
\ref{varieties-lemma-bijection-irreducible-components}，
Cohomology of Schemes，補題
\ref{coherent-lemma-flat-base-change-cohomology}，
補題 \ref{curves-lemma-sanity-check-duality}，ならびに
$K$ が $k$ 上幾何学的に整であるという初等的事実によって許される．

\medskip\noindent
矛盾を得るため，$\mathcal{L}$ は大域生成でないと仮定する．
すると連接イデアル層 $\mathcal{I} \subset \mathcal{O}_X$ を，
$H^0(X, \mathcal{I}\mathcal{L}) = H^0(X, \mathcal{L})$ であり，
$\mathcal{O}_X/\mathcal{I}$ が零でなく，台が次元 $0$ となるように選べる．
例えば，$\mathcal{I}$ を，任意の閉点であって $\mathcal{L}$ の大域切断の
共通零点軌跡にあるもののイデアル層とすればよい．短完全列
$$
0 \to \mathcal{I}\mathcal{L} \to \mathcal{L} \to 
\mathcal{L}/\mathcal{I}\mathcal{L} \to 0
$$
を考える．$\mathcal{L}/\mathcal{I}\mathcal{L}$ の台は次元 $0$ なので，
$\mathcal{L}/\mathcal{I}\mathcal{L}$ は大域切断で生成される
（Varieties，補題 \ref{varieties-lemma-chi-tensor-finite}）．
短完全列と
$H^0(X, \mathcal{I}\mathcal{L}) = H^0(X, \mathcal{L})$
から単射
$H^0(X, \mathcal{L}/\mathcal{I}\mathcal{L}) \to
H^1(X, \mathcal{I}\mathcal{L})$
を得る．

\medskip\noindent
$k$-ベクトル空間 $H^1(X, \mathcal{I}\mathcal{L})$ は
$\Hom(\mathcal{I}\mathcal{L}, \omega_X)$ の双対であることを思い出そう．
$\varphi : \mathcal{I}\mathcal{L} \to \omega_X$ を選ぶ．
補題 \ref{curves-lemma-vanishing-on-gorenstein} により
$H^1(X, \mathcal{L}) = 0$ である．したがって
$$
\dim_k H^0(X, \mathcal{I}\mathcal{L}) = \dim_k H^0(X, \mathcal{L}) =
\deg(\mathcal{L}) + \chi(\mathcal{O}_X) > \dim_k H^1(X, \mathcal{O}_X) =
\dim_k H^0(X, \omega_X)
$$
である．$\varphi$ は大域切断上で単射でなく，特に
$\varphi$ は単射でないと結論する．任意の一般点
$\eta \in X$ で，$X$ の既約成分の一般点であるものに対して，
$V_\eta \subset \Hom(\mathcal{I}\mathcal{L}, \omega_X)$ を，
$k$-部分ベクトル空間であって，$\varphi$ で $\eta$ において零となるもの
からなるものとする．
$\mathcal{I}\mathcal{L}$ のすべての随伴点は $X$ の一般点なので，
上のことから
$\Hom(\mathcal{I}\mathcal{L}, \omega_X) = \bigcup V_\eta$
を得る．$X$ は有限個の一般点をもち，$k$ は無限なので，
$\Hom(\mathcal{I}\mathcal{L}, \omega_X) = V_\eta$ が
ある $\eta$ に対して成り立つと結論する．
$\eta \in C \subset X$ を対応する既約成分とする．
$Y \subset X$ を $X$ の他の既約成分の合併とする．
すると $Y$ は空でない被約閉部分概形であり，$X$ と等しくない．
$\mathcal{J} \subset \mathcal{O}_X$ を $Y$ のイデアル層とする．
$\mathcal{J}$ の台は $C$ であることを心に留めておく．

\medskip\noindent
$\varphi : \mathcal{I}\mathcal{L} \to \omega_X$ を任意にとる．
$\mathcal{J}\mathcal{I}\mathcal{L}$ は（$\mathcal{L}$ の部分加群として）
埋め込まれた随伴点をもたず，$\varphi$ は一般点 $\eta$（$\mathcal{J}$ の
台の一般点）で零なので，$\varphi$ は
$$
\mathcal{I}\mathcal{L} \to
\mathcal{I}\mathcal{L}/\mathcal{J}\mathcal{I}\mathcal{L} \to \omega_X
$$
と分解する．$\mathcal{I}\mathcal{L}/\mathcal{J}\mathcal{I}\mathcal{L}$ を
$Y$ 上の連接層の順像と見なせる．記号の濫用により，その層も同じ記号で表す．
補題 \ref{curves-lemma-closed-immersion-dim-1-CM} により
$\omega_Y = \SheafHom(\mathcal{O}_Y, \omega_X)$ なので，分解
$$
\mathcal{I}\mathcal{L} \to
\mathcal{I}\mathcal{L}/
\mathcal{J}\mathcal{I}\mathcal{L} 
\xrightarrow{\varphi'} \omega_Y \to \omega_X
$$
を得る．これは $\varphi$ の分解である．$\mathcal{I}' \subset \mathcal{O}_Y$ を，
$\mathcal{I} \subset \mathcal{O}_X$ の像とする．全射
$\mathcal{I}\mathcal{L}/\mathcal{J}\mathcal{I}\mathcal{L} \to
\mathcal{I}'\mathcal{L}|_Y$ で，核が閉点に台をもつものが存在する．
$\omega_Y$ は $Y$ 上の Cohen--Macaulay 加群なので，写像 $\varphi'$ は
写像 $\varphi'' : \mathcal{I}'\mathcal{L}|_Y \to \omega_Y$ を経由して分解する．
したがって可換図式
$$
\vcenter{
\xymatrix{
0 \ar[r] &
\mathcal{I}\mathcal{L} \ar[r] \ar[d] &
\mathcal{L} \ar[r] \ar[d] &
\mathcal{L}/\mathcal{I}\mathcal{L} \ar[r] \ar[d] &
0 \\
0 \ar[r] &
\mathcal{I}'\mathcal{L}|_Y \ar[r] &
\mathcal{L}|_Y \ar[r] &
\mathcal{L}|_Y/\mathcal{I}'\mathcal{L}|_Y \ar[r] &
0
}
}
\quad\text{かつ}\quad
\vcenter{
\xymatrix{
\mathcal{I}\mathcal{L} \ar[r]_\varphi \ar[d] & \omega_X \\
\mathcal{I}'\mathcal{L}|_Y \ar[r]^{\varphi''} & \omega_Y \ar[u]
}
}
$$
を得る．ここから次のように証明を終えられる．
各 $\varphi$ に対して $\varphi''$ が存在し，さらに
$\omega_X \in \textit{Coh}(\mathcal{O}_X)$ は関手
$\mathcal{F} \mapsto \Hom_k(H^1(X, \mathcal{F}), k)$ を表現するので，
写像
$H^1(X, \mathcal{I}\mathcal{L}) \to
H^1(Y, \mathcal{I}'\mathcal{L}|_Y)$
は単射である．境界写像
$H^0(X, \mathcal{L}/\mathcal{I}\mathcal{L}) \to
H^1(X, \mathcal{I}\mathcal{L})$
は単射なので，合成
$$
H^0(X, \mathcal{L}/\mathcal{I}\mathcal{L}) \to
H^0(X, \mathcal{L}|_Y/\mathcal{I}'\mathcal{L}|_Y) \to
H^1(X, \mathcal{I}'\mathcal{L}|_Y)
$$
は単射であると結論する．さらに
$\mathcal{L}/\mathcal{I}\mathcal{L} \to
\mathcal{L}|_Y/\mathcal{I}'\mathcal{L}|_Y$
は，台が次元 $0$ の連接加群の全射なので，第一の写像
$H^0(X, \mathcal{L}/\mathcal{I}\mathcal{L}) \to
H^0(X, \mathcal{L}|_Y/\mathcal{I}'\mathcal{L}|_Y)$
は全射（したがって全単射）である．しかし帰納法により
$\mathcal{L}|_Y$ は大域生成である（$Y$ が非連結でももちろん成り立つ）ので，
境界写像
$$
H^0(X, \mathcal{L}|_Y/\mathcal{I}'\mathcal{L}|_Y) \to
H^1(X, \mathcal{I}'\mathcal{L}|_Y)
$$
は単射ではありえない．この矛盾によって証明が完了する．
\end{proof}





\section{有理尾部の収縮}
\label{curves-section-contracting-rational-tails}

\noindent
本節では，標準層の正値性を改善するため概形を収縮する場合のうち，
最も単純なものを論じる．

\begin{example}[有理尾部の収縮]
\label{curves-example-rational-tail}
$k$ を体とする．$X$ を $k$ 上の固有概形で，次元 $1$ かつ
$H^0(X, \mathcal{O}_X) = k$ であるものとする．$X$ の特異点は
高々節点的であると仮定する．
{\it 有理尾部}とは，既約成分 $C \subset X$（整閉部分概形と見る）で，
次の性質をもつものをいう．
\begin{enumerate}
\item $X' \not = \emptyset$．ここで $X' \subset X$ は
$X \setminus C$ の概形論的閉包である，
\item 概形論的交わり $C \cap X'$ は一つの被約点 $x$ である，
\item $H^0(C, \mathcal{O}_C)$ は $x$ の剰余体へ同型に写る，かつ
\item $C$ の種数は零である．
\end{enumerate}
$X$ の既約成分で $x$ を通るものが少なくとも二つあるので，
$x$ は節点であると結論する．
$k' = H^0(C, \mathcal{O}_C) = \kappa(x)$ とおく．
このとき $k'/k$ は体の有限分離拡大である（補題 \ref{curves-lemma-nodal}）．
標準的な射
$$
c : X \longrightarrow X'
$$
で，$X'$ 上恒等写像を誘導し，$C$ を $x \in X'$ へ，
標準的な射 $C \to \Spec(k') = x$ によって写すものが存在する．
これは Morphisms，補題 \ref{morphisms-lemma-scheme-theoretic-union} から従う．
実際，$X$ は $C$ と $X'$ の概形論的合併である（$X$ は被約である）．
さらに，
$$
c_*\mathcal{O}_X = \mathcal{O}_{X'}
\quad\text{かつ}\quad
R^1c_*\mathcal{O}_X = 0
$$
であると主張する．これを見るため，$i_C : C \to X$，
$i_{X'} : X' \to X$，および $i_x : x \to X$ を埋め込みとし，
Morphisms，補題 \ref{morphisms-lemma-scheme-theoretic-union} の完全列
$$
0 \to \mathcal{O}_X \to
i_{C, *}\mathcal{O}_C \oplus i_{X', *}\mathcal{O}_{X'} \to
i_{x, *}\kappa(x) \to 0
$$
を用いる．高次順像の長完全列を見ると，
$H^0(C, \mathcal{O}_C) = k'$ および $H^1(C, \mathcal{O}_C) = 0$ を
示せば十分であり，これは仮定から従う．
$X'$ もまた $k$ 上の固有概形で，次元 $1$，特異点が高々節点的であり
（補題 \ref{curves-lemma-closed-subscheme-nodal-curve}），
$H^0(X', \mathcal{O}_{X'}) = k$ を満たし，$X'$ は $X$ と同じ種数をもつことに注意する．
$c : X \to X'$ を{\it 有理尾部の収縮}と呼ぶ．
\end{example}

\begin{lemma}
\label{curves-lemma-rational-tail-negative}
$k$ を体とする．$X$ を $k$ 上の固有概形で，次元 $1$ かつ
$H^0(X, \mathcal{O}_X) = k$ であるものとする．$X$ の特異点は
高々節点的であると仮定する．$C \subset X$ を有理尾部とする
（例 \ref{curves-example-rational-tail}）．このとき $\deg(\omega_X|_C) < 0$ である．
\end{lemma}

\begin{proof}
$X' \subset X$ を例におけるものとする．このとき短完全列
$$
0 \to \omega_C \to \omega_X|_C \to \mathcal{O}_{C \cap X'} \to 0
$$
を得る．補題 \ref{curves-lemma-closed-subscheme-reduced-gorenstein}，
\ref{curves-lemma-facts-about-nodal-curves}，および
\ref{curves-lemma-closed-subscheme-nodal-curve} を参照せよ．
例と同じ $k'$ に対して，$\deg(\omega_C) = -2[k' : k]$ である．
実際，命題 \ref{curves-proposition-projective-line} により
$C \cong \mathbf{P}^1_{k'}$ であり，また
$\deg(C \cap X') = [k' : k]$ である．
したがって $\deg(\omega_X|_C) = -[k' : k]$ であり，これは負である．
\end{proof}

\begin{lemma}
\label{curves-lemma-rational-tail-field-extension}
$k$ を体とする．$X$ を $k$ 上の固有概形で，次元 $1$ かつ
$H^0(X, \mathcal{O}_X) = k$ であるものとする．$X$ の特異点は
高々節点的であると仮定する．$C \subset X$ を有理尾部とする
（例 \ref{curves-example-rational-tail}）．
任意の体拡大 $K/k$ に対して，基底変換 $C_K \subset X_K$ は
有限個の有理尾部の非交和である．
\end{lemma}

\begin{proof}
$x \in C$ と $k' = \kappa(x)$ を例におけるものとする．
命題 \ref{curves-proposition-projective-line} により
$C \cong \mathbf{P}^1_{k'}$ であることに注意する．
$k'/k$ は有限分離拡大なので，
$k' \otimes_k K = K'_1 \times \ldots \times K'_n$
は有限分離拡大 $K'_i/K$ の有限積である．
$C_i = \mathbf{P}^1_{K'_i}$ とおき，$x_i \in C_i$ を
$x$ の逆像とする．すると $C_K = \coprod C_i$ かつ
$X'_K \cap C_i = x_i$ であり，所望の結論を得る．
\end{proof}

\begin{lemma}
\label{curves-lemma-no-rational-tail}
$k$ を体とする．$X$ を $k$ 上の固有概形で，次元 $1$ かつ
$H^0(X, \mathcal{O}_X) = k$ であるものとする．$X$ の特異点は
高々節点的であると仮定する．$X$ が有理尾部をもたないならば
（例 \ref{curves-example-rational-tail}），任意の被約連結閉部分概形
$Y \subset X$ で，$Y \not = X$ かつ次元 $1$ のものに対して
$\deg(\omega_X|_Y) \geq \dim_k H^1(Y, \mathcal{O}_Y)$
が成り立つ．
\end{lemma}

\begin{proof}
$Y \subset X$ を主張におけるものとする．このとき
$k' = H^0(Y, \mathcal{O}_Y)$ は体であり，$k$ の有限拡大である．また
$[k' : k]$ は，以下に現れる $Y$ と $Y$ 上の連接層に付随する
すべての数値的不変量を割り切る；Varieties，補題
\ref{varieties-lemma-divisible} を参照せよ．
$Z \subset X$ を補題 \ref{curves-lemma-closed-subscheme-reduced-gorenstein}
におけるものとする．この補題および補題
\ref{curves-lemma-facts-about-nodal-curves} と
\ref{curves-lemma-closed-subscheme-nodal-curve} の結果を，以下では断らずに用いる．
すると短完全列
$$
0 \to \omega_Y \to \omega_X|_Y \to \mathcal{O}_{Y \cap Z} \to 0
$$
を得る．補題 \ref{curves-lemma-closed-subscheme-reduced-gorenstein} を参照せよ．
したがって
$$
\deg(\omega_X|_Y) = \deg(Y \cap Z) + \deg(\omega_Y) =
\deg(Y \cap Z) - 2\chi(Y, \mathcal{O}_Y)
$$
を得る．ゆえに補題が偽ならば，
$$
2[k' : k] > \deg(Y \cap Z) + \dim_k H^1(Y, \mathcal{O}_Y)
$$
となる．$Y \cap Z$ は空でなく，上述の可除性もあるので，これが起こりうるのは
$Y \cap Z$ が一つの $k'$-有理点で，$Y$ の滑らかな軌跡にあり，
$H^1(Y, \mathcal{O}_Y) = 0$ の場合に限る．
$Y$ が既約ならば，これは $Y$ が有理尾部であることを含意する．
$Y$ が可約ならば，$\deg(\omega_X|_Y) = -[k' : k]$ なので，
ある既約成分 $C$（$Y$ の既約成分）で $\deg(\omega_X|_C) < 0$ となるものが存在する；
Varieties，補題 \ref{varieties-lemma-degree-in-terms-of-components} を参照せよ．
すると上の解析を $C$ に適用して，$C$ が有理尾部であることを得る．
\end{proof}

\begin{lemma}
\label{curves-lemma-no-rational-tail-semiample-genus-geq-2}
$k$ を体とする．$X$ を $k$ 上の固有概形で，次元 $1$ かつ
$H^0(X, \mathcal{O}_X) = k$ であるものとする．$X$ の特異点は
高々節点的であると仮定する．$X$ は有理尾部をもたないと仮定する
（例 \ref{curves-example-rational-tail}）．このとき次が成り立つ．
\begin{enumerate}
\item $X$ の種数が $0$ ならば，$X$ は既約平面二次曲線と同型であり，
$\omega_X^{\otimes -1}$ は非常に豊富である，
\item $X$ の種数が $1$ ならば，$\omega_X \cong \mathcal{O}_X$ である，
\item $X$ の種数が $\geq 2$ ならば，$\omega_X^{\otimes m}$ は
$m \geq 2$ に対して大域生成である．
\end{enumerate}
\end{lemma}

\begin{proof}
補題 \ref{curves-lemma-facts-about-nodal-curves} により $X$ は Gorenstein であり，
すなわち $\omega_X$ は可逆 $\mathcal{O}_X$-加群である．

\medskip\noindent
$X$ の種数が零ならば $\deg(\omega_X) < 0$ である．したがって
$X$ が複数の既約成分をもてば，補題 \ref{curves-lemma-no-rational-tail} に
矛盾する．既約の場合，$X$ は既約平面二次曲線と同型であり，
補題 \ref{curves-lemma-genus-zero} により $\omega_X^{\otimes -1}$ は非常に豊富である．

\medskip\noindent
$X$ の種数が $1$ ならば，$\omega_X$ は大域切断をもち，
すべての既約成分に対して $\deg(\omega_X|_C) = 0$ である．
実際，補題 \ref{curves-lemma-no-rational-tail} により
$\deg(\omega_X|_C) \geq 0$ がすべての既約成分 $C$ に対して成り立ち，補題
\ref{curves-lemma-genus-gorenstein} によりこれらの数の和は $0$ である．また
Varieties，補題 \ref{varieties-lemma-degree-in-terms-of-components} を適用できる．
すると Varieties，補題 \ref{varieties-lemma-no-sections-dual-nef} により
$\omega_X \cong \mathcal{O}_X$ である．

\medskip\noindent
種数 $g$（$X$ の種数）は $2$ 以上であると仮定する．
$X$ が既約ならば，補題 \ref{curves-lemma-degree-more-than-2g} により終わる．
$X$ は可約であると仮定する．補題 \ref{curves-lemma-no-rational-tail} により，
補題 \ref{curves-lemma-global-generation} の不等式は，主張における
任意の $Y \subset X$ に対して，$Y = X$ の場合を除いて成り立つ．
補題 \ref{curves-lemma-global-generation} の証明を解析すると，（可約の場合）
$Y = X$ に対して用いられる不等式は
$$
\deg(\omega_X^{\otimes m}) > -2 \chi(\mathcal{O}_X)
\quad\text{かつ}\quad
\deg(\omega_X^{\otimes m}) + \chi(\mathcal{O}_X) > \dim_k H^1(X, \mathcal{O}_X)
$$
だけであることが分かる．これらはいずれも $g \geq 2$ と $m \geq 2$ の
仮定の下で成り立つので，結論を得る．
\end{proof}

\begin{lemma}
\label{curves-lemma-contracting-rational-tails}
$k$ を体とする．$X$ を $k$ 上の固有概形で，次元 $1$ かつ
$H^0(X, \mathcal{O}_X) = k$ であるものとする．$X$ の特異点は
高々節点的であると仮定する．有理尾部の収縮（例
\ref{curves-example-rational-tail}）の列
$$
X = X_0 \to X_1 \to \ldots \to X_n = X'
$$
を，有理尾部がなくなるまで考える．このとき
\begin{enumerate}
\item $X$ の種数が $0$ ならば，$X'$ は既約平面二次曲線である，
\item $X$ の種数が $1$ ならば，$\omega_{X'} \cong \mathcal{O}_X$ である，
\item $X$ の種数が $> 1$ ならば，$\omega_{X'}^{\otimes m}$ は
$m \geq 2$ に対して大域生成である．
\end{enumerate}
$X$ の種数が $\geq 1$ ならば，射 $X \to X'$ は選択によらず，
この射の形成は基礎体の拡大と可換である．
\end{lemma}

\begin{proof}
有理尾部がなくなるまで順次収縮する．すると (1)，(2)，(3) は補題
\ref{curves-lemma-no-rational-tail-semiample-genus-geq-2} から従う．

\medskip\noindent
一意性．$f : X \to X'$ が選択によらないことを見るには，任意の有理尾部
$C \subset X$ が $X \to X'$ により一点へ写ることを示せば十分である；
細部は省略する．そうでなければ，切断
$s \in \Gamma(X', \omega_{X'}^{\otimes 2})$ で，既約成分
$f(C)$ の一般点で消えないものを見いだせる．
各収縮 $X_i \to X_{i + 1}$ には切断 $X_{i + 1} \to X_i$ があるので，
切断 $X' \to X$（$f$ の切断）が存在する．すると完全列
$$
0 \to \omega_{X'} \to \omega_X \to \omega_X|_{X''} \to 0
$$
を得る．ここで $X'' \subset X$ は $f$ により収縮される既約成分の合併である．
補題 \ref{curves-lemma-closed-subscheme-reduced-gorenstein} を参照せよ．
したがって写像 $\omega_{X'}^{\otimes 2} \to \omega_X^{\otimes 2}$ を得る．
$s$ の像を取れば，$\omega_X^{\otimes 2}$ の切断で，$C$ の一般点で
消えないものを得る．これは $\omega_X$ の有理尾部への制限が負の次数をもつ
という事実（補題 \ref{curves-lemma-rational-tail-negative}）に矛盾する．

\medskip\noindent
基礎体の拡大に関する主張は補題
\ref{curves-lemma-rational-tail-field-extension} から従う．細部は省略する．
\end{proof}





\section{有理橋の収縮}
\label{curves-section-contracting-rational-bridges}

\noindent
本節では，標準層の正値性を改善するため概形を収縮する場合のうち，
節 \ref{curves-section-contracting-rational-tails} で論じた場合に次いで
単純なものを論じる．

\begin{example}[有理橋の収縮]
\label{curves-example-rational-bridge}
$k$ を体とする．$X$ を $k$ 上の固有概形で，次元 $1$ かつ
$H^0(X, \mathcal{O}_X) = k$ であるものとする．$X$ の特異点は
高々節点的であると仮定する．
{\it 有理橋}とは，既約成分 $C \subset X$（整閉部分概形と見る）で，
次の性質をもつものをいう．
\begin{enumerate}
\item $X' \not = \emptyset$．ここで $X' \subset X$ は
$X \setminus C$ の概形論的閉包である，
\item 概形論的交わり $C \cap X'$ は
次数 $2$ を $H^0(C, \mathcal{O}_C)$ 上にもつ，かつ
\item $C$ の種数は零である．
\end{enumerate}
$k' = H^0(C, \mathcal{O}_C)$ および
$k'' = H^0(C \cap X', \mathcal{O}_{C \cap X'})$ とおく．
このとき $k'$ は体であり（Varieties，補題
\ref{varieties-lemma-proper-geometrically-reduced-global-sections}），
$\dim_{k'}(k'') = 2$ である．$X$ の既約成分が少なくとも二つ，
$C \cap X'$ の各点を通るので，これらの点は $X$ の節点であり，
$C$ と $X'$ の双方の滑らかな点であると結論する
（補題 \ref{curves-lemma-closed-subscheme-nodal-curve}）．
したがって $k'/k$ は体の有限分離拡大であり，$k''/k'$ は
体の次数 $2$ の分離拡大であるか，または
$k'' = k' \times k'$ である（補題 \ref{curves-lemma-nodal}）．
節 \ref{curves-section-pushouts} により押し出し
$$
\xymatrix{
C \cap X' \ar[r] \ar[d] &
X' \ar[d]^a \\
\Spec(k') \ar[r] &
Y
}
$$
が存在し，多くのよい性質をもつ（以下ではすべて断らずに用いる）．
$y \in Y$ を $\Spec(k') \to Y$ の像とする．このとき
$$
\mathcal{O}_{Y, y}^\wedge \cong k'[[s, t]]/(st)
\quad\text{または}\quad
\mathcal{O}_{Y, y}^\wedge \cong
\{f \in k''[[s]] : f(0) \in k'\}
$$
であり，どちらであるかは $C \cap X'$ が $2$ 個の点をもつか
$1$ 個の点をもつかによる．これは補題
\ref{curves-lemma-complete-local-ring-pushout} と，
$\mathcal{O}_{X', p} \cong \kappa(p)[[t]]$ が
$p \in C \cap X'$ に対して成り立つという事実から従う；More on Algebra，補題
\ref{more-algebra-lemma-power-series-over-residue-field} を参照せよ．
したがって $y \in Y$ は節点である；補題 \ref{curves-lemma-nodal} と
\ref{curves-lemma-2-branches-delta-1}，特に補題
\ref{curves-lemma-2-branches-delta-1} の (2) $\Rightarrow$ (1) の証明における
場合 II の議論を参照せよ．したがって $Y$ の特異点は高々節点的である．

\medskip\noindent
上の可換図式を次の図式へ拡張できる：
$$
\xymatrix{
C \cap X' \ar[r] \ar[d] &
X' \ar[d]^a \ar[r] & X \ar[ld]^c &
C \ar[ld] \ar[l] \\
\Spec(k') \ar[r] &
Y &
\Spec(k') \ar[l]
}
$$
ここで下段の二つの水平矢印は同じである．実際，$X$ は $X'$ と $C$ の
概形論的合併であり（したがって Morphisms，補題
\ref{morphisms-lemma-scheme-theoretic-union} により押し出しである），
射 $C \to Y$ と $X' \to Y$ は $C \cap X'$ 上一致する．最後に，
$$
c_*\mathcal{O}_X = \mathcal{O}_Y
\quad\text{かつ}\quad
R^1c_*\mathcal{O}_X = 0
$$
であると主張する．これを見るには Morphisms，補題
\ref{morphisms-lemma-scheme-theoretic-union} の完全列
$$
0 \to \mathcal{O}_X \to
\mathcal{O}_C \oplus \mathcal{O}_{X'} \to
\mathcal{O}_{C \cap X'} \to 0
$$
を用いる．高次順像の長完全列は
$$
0 \to c_*\mathcal{O}_X \to c_*\mathcal{O}_C \oplus c_*\mathcal{O}_{X'} \to
c_*\mathcal{O}_{C \cap X'} \to
R^1c_*\mathcal{O}_X \to
R^1c_*\mathcal{O}_C \oplus R^1c_*\mathcal{O}_{X'}
$$
である．$c|_{X'} = a$ はアフィンなので
$R^1c_*\mathcal{O}_{X'} = 0$ である．
$c|_C$ は $C \to \Spec(k') \to X$ と分解し，$C$ の種数は零なので，
$R^1c_*\mathcal{O}_C = 0$ を得る．
$\mathcal{O}_{X'} \to \mathcal{O}_{C \cap X'}$ は全射で，
$c|_{X'}$ はアフィンなので，
$c_*\mathcal{O}_{X'} \to c_*\mathcal{O}_{C \cap X'}$
は全射である．これは $R^1c_*\mathcal{O}_X = 0$ を示す．最後に，
完全列と，More on Morphisms，命題
\ref{more-morphisms-proposition-pushout-along-closed-immersion-and-integral}
における押し出しの構造層の記述により
$\mathcal{O}_Y = c_*\mathcal{O}_X$ である．

\medskip\noindent
以上から，$Y$ もまた $k$ 上の固有概形で，次元 $1$ かつ
$H^0(Y, \mathcal{O}_Y) = k$ であり，特異点は高々節点的であり
（補題 \ref{curves-lemma-closed-subscheme-nodal-curve}），
$Y$ は $X$ と同じ種数をもつ．
$c : X \to Y$ を{\it 有理橋の収縮}と呼ぶ．
\end{example}

\begin{lemma}
\label{curves-lemma-rational-bridge-zero}
$k$ を体とする．$X$ を $k$ 上の固有概形で，次元 $1$ かつ
$H^0(X, \mathcal{O}_X) = k$ であるものとする．$X$ の特異点は
高々節点的であると仮定する．$C \subset X$ を有理橋とする
（例 \ref{curves-example-rational-bridge}）．このとき $\deg(\omega_X|_C) = 0$ である．
\end{lemma}

\begin{proof}
$X' \subset X$ を例におけるものとする．このとき短完全列
$$
0 \to \omega_C \to \omega_X|_C \to \mathcal{O}_{C \cap X'} \to 0
$$
を得る．補題 \ref{curves-lemma-closed-subscheme-reduced-gorenstein}，
\ref{curves-lemma-facts-about-nodal-curves}，および
\ref{curves-lemma-closed-subscheme-nodal-curve} を参照せよ．
例と同じ $k''/k'/k$ に対して，
$\deg(\omega_C) = -2[k' : k]$ である．これは $C$ の種数が $0$ だからである
（補題 \ref{curves-lemma-rr}）．また
$\deg(C \cap X') = [k'' : k] = 2[k' : k]$ である．
したがって $\deg(\omega_X|_C) = 0$ である．
\end{proof}

\begin{lemma}
\label{curves-lemma-rational-bridge-field-extension}
$k$ を体とする．$X$ を $k$ 上の固有概形で，次元 $1$ かつ
$H^0(X, \mathcal{O}_X) = k$ であるものとする．$X$ の特異点は
高々節点的であると仮定する．$C \subset X$ を有理橋とする
（例 \ref{curves-example-rational-bridge}）．任意の体拡大 $K/k$ に対して，
基底変換 $C_K \subset X_K$ は有限個の有理橋の非交和である．
\end{lemma}

\begin{proof}
$k''/k'/k$ を例におけるものとする．
$k'/k$ は有限分離拡大なので，
$k' \otimes_k K = K'_1 \times \ldots \times K'_n$
は有限分離拡大 $K'_i/K$ の有限積である．対応する積分解
$k'' \otimes_k K = \prod K''_i$ は次数 $2$ の分離代数拡大
$K''_i/K'_i$ を与える．$C_i = C_{K'_i}$ とおく．すると
$C_K = \coprod C_i$ であり，したがって各 $C_i$ は種数 $0$ である
（$K'_i$ 上の曲線と見る）．実際，平坦基底変換により
$H^1(C_K, \mathcal{O}_{C_K}) = 0$ だからである．最後に，
$X'_K \cap C_i = \Spec(K''_i)$ は次数 $2$ を $K'_i$ 上にもち，
所望の結論を得る．
\end{proof}

\begin{lemma}
\label{curves-lemma-rational-bridge-canonical}
$c : X \to Y$ を有理橋の収縮とする
（例 \ref{curves-example-rational-bridge}）．このとき
$c^*\omega_Y \cong \omega_X$ である．
\end{lemma}

\begin{proof}
これは直接計算でも証明できるが，ここでは次の特徴づけを用いる．
すなわち，$\omega_X$ は連接 $\mathcal{O}_X$-加群であり，関手
$\textit{Coh}(\mathcal{O}_X) \to \textit{Sets}$，
$\mathcal{F} \mapsto \Hom_k(H^1(X, \mathcal{F}), k) =
H^1(X, \mathcal{F})^\vee$
を表現する；補題
\ref{curves-lemma-duality-dim-1-CM} または Duality for Schemes，補題
\ref{duality-lemma-dualizing-module-proper-over-A} を参照せよ．

\medskip\noindent
正確に述べると，$\mathcal{C}_Y$ を，可逆 $\mathcal{O}_Y$-加群を対象とし，
$\mathcal{O}_Y$-加群準同型を射とする圏とする．
$\mathcal{C}_X$ を，可逆 $\mathcal{O}_X$-加群 $\mathcal{L}$ で
$\mathcal{L}|_C \cong \mathcal{O}_C$ を満たすものを対象とし，
$\mathcal{O}_Y$-加群準同型を射とする圏とする．関手
$$
c^* : \mathcal{C}_Y \to \mathcal{C}_X
$$
は圏同値であると主張する．実際，More on Morphisms，補題
\ref{more-morphisms-lemma-bijection-on-Pic} により本質的全射である．
次に射影公式（Cohomology，補題 \ref{cohomology-lemma-projection-formula}）から
$c_*c^*\mathcal{N} = \mathcal{N}$ が従い，したがって $c^*$ は
$c_*$ を準逆とする圏同値である．

\medskip\noindent
$\omega_X$ は $\mathcal{C}_X$ の対象であると主張する．実際，短完全列
$$
0 \to \omega_C \to \omega_X|_C \to \mathcal{O}_{C \cap X'} \to 0
$$
がある．補題 \ref{curves-lemma-closed-subscheme-reduced-gorenstein} を参照せよ．
次数を取ると $\deg(\omega_X|_C) = 0$ を得る（小さな細部は省略する）．
したがって補題 \ref{curves-lemma-genus-zero-pic} により $\omega_X|_C$ は自明であり，
$\omega_X$ は $\mathcal{C}_X$ の対象である．

\medskip\noindent
$R^1c_*\mathcal{O}_X = 0$ なので，射影公式により
$R^1c_*c^*\mathcal{N} = 0$ が $\mathcal{N} \in \Ob(\mathcal{C}_Y)$ に
対して成り立つ．したがって Leray スペクトル系列
（Cohomology，補題 \ref{cohomology-lemma-apply-Leray}）により，圏と関手の図式
$$
\xymatrix{
\mathcal{C}_Y \ar[rr]_{c^*} \ar[dr]_{H^1(Y, -)^\vee} & &
\mathcal{C}_X \ar[ld]^{H^1(X, -)^\vee} \\
& \textit{Sets}
}
$$
は可換である．$\omega_Y \in \Ob(\mathcal{C}_Y)$ は南東向きの矢印を表現し，
$\omega_X \in \Ob(\mathcal{C}_X)$ も南東向きの矢印を表現するので，
Yoneda の補題（Categories，補題 \ref{categories-lemma-yoneda}）により結論する．
\end{proof}

\begin{lemma}
\label{curves-lemma-no-rational-bridge-ample-genus-geq-2}
$k$ を体とする．$X$ を $k$ 上の固有概形で，次元 $1$ かつ
$H^0(X, \mathcal{O}_X) = k$ であるものとする．次を仮定する．
\begin{enumerate}
\item $X$ の特異点は高々節点的である，
\item $X$ は有理尾部をもたない
（例 \ref{curves-example-rational-tail}），
\item $X$ は有理橋をもたない
（例 \ref{curves-example-rational-bridge}），
\item 種数 $g$（$X$ の種数）は $\geq 2$ である．
\end{enumerate}
このとき $\omega_X$ は豊富である．
\end{lemma}

\begin{proof}
$\deg(\omega_X|_C) > 0$ を，すべての既約成分 $C$（$X$ の既約成分）について
示せば十分である；Varieties，補題
\ref{varieties-lemma-ampleness-in-terms-of-degrees-components} を参照せよ．
$X = C$ が既約ならば，これは $g \geq 2$ と補題
\ref{curves-lemma-genus-gorenstein} から従う．
そうでなければ $k' = H^0(C, \mathcal{O}_C)$ とおく．これは体で，
$k$ の有限拡大であり，$[k' : k]$ は以下に現れる $C$ と
$C$ 上の連接層に付随するすべての数値的不変量を割り切る；
Varieties，補題 \ref{varieties-lemma-divisible} を参照せよ．
$X' \subset X$ を，補題 \ref{curves-lemma-closed-subscheme-reduced-gorenstein}
における $X \setminus C$ の閉包とする．この補題および補題
\ref{curves-lemma-facts-about-nodal-curves} と
\ref{curves-lemma-closed-subscheme-nodal-curve} の結果を，以下では断らずに用いる．
すると短完全列
$$
0 \to \omega_C \to \omega_X|_C \to \mathcal{O}_{C \cap X'} \to 0
$$
を得る．補題 \ref{curves-lemma-closed-subscheme-reduced-gorenstein} を参照せよ．
したがって
$$
\deg(\omega_X|_C) = \deg(C \cap X') + \deg(\omega_C) =
\deg(C \cap X') - 2\chi(C, \mathcal{O}_C)
$$
を得る．ゆえに補題が偽ならば，
$$
2[k' : k] \geq \deg(C \cap X') + 2\dim_k H^1(C, \mathcal{O}_C)
$$
となる．$C \cap X'$ は空でなく，上述の可除性もあるので，
これが起こりうるのは次のいずれかの場合に限る．
\begin{enumerate}
\item[(a)] $C \cap X'$ は一つの $k'$-有理点で，$C$ の点であり，
$H^1(C, \mathcal{O}_C) = 0$ である，または
\item[(b)] $C \cap X'$ は次数 $2$ を $k'$ 上にもち，
$H^1(C, \mathcal{O}_C) = 0$ である．
\end{enumerate}
第一の場合は $C$ が有理尾部であることを，第二の場合は
$C$ が有理橋であることを意味する．両方とも排除されているので証明が完了する．
\end{proof}

\begin{lemma}
\label{curves-lemma-contracting-rational-bridges}
$k$ を体とする．$X$ を $k$ 上の固有概形で，次元 $1$ かつ
$H^0(X, \mathcal{O}_X) = k$，種数 $g \geq 2$ を満たすものとする．
$X$ の特異点は高々節点的であり，$X$ は有理尾部をもたないと仮定する．
有理橋の収縮（例 \ref{curves-example-rational-bridge}）の列
$$
X = X_0 \to X_1 \to \ldots \to X_n = X'
$$
を，有理橋がなくなるまで考える．このとき $\omega_{X'}$ は豊富である．
射 $X \to X'$ は選択によらず，この射の形成は基礎体の拡大と可換である．
\end{lemma}

\begin{proof}
有理橋がなくなるまで順次収縮する．すると補題
\ref{curves-lemma-no-rational-bridge-ample-genus-geq-2} により
$\omega_{X'}$ は豊富である．

\medskip\noindent
$f : X \to X'$ を合成とする．補題
\ref{curves-lemma-rational-bridge-canonical} と帰納法により
$f^*\omega_{X'} = \omega_X$ である．
有理橋の収縮について成り立つので
$f_*\mathcal{O}_X = \mathcal{O}_{X'}$ である．したがって射影公式は，
$f_*f^*\mathcal{L} = \mathcal{L}$ が任意の可逆
$\mathcal{O}_{X'}$-加群 $\mathcal{L}$ に対して成り立つことを示す．ゆえに
$$
\Gamma(X', \omega_{X'}^{\otimes m}) = \Gamma(X, \omega_X^{\otimes m})
$$
がすべての $m$ に対して成り立つ．Morphisms，補題
\ref{morphisms-lemma-proper-ample-is-proj} により，$X'$ はこれらの直和の
Proj なので，射 $X \to X'$ は完全に標準的であると結論する．

\medskip\noindent
$K/k$ を体の拡大とする．このとき
$\omega_{X_K}$ は $\omega_X$ の引き戻しであり
（補題 \ref{curves-lemma-sanity-check-duality}），
$\Gamma(X, \omega_X^{\otimes m}) \otimes_k K$ は
$\Gamma(X_K, \omega_{X_K}^{\otimes m})$ に等しい．これは
Cohomology of Schemes，補題
\ref{coherent-lemma-flat-base-change-cohomology} による．
したがって上の議論により，$f : X \to X'$ の形成は
$K/k$ による基底変換と可換である．細部は省略する．
\end{proof}





\section{安定曲線への収縮}
\label{curves-section-contracting-to-stable}

\noindent
本節では，節 \ref{curves-section-contracting-rational-tails} および
\ref{curves-section-contracting-rational-bridges} で得た収縮射を組み合わせる．
すなわち，$k$ を体とし，$X$ を $k$ 上の固有概形で，次元 $1$，
$H^0(X, \mathcal{O}_X) = k$，種数 $g \geq 2$ を満たすものとする．
$X$ の特異点は高々節点的であると仮定する．補題
\ref{curves-lemma-contracting-rational-tails} の射と補題
\ref{curves-lemma-contracting-rational-bridges} の射を合成すると，射
$$
c : X \longrightarrow Y
$$
を得る．ここで $Y$ もまた $k$ 上の固有概形で，次元 $1$，
特異点が高々節点的であり，$k = H^0(Y, \mathcal{O}_Y)$，
種数 $g$ を満たす．さらに
$\mathcal{O}_Y = c_*\mathcal{O}_X$ かつ $R^1c_*\mathcal{O}_X = 0$
であり，$\omega_Y$ は $Y$ 上豊富である．補題
\ref{curves-lemma-characterize-contraction-to-stable} は，実際にこれらの条件が
この射を特徴づけることを示す．

\begin{lemma}
\label{curves-lemma-contract-gorenstein-canonical}
$k$ を体とする．$c : X \to Y$ を $k$ 上の固有概形の射とする．
次を仮定する．
\begin{enumerate}
\item $\mathcal{O}_Y = c_*\mathcal{O}_X$ かつ $R^1c_*\mathcal{O}_X = 0$，
\item $X$ と $Y$ は被約かつ Gorenstein で，次元 $1$ である，
\item $\exists\ m \in \mathbf{Z}$ で，
$H^1(X, \omega_X^{\otimes m}) = 0$ かつ $\omega_X^{\otimes m}$ が
大域切断で生成されるものがある．
\end{enumerate}
このとき $c^*\omega_Y \cong \omega_X$ である．
\end{lemma}

\begin{proof}
$c$ のファイバーは More on Morphisms，定理
\ref{more-morphisms-theorem-stein-factorization-Noetherian}
により幾何学的に連結である．特に $c$ は全射である．
閉点 $y = y_1, \ldots, y_r$（$Y$ の閉点）で $X_y$ が次元 $1$ をもつものは
有限個であり，$Y \setminus \{y_1, \ldots, y_r\}$ 上では射 $c$ は同型である．
細部は省略する．ヒント：$\{y_1, \ldots, y_r\}$ の外では射 $c$ は有限である；
Cohomology of Schemes，補題
\ref{coherent-lemma-characterize-finite} を参照せよ．

\medskip\noindent
写像 $b : c^*\omega_Y \to \omega_X$ を注意深く構成しよう．
$f : X \to \Spec(k)$ および $g : Y \to \Spec(k)$ を構造射とする．
補題 \ref{curves-lemma-duality-dim-1} とその証明により，
$f^!k = \omega_X[1]$ かつ $g^!k = \omega_Y[1]$ である．
すると $f^! = c^! \circ g^!$ であり，したがって
$c^!\omega_Y = \omega_X$ である．ゆえに，関手的同型
$$
\Hom_{D(\mathcal{O}_X)}(\mathcal{F}, \omega_X)
\longrightarrow
\Hom_{D(\mathcal{O}_Y)}(Rc_*\mathcal{F}, \omega_Y)
$$
が連接 $\mathcal{O}_X$-加群 $\mathcal{F}$ に対して存在する．これは
$c^!$ の定義による\footnote{Duality for Schemes，補題
\ref{duality-lemma-twisted-inverse-image} の右随伴を
$D^+_\QCoh(\mathcal{O}_Y)$ に制限したものとして．}．
この同型は跡写像 $t : Rc_*\omega_X \to \omega_Y$
（随伴の余単位）によって誘導される．射影公式
（Cohomology，補題 \ref{cohomology-lemma-projection-formula}）により，
標準写像 $a : \omega_Y \to Rc_*c^*\omega_Y$ は同型である．
以上を組み合わせると，標準写像
$b : c^*\omega_Y \to \omega_X$ で
$$
t \circ Rc_*(b) = a^{-1}
$$
を満たすものが存在する．特に $b$ を
$c^{-1}(Y \setminus \{y_1, \ldots, y_r\})$ に制限すると同型である
（可逆加群間の写像であり，別の写像との合成が同型 $a^{-1}$ になるからである）．

\medskip\noindent
(3) における $m \in \mathbf{Z}$ を選び，写像
$$
b^{\otimes m} :
\Gamma(Y, \omega_Y^{\otimes m})
\longrightarrow
\Gamma(X, \omega_X^{\otimes m})
$$
を考える．$Y$ は被約であり，構成時に述べた $b$ の最後の性質により，
この写像は単射である．Riemann--Roch（補題 \ref{curves-lemma-rr}）により
$\chi(X, \omega_X^{\otimes m}) =\chi(Y, \omega_Y^{\otimes m})$ である．
したがって
$$
\dim_k \Gamma(Y, \omega_Y^{\otimes m}) \geq
\dim_k \Gamma(X, \omega_X^{\otimes m}) = \chi(X, \omega_X^{\otimes m})
$$
であり，$b^{\otimes m}$ は大域切断上で同型を誘導すると結論する．ゆえに
$b^{\otimes m} : c^*\omega_Y^{\otimes m} \to \omega_X^{\otimes m}$
は全射である．実際，$\omega_X^{\otimes m}$ の生成元は像に入る．
したがって $b^{\otimes m}$ は同型であり，ゆえに $b$ は同型である．
\end{proof}

\begin{lemma}
\label{curves-lemma-characterize-contraction-to-stable}
$k$ を体とする．$X$ を $k$ 上の固有概形で，次元 $1$，
$H^0(X, \mathcal{O}_X) = k$，種数 $g \geq 2$ を満たすものとする．
$X$ の特異点は高々節点的であると仮定する．
射
$$
c : X \longrightarrow Y
$$
で，$k$ 上の概形の射であり，次の性質をもつものが，
一意な同型を除いて一意に存在する：
\begin{enumerate}
\item $Y$ は $k$ 上固有で，$\dim(Y) = 1$ であり，$Y$ の特異点は高々節点的である，
\item $\mathcal{O}_Y = c_*\mathcal{O}_X$ かつ $R^1c_*\mathcal{O}_X = 0$，かつ
\item $\omega_Y$ は $Y$ 上豊富である．
\end{enumerate}
\end{lemma}

\begin{proof}
存在：序論で述べたように，補題
\ref{curves-lemma-contracting-rational-tails} と
\ref{curves-lemma-contracting-rational-bridges} を組み合わせれば，
列挙したすべての性質をもつ射が存在する．さらに，これは合成
$$
X \to X_1 \to X_2 \ldots \to X_n \to X_{n + 1} \to \ldots \to X_{n + n'}
$$
として書ける．ここで最初の $n$ 個の射は有理尾部の収縮であり，
最後の $n'$ 個の射は有理橋の収縮である．性質 (2) は，有理尾部の各収縮
（例 \ref{curves-example-rational-tail}）および有理橋の各収縮
（例 \ref{curves-example-rational-bridge}）について成り立つことに注意する．
この性質が射の合成に継承されることは容易に分かる．

\medskip\noindent
一意性：$c : X \to Y$ を条件 (1)，(2)，(3) を満たす射とする．
一意な同型 $X_{n + n'} \to Y$ が存在し，射
$X \to X_{n + n'}$ および $c$ と両立することを示す．

\medskip\noindent
証明を始める前に，$c$ についていくつか観察する．まず，$c$ のファイバーは
More on Morphisms，定理
\ref{more-morphisms-theorem-stein-factorization-Noetherian}
により幾何学的に連結である．特に $c$ は全射である．
閉点 $y \in Y$ に対して，ファイバー $X_y$ は
$$
H^1(X_y, \mathcal{O}_{X_y}) = 0
\quad\text{かつ}\quad
H^0(X_y, \mathcal{O}_{X_y}) = \kappa(y)
$$
を満たす．第一の等式は More on Morphisms，補題
\ref{more-morphisms-lemma-check-h1-fibre-zero} により，第二の等式は
More on Morphisms，補題
\ref{more-morphisms-lemma-h1-fibre-zero-check-h0-kappa} による．
したがって，$X_y = x$ で，$x$ が $X$ の点のうち $y$ へ写る一意なもので
$y$ と同じ剰余体をもつか，または $X_y$ が次元 $1$ の固有概形で，
$\kappa(y)$ 上のものである．第二の場合，$X_y$ は Cohen--Macaulay である
（補題 \ref{curves-lemma-automatic}）．しかし $X$ は被約なので，$X_y$ は
そのすべての一般点で被約でなければならず（細部は省略する），
したがって Properties，補題 \ref{properties-lemma-criterion-reduced} により
$X_y$ は被約である．ゆえに $X_y$ の特異点は高々節点的である
（補題 \ref{curves-lemma-closed-subscheme-nodal-curve}）．
$X_y$ の種数は零であることに注意する（上を参照）．
最後に，点 $y$ で，ファイバー $X_y$ が次元 $1$ をもつものは有限個であり，
$\{y_1, \ldots, y_r\}$ と書ける．また
$c^{-1}(Y \setminus \{y_1, \ldots, y_r\})$ は
$Y \setminus \{y_1, \ldots, y_r\}$ へ $c$ により同型に写る．細部は省略する．
ヒント：$\{y_1, \ldots, y_r\}$ の外では射 $c$ は有限である；
Cohomology of Schemes，補題 \ref{coherent-lemma-characterize-finite} を参照せよ．

\medskip\noindent
$C \subset X$ を有理尾部とする．$c$ は $C$ を一点へ写すと主張する．
そうでないと仮定して矛盾を導く．すると $C$ の像は既約成分
$D \subset Y$ である．$H^0(C, \mathcal{O}_C) = k'$ は $k$ の
有限分離拡大であり，$C$ は $k'$-有理点 $x$ をもつことを思い出そう．
この点は $C$ と $X$ の「残りの部分」との一意な交点でもある．
上の一般的議論から，
$C \setminus \{x\} \subset c^{-1}(Y \setminus \{y_1, \ldots, y_r\})$
は開集合 $V$（$D$ の開集合）へ同型に写ると結論する．
$y = c(x) \in D$ とおく．$y$ は $D$ のうち $Y$ の「残りの部分」と交わる
唯一の点であることに注意する．
$y \not \in \{y_1, \ldots, y_r\}$ ならば $C \cong D$ であり，
$D$ が $Y$ の有理尾部であることは明らかである．これは
$\omega_Y$ の豊富性に矛盾する（補題 \ref{curves-lemma-rational-tail-negative}）．
したがって $y \in \{y_1, \ldots, y_r\}$ かつ $\dim(X_y) = 1$ である．
すると $x \in X_y \cap C$ であり，$x$ は $X_y$ と $C$ の滑らかな点である
（補題 \ref{curves-lemma-closed-subscheme-nodal-curve}）．
$y \in D$ が $D$ の特異点ならば，$y$ は節点であり，このとき
$Y = D$ である（補題 \ref{curves-lemma-closed-subscheme-nodal-curve} により，
$Y$ の別の成分が $y$ を通ることはできない）．
すると $X = X_y \cup C$ であり，これは $g = 0$ を意味する．実際，
例 \ref{curves-example-rational-tail} の議論により，その種数は
$X_y$ の種数に等しいからである．これは矛盾である．
$y \in D$ が $D$ の滑らかな点ならば，$C \to D$ は同型である
（非特異射影模型は一意であり，$C$ と $D$ は双有理だからである；
節 \ref{curves-section-curves-function-fields} を参照せよ）．すると
$D$ は $Y$ の有理尾部であり，$\omega_Y$ の豊富性に矛盾する．

\medskip\noindent
$n \geq 1$ と仮定する．$C \subset X$ が $X \to X_1$ により収縮される
有理尾部ならば，前段落により $C$ は $Y$ の一点へ写る．したがって
$c : X \to Y$ は $X \to X_1$ を経由して分解する
（$X$ は $C$ と $X_1$ の押し出しである；例
\ref{curves-example-rational-tail} の議論を参照せよ）．
$X$ を $X_1$ で置き換えると $n$ は減少する．帰納法により
$n = 0$，すなわち $X$ は有理尾部をもたないと仮定してよい．

\medskip\noindent
$n = 0$，すなわち $X$ は有理尾部をまったくもたないと仮定する．
補題 \ref{curves-lemma-no-rational-tail-semiample-genus-geq-2} により，
$\omega_X^{\otimes 2}$ と $\omega_X^{\otimes 3}$ は大域生成である．
補題 \ref{curves-lemma-vanishing-twist} により
$H^1(X, \omega_X^{\otimes 3}) = 0$ が従う．
$m = 3$ として補題 \ref{curves-lemma-contract-gorenstein-canonical} を適用すると，
$c^*\omega_Y \cong \omega_X$ を得る．また補題
\ref{curves-lemma-rational-bridge-canonical} と帰納法により
$\omega_X = (X \to X_{n'})^*\omega_{X_{n'}}$ である．
$c$ と $X \to X_{n'}$ の双方に射影公式を適用すると，
$$
\Gamma(X_{n'}, \omega_{X_{n'}}^{\otimes m}) =
\Gamma(X, \omega_X^{\otimes m}) =
\Gamma(Y, \omega_Y^{\otimes m})
$$
がすべての $m$ に対して成り立つと結論する．Morphisms，補題
\ref{morphisms-lemma-proper-ample-is-proj} により，$X_{n'}$ と $Y$ は
これらの直和の Proj なので，所望の標準同型 $X_{n'} = Y$ が存在すると結論する．
この図式を可換にする同型が一意であることの確認は省略する．
\end{proof}

\begin{lemma}
\label{curves-lemma-tricanonical}
$k$ を体とする．$X$ を $k$ 上の固有概形で，次元 $1$，
$H^0(X, \mathcal{O}_X) = k$，種数 $g \geq 2$ を満たすものとする．
$X$ の特異点は高々節点的で，$\omega_X$ は豊富であると仮定する．このとき
$\omega_X^{\otimes 3}$ は非常に豊富であり，
$H^1(X, \omega_X^{\otimes 3}) = 0$ である．
\end{lemma}

\begin{proof}
Varieties，補題 \ref{varieties-lemma-ampleness-in-terms-of-degrees-components} と
補題 \ref{curves-lemma-rational-tail-negative} および
\ref{curves-lemma-rational-bridge-zero} を組み合わせると，$X$ は有理尾部も有理橋も
含まないことが分かる．すると補題 \ref{curves-lemma-contracting-rational-tails} により
$\omega_X^{\otimes 3}$ は大域生成である．
$k$-基底 $s_0, \ldots, s_n$（$H^0(X, \omega_X^{\otimes 3})$ の基底）を選ぶ．
射
$$
\varphi_{\omega_X^{\otimes 3}, (s_0, \ldots, s_n)} :
X \longrightarrow \mathbf{P}^n_k
$$
を得る．Constructions，節 \ref{constructions-section-projective-space} を参照せよ．
補題は，この射が閉埋め込みであると主張している．これを確認するため，
$k$ をその代数閉包で置き換えてよい；Descent，補題
\ref{descent-lemma-descending-property-closed-immersion} を参照せよ．
したがって $k$ は代数閉であると仮定してよい．

\medskip\noindent
$k$ は代数閉であると仮定する．Varieties，補題
\ref{varieties-lemma-variant-separate-points-tangent-vectors}
を用いて補題を証明する．$Z \subset X$ を，次数 $2$ を $Z$ 上にもつ閉部分概形で，
イデアル層 $\mathcal{I} \subset \mathcal{O}_X$ をもつものとする．
写像
$$
H^0(X, \mathcal{L}) \to H^0(Z, \mathcal{L}|_Z)
$$
が全射であることを示さなければならない．したがって
$H^1(X, \mathcal{I}\mathcal{L}) = 0$ を示せば十分である．
このため補題 \ref{curves-lemma-vanishing-on-gorenstein} を用いる．
したがって，
$$
3\deg(\omega_X|_Y) > -2\chi(Y, \mathcal{O}_Y) + \deg(Z \cap Y)
$$
をすべての被約連結閉部分概形 $Y \subset X$ に対して示せば十分である．
$k$ は代数閉であり，$Y$ は連結かつ被約なので，
$H^0(Y, \mathcal{O}_Y) = k$ である（Varieties，補題
\ref{varieties-lemma-proper-geometrically-reduced-global-sections}）．
したがって $\chi(Y, \mathcal{O}_Y) = 1 - \dim H^1(Y, \mathcal{O}_Y)$ である．
ゆえに
$$
3\deg(\omega_X|_Y) > -2 + 2\dim H^1(Y, \mathcal{O}_Y) + \deg(Z \cap Y)
$$
を示さなければならない．補題 \ref{curves-lemma-no-rational-tail} により，
これは $Y = X$ または $\deg(\omega_X|_Y) = 0$ の場合を除いて真である．
$\omega_X$ は豊富なので第二の場合は起こらない
（この証明で最初に引用した補題を参照せよ）．最後に $Y = X$ ならば，
Riemann--Roch（補題 \ref{curves-lemma-rr}）と $g \geq 2$ を用いて
不等式が成り立つことが分かる．$Z = \emptyset$ とした同じ議論により
$H^1(X, \omega_X^{\otimes 3}) = 0$ である．
\end{proof}






\section{ベクトル場}
\label{curves-section-vector-fields}

\noindent
本節では，曲線上のベクトル場の空間を調べる．
ベクトル場は無限小自己同型に対応する（More on Morphisms，第
\ref{more-morphisms-section-action-by-derivations} 節を参照）ので，
モジュライ理論において重要な役割を果たす．

\medskip\noindent
$k$ を代数閉体とする．
$X$ を $k$ 上有限型の概形とする．
$x \in X$ を閉点とする．
$D \in \text{Der}_k(\mathcal{O}_X, \mathcal{O}_X)$ が
{\it $x$ を固定する}とは，$D(\mathcal{I}) \subset \mathcal{I}$ が
成り立つことをいう．ここで $\mathcal{I} \subset \mathcal{O}_X$ は
$x$ のイデアル層である．

\begin{lemma}
\label{curves-lemma-smooth-vector-fields}
$k$ を代数閉体とする．
$X$ を $k$ 上の滑らかで固有な連結曲線とする．
$g$ を $X$ の種数とする．
\begin{enumerate}
\item $g \geq 2$ ならば，
$\text{Der}_k(\mathcal{O}_X, \mathcal{O}_X)$ は零である．
\item $g = 1$ かつ
$D \in \text{Der}_k(\mathcal{O}_X, \mathcal{O}_X)$ が零でなければ，
$D$ は $X$ のいかなる閉点も固定しない．
\item $g = 0$ かつ
$D \in \text{Der}_k(\mathcal{O}_X, \mathcal{O}_X)$ が零でなければ，
$D$ は高々 $2$ 個の $X$ の閉点を固定する．
\end{enumerate}
\end{lemma}

\begin{proof}
普遍 $k$-導分 $d : \mathcal{O}_X \to \Omega_{X/k}$ が存在するので，
$D = \theta \circ d$ となることを思い出そう．ここで
ある $\mathcal{O}_X$-線形写像
$\theta : \Omega_{X/k} \to \mathcal{O}_X$ を用いている．補題
\ref{curves-lemma-duality-dim-1} により
$\Omega_{X/k} \cong \omega_X$ である．Riemann--Roch により
$\deg(\omega_X) = 2g - 2$ である（補題 \ref{curves-lemma-rr}）．
したがって $\theta$ は，$g > 1$ ならば，Varieties，補題
\ref{varieties-lemma-check-invertible-sheaf-trivial} により
零でなければならない．これで (1) が示された．
$g = 1$ ならば零でない $\theta$ はどこでも消えず，
$g = 0$ ならば零でない $\theta$ は次数 $2$ の因子に沿って消える．
ゆえに，$\theta$ が閉点 $x \in X$ において消えることと，
上で定義した意味で $D$ が $x$ を固定することとが同値であると示せば，
(2) と (3) が従う．$z \in \mathcal{O}_{X, x}$ を一意化元とする．
すると $dz$ は $\Omega_{X, x}$ の基底元である（補題
\ref{curves-lemma-uniformizer-works} を参照）．
$D(z) = \theta(dz)$ であるから結論を得る．
\end{proof}

\begin{lemma}
\label{curves-lemma-nodal-vector-fields}
$k$ を代数閉体とする．
$X$ を $1$ 次元の $k$ 上固有かつ連結な概形で，その特異点が
高々節点的であるものとする．$\nu : X^\nu \to X$ を正規化とする．
$S \subset X^\nu$ を，$\nu$ が同型でない点の集合とする．このとき
$$
\text{Der}_k(\mathcal{O}_X, \mathcal{O}_X) =
\{D' \in \text{Der}_k(\mathcal{O}_{X^\nu}, \mathcal{O}_{X^\nu}) \mid
D' \text{ はすべての }x^\nu \in S\text{ を固定する}\}
$$
\end{lemma}

\begin{proof}
$x \in X$ を節点とする．$x', x'' \in X^\nu$ を $x$ の逆像とする．
（$k$ は代数閉体なので，すべての節点は分裂節点である．定義
\ref{curves-definition-split-node} および補題 \ref{curves-lemma-split-node} を参照．）
% P07-ERR-0378
$u \in \mathcal{O}_{X^\nu, x'}$ および
$v \in \mathcal{O}_{X^\nu, x''}$ を一意化元とする．次の完全列がある：
$$
0 \to \mathcal{O}_{X, x} \to
\mathcal{O}_{X^\nu, x'} \times \mathcal{O}_{X^\nu, x''} \to k \to 0
$$
これは補題 \ref{curves-lemma-multicross} から従う．したがって $u$ と $v$ を
$\mathcal{O}_{X, x}$ の元とみなせ，それらは $uv = 0$ を満たす．

\medskip\noindent
$D \in \text{Der}_k(\mathcal{O}_X, \mathcal{O}_X)$ とする．このとき
$0 = D(uv) = vD(u) + uD(v)$ である．$(u)$ は $v$ の
$\mathcal{O}_{X, x}$ における零化イデアルであり，逆も同様なので，
$D(u) \in (u)$ かつ $D(v) \in (v)$ であることが分かる．
$\mathcal{O}_{X^\nu, x'} = k + (u)$ であるから，$D$ を
$\mathcal{O}_{X^\nu, x'}$ へ延長でき，しかもその延長は $x'$ を固定する．
これにより，等式の右辺に属する $D'$ が得られる．逆に，
$D'$ が与えられ，それが $x'$ と $x''$ を固定するとき，$D'$ は部分環
$\mathcal{O}_{X, x} \subset
\mathcal{O}_{X^\nu, x'} \times \mathcal{O}_{X^\nu, x''}$
を保つことが分かり，これによって等式の右辺から左辺へ移る．
\end{proof}

\begin{lemma}
\label{curves-lemma-stable-vector-fields}
$k$ を代数閉体とする．
$X$ を $1$ 次元の $k$ 上固有かつ連結な概形で，その特異点が
高々節点的であるものとする．$X$ の種数は少なくとも $2$ であり，
$X$ は有理尾部も有理橋ももたないと仮定する．このとき
$\text{Der}_k(\mathcal{O}_X, \mathcal{O}_X) = 0$.
\end{lemma}

\begin{proof}
$D \in \text{Der}_k(\mathcal{O}_X, \mathcal{O}_X)$ とする．
$X^\nu$ を $X$ の正規化とする．
$D' \in \text{Der}_k(\mathcal{O}_{X^\nu}, \mathcal{O}_{X^\nu})$ を，
補題 \ref{curves-lemma-nodal-vector-fields} により $D$ に対応する元とする．
$C \subset X^\nu$ を既約成分とする．$C$ の種数が $> 1$ ならば，
補題 \ref{curves-lemma-smooth-vector-fields} の (1) により
$D'|_{\mathcal{O}_C} = 0$ である．$C$ の種数が $1$ ならば，
閉点 $c$ が $C$ 上に少なくとも一つあり，それは $X$ の節点へ写る
（そうでなければ $X \cong C$ の種数が $1$ となる）．
上の対応により，$D'|_{\mathcal{O}_C}$ は $c$ を固定するので，
補題 \ref{curves-lemma-smooth-vector-fields} の (2) により零である．
最後に，$C$ の種数が零であるとする．このとき，少なくとも $3$ 個の
互いに異なる閉点 $c_1, c_2, c_3 \in C$ が存在し，それらは $X$ の節点へ写る．
実際，そうでなければ，$X$ は $C$ の二点を貼り合わせたもの
（$C$ の二点が同じ節点へ写る）であるか，$C$ は有理橋
（二点が $X$ の異なる節点へ写る）であるか，または $C$ は有理尾部
（一点が $X$ の節点へ写る）である．
% P07-ERR-0379
$C$ の種数は $\geq 2$ であり，有理橋も有理尾部ももたないので，
これら三つの可能性はいずれも許されない．したがって
$D'|_{\mathcal{O}_C}$ は $c_1, c_2, c_3$ を固定するので，
補題 \ref{curves-lemma-smooth-vector-fields} の (3) により零である．
\end{proof}
